
\documentclass[UTF8,openany,zihao=-4]{ctexbook}
\usepackage[a4paper,top=1.45cm,bottom=1.45cm,left=1.35cm,right=1.35cm,headsep=0.25cm,footskip=0.65cm]{geometry}
\usepackage{amsmath,amssymb}
\usepackage{enumitem}
\usepackage{booktabs,longtable,array,tabularx,multicol}
\usepackage{xcolor}
\usepackage{hyperref}
\usepackage{fancyhdr}
\usepackage{lastpage}
\usepackage{titlesec}
\hypersetup{colorlinks=true,linkcolor=black,urlcolor=black,citecolor=black}
\setlength{\parindent}{0em}
\setlength{\parskip}{0.18em}
\linespread{1.02}
\raggedbottom
\pagestyle{fancy}
\setlength{\headheight}{15pt}
\fancyhf{}
\fancyhead[L]{操作系统原理高密度理解讲义}
\fancyhead[R]{\leftmark}
\fancyfoot[C]{\thepage/\pageref{LastPage}}
\renewcommand{\headrulewidth}{0.25pt}
\ctexset{
  chapter = {format = \huge\bfseries},
  section = {format = \Large\bfseries},
  subsection = {format = \large\bfseries}
}
\titlespacing*{\chapter}{0pt}{-0.6em}{0.8em}
\titlespacing*{\section}{0pt}{0.55em}{0.22em}
\titlespacing*{\subsection}{0pt}{0.35em}{0.12em}
\newlist{points}{enumerate}{2}
\setlist[points,1]{label=\arabic*. ,leftmargin=2.0em,itemsep=0.06em,topsep=0.06em,parsep=0em,partopsep=0em}
\setlist[points,2]{label=(\arabic*) ,leftmargin=1.8em,itemsep=0.02em,topsep=0.02em,parsep=0em,partopsep=0em}
\setlist[itemize]{itemsep=0.04em,topsep=0.04em,parsep=0em,partopsep=0em}
\setlist[enumerate]{itemsep=0.04em,topsep=0.04em,parsep=0em,partopsep=0em}
\newcommand{\concept}[1]{\subsection{#1}}
\newcommand{\key}[1]{\textbf{#1}}
\newcommand{\term}[1]{\textbf{#1}}
\newcommand{\code}[1]{\texttt{#1}}
\newcommand{\zhushi}[2]{\par\noindent{\kaishu\small #1：#2}\par}
\newcommand{\plain}[1]{\zhushi{大白话}{#1}}
\newcommand{\exam}[1]{\zhushi{可能怎么考}{#1}}
\newcolumntype{Y}{>{\raggedright\arraybackslash}X}
\title{\bfseries 操作系统原理高密度理解讲义\\\large 按 PPT 顺序：知识点 + 大白话 + 可能怎么考}
\author{基于课程 PPT、往年试卷与个人 os.docx 辅助记忆整理}
\date{\today}

\begin{document}
\maketitle
\tableofcontents
\chapter*{使用说明}
\addcontentsline{toc}{chapter}{使用说明}
\begin{points}
  \item 本讲义按 PPT 顺序组织，尽量覆盖从绪论到设备管理的完整考试知识链。
  \item 每个知识点先用普通字体写“需要背和会用的知识”，再用楷体写“\textbf{大白话}”和“\textbf{可能怎么考}”。
  \item 背诵时先背普通字体；做题前看楷体部分，建立题目识别直觉。
  \item 版式改为高密度文本，不使用花哨色块；重点靠编号、粗体和楷体区分。
\end{points}

\chapter{操作系统绪论}

\section{操作系统的概念与定位}
\concept{计算机系统的层次}
\begin{points}
  \item 计算机系统由\term{硬件}和\term{软件}组成。硬件包括 CPU、内存、I/O 设备等；软件包括系统软件和应用软件。
  \item 操作系统是配置在裸机上的\key{第一层软件}，位于硬件与其他软件之间，是其他软件运行的基础。
  \item 编译器、数据库管理系统、编辑程序也属于系统软件，但它们不是操作系统本身。
  \item 从层次上看：硬件/裸机 $\rightarrow$ 操作系统 $\rightarrow$ 系统程序 $\rightarrow$ 应用程序 $\rightarrow$ 用户。
\end{points}
\plain{硬件不会直接服务应用，OS 是硬件上第一层“扩展机”，把裸机包装成可管理、可调用的机器。}
\exam{常考“谁是第一层软件”“编译器/数据库是不是 OS”“OS 位于哪两层之间”。}

\concept{为什么需要操作系统}
\begin{points}
  \item \term{从用户视角看}：操作系统是用户和计算机硬件之间的接口，为用户提供方便、安全、高效的使用环境。
  \item \term{从系统视角看}：操作系统是计算机系统资源的管理者，负责分配、回收、保护和调度资源。
  \item 操作系统通过三种思想扩展计算机能力：\key{资源复用、资源虚化、资源抽象}。
  \item 复用解决“资源数量不够”：CPU 轮流给多个进程用是时间复用；内存分块给多个进程用是空间复用。
  \item 虚化解决“一个物理资源变成多个逻辑资源”：一个 CPU 虚拟成多个虚拟处理机，一台打印机通过 SPOOLing 虚拟成多台逻辑打印机。
  \item 抽象解决“硬件太复杂”：磁盘扇区被抽象成文件，CPU 执行流被抽象成进程，物理设备被抽象成逻辑设备名。
\end{points}
\plain{没有 OS，用户要直接面对 CPU、内存、磁盘、设备；OS 的价值就是管资源、挡危险、把复杂硬件变简单。}
\exam{常考从用户视角/系统视角回答，或解释资源复用、虚化、抽象。}


\concept{操作系统定义}
\begin{points}
  \item 操作系统没有唯一的精确定义，但可从两个角度理解。
  \item \term{控制程序}：控制程序执行，防止错误和不当使用，协调应用与硬件。
  \item \term{资源管理器}：管理处理机、存储器、I/O 设备、文件、磁盘空间等软硬件资源。
  \item 简洁表述：操作系统是位于裸机和应用程序之间的系统软件，对下管理控制计算机软硬件资源，对上提供方便、安全、高效的接口和运行环境。
\end{points}
\plain{定义题不要追求一句神句，抓住“系统软件、资源管理者、用户接口/扩展机”三个关键词就稳。}
\exam{常考简答“什么是操作系统”，按“位置-对象-作用-接口”写。}


\section{操作系统的形成与发展}
\concept{发展阶段总览}
\begin{points}
  \item 操作系统发展顺序：\key{手工操作 $\rightarrow$ 批处理系统 $\rightarrow$ 多道批处理系统 $\rightarrow$ 分时系统 $\rightarrow$ 实时系统 $\rightarrow$ 通用操作系统}。
  \item 每个阶段都是在解决上一阶段的主要矛盾：人机矛盾、CPU 与 I/O 速度矛盾、交互需求、实时响应需求。
\end{points}
\plain{发展史就是不断解决矛盾：人机慢、CPU 等 I/O、用户要交互、任务有截止期限。}
\exam{常考按阶段排序，或问某阶段解决了什么问题。}


\concept{手工操作阶段}
\begin{points}
  \item 用户独占机器，人工装入程序、启动运行、取走结果。
  \item 主要矛盾：人太慢，CPU 太快，I/O 太慢，导致资源利用率很低。
  \item 没有真正的操作系统，只是人手工控制机器。
\end{points}
\plain{人手工装程序，CPU 大量等人和 I/O，所以效率极低。}
\exam{常考缺点：用户独占资源、CPU 利用率低、没有自动作业控制。}


\concept{批处理系统}
\begin{points}
  \item 把多个作业组织成一批，由监督程序 Monitor 自动控制，一个接一个执行。
  \item 优点：减少人工干预，提高吞吐量。
  \item 缺点：缺少交互性，用户提交作业后不能直接干预运行过程。
  \item 关键词：\key{自动、成批、无交互、吞吐量、周转时间}。
\end{points}
\plain{把一堆作业交给监督程序自动跑，减少人工干预，但用户不能边运行边改。}
\exam{常考关键词“成批、自动、无交互、吞吐量”。}


\concept{联机批处理与脱机批处理}
\begin{points}
  \item 联机批处理的输入输出在监督程序控制下与主机同步进行，实现简单，但 I/O 速度慢，CPU 更容易等待。
  \item 脱机批处理把输入输出任务交给外围机完成，主机主要负责计算和从输入带/向输出带取送作业。
  \item 脱机方式的优点是减少主机空闲时间、提高 I/O 效率，并形成一定的异步能力。
  \item 直觉：主机不必边算边亲自管慢速输入输出，而是先把作业和结果放到中间介质，再分工处理。
\end{points}
\plain{联机是主机边算边直接管输入输出，脱机是把慢速 I/O 先交给外围机，主机专心算。}
\exam{常考联机批处理和脱机批处理的区别，以及脱机批处理为什么能提高效率。}


\concept{单道批处理系统特征}
\begin{points}
  \item 单道批处理系统内存中始终只保留一道作业运行。
  \item 其特征包括：\key{自动性、顺序性、单道性}。
  \item 自动性指作业自动依次运行；顺序性指作业完成顺序与进入顺序一致；单道性指任一时刻内存只有一道作业。
  \item 它仍然无法解决 CPU 等 I/O 的主要浪费问题。
\end{points}
\plain{单道批处理虽然已经自动切换作业，但内存里仍只有一道作业，所以一个作业等 I/O 时 CPU 还是可能闲着。}
\exam{常考单道批处理的三个特征，以及它为什么资源利用率仍不高。}


\concept{多道批处理系统}
\begin{points}
  \item 多个作业同时进入内存，交替运行；当一个作业等待 I/O 时，CPU 可执行另一个作业。
  \item 核心目标：解决 CPU 等待 I/O 的浪费，提高 CPU 和 I/O 设备利用率。
  \item 代价：需要进程调度、内存分配、资源保护、中断处理、上下文切换等机制，系统开销更大。
  \item 直觉例子：作业 A 等磁盘读入时，作业 B 用 CPU 计算；磁盘和 CPU 可以并行工作。
\end{points}
\plain{多个作业进内存，一个等 I/O，另一个用 CPU，让 CPU 和设备都别闲着。}
\exam{常考为什么能提高效率、为什么需要中断/调度/存储保护。}


\concept{多道批处理系统特征}
\begin{points}
  \item \term{多道性}：内存中可同时驻留多道相互独立的程序。
  \item \term{宏观上并行}：从较长时间范围看，多道程序都在向前推进。
  \item \term{微观上串行}：在单 CPU 条件下，任一时刻真正占用处理机执行的通常只有一道程序。
  \item \term{无交互性}：用户提交作业后通常不能像分时系统那样随时交互控制。
\end{points}
\plain{多道批处理是“宏观看大家都在跑，微观看 CPU 还是轮流给各作业用”，别把它和真正并行或分时交互混在一起。}
\exam{常考多道批处理的典型特征，尤其“宏观并行、微观串行、无交互”。}


\concept{多道程序设计的优点与缺点}
\begin{points}
  \item 优点：提高 CPU 利用率、内存利用率和 I/O 设备利用率，从而提高系统吞吐量。
  \item 其本质原因是：当某道程序等待 I/O 时，CPU 可转去执行另一道程序，系统中多种资源能更充分并行工作。
  \item 主要缺点：单个作业的周转时间可能延长，因为作业之间要竞争 CPU 和其他资源。
  \item 因而多道程序设计更偏向整体效率优化，不一定让每个单独作业都更快完成。
\end{points}
\plain{多道程序设计优化的是“整机忙不忙、单位时间做完多少活”，不是保证“每个作业都更早结束”；这正是吞吐量和周转时间之间的典型权衡。}
\exam{PPT 直接把 CPU 利用率、平均周转时间、系统吞吐量放在同一题里比较，答题时要明确写出“整体效率提高，但周转时间未必缩短”。}


\concept{多道程序设计的道数问题}
\begin{points}
  \item 多道程序度不是越大越好，它受内存容量、设备数量、作业特性以及响应时间要求等因素限制。
  \item 若程序平均等待 I/O 的时间占运行时间比例为 $p$，内存中有 $n$ 道同类程序，则 CPU 空闲概率近似为 $p^n$。
  \item 因而 CPU 利用率近似为 $1-p^n$，说明随着道数增加，CPU 利用率先提升，但提升幅度会逐渐减小。
  \item 若系统关键瓶颈不是 CPU，而是打印机、磁盘或内存容量，则盲目增加道数未必继续提高效率，反而可能加重竞争与响应延迟。
\end{points}
\plain{这节不是让你死背公式，而是理解“加作业确实能减少 CPU 发呆，但收益会递减，而且别的资源也可能先撑不住”。}
\exam{计算题常直接给 $p$ 和 $n$ 让你算 CPU 利用率，也可能追问“为什么内存增加后利用率继续上升但边际收益变小”。}


\concept{分时系统}
\begin{points}
  \item 多个用户通过终端同时使用计算机，系统把 CPU 时间划分为时间片，轮流响应各用户请求。
  \item 追求目标：较快响应用户，使每个用户感觉自己独占机器。
  \item 特征：\key{多路性、交互性、独立性、及时性}。
  \item 注意：分时系统的及时性是“用户命令较快响应”，不是实时系统那种严格截止期限。
  \item 若时间片为 $Q$，用户数为 $N$，粗略响应周期可理解为 $T=Q\times N$。用户越多，每个用户轮到 CPU 的间隔越长。
\end{points}
\plain{把 CPU 切成很多时间片轮流给用户，目标是让每个终端都感觉机器在响应自己。}
\exam{常考分时四特征，尤其“及时性”不等于硬实时。}


\concept{实时系统}
\begin{points}
  \item 实时系统强调在规定时间内响应外部事件。不是单纯“速度快”，而是响应必须\key{及时、可预测、可靠}。
  \item 硬实时：错过截止时间会造成严重失败，例如飞控、安全控制。
  \item 软实时：偶尔超时可以容忍，例如多媒体播放中偶尔卡顿。
  \item 关键词：\key{及时性、可靠性、截止时间、硬实时、软实时}。
\end{points}
\plain{实时不是“跑得快”而是“截止时间前必须响应”。}
\exam{常考硬实时/软实时区别、实时系统追求可靠性和可预测性。}


\section{操作系统的基本特征}
\concept{并发性}
\begin{points}
  \item 并发指若干事件在\key{同一时间间隔内}发生。
  \item 并行指若干事件在\key{同一时刻}发生。
  \item 单 CPU 上多个进程可以并发，但不能并行；多 CPU/多核上多个进程可以真正并行。
  \item 课件常用判断：并行一定是并发，但并发不一定是并行；并行是并发的特例。
\end{points}
\plain{核心是“看起来同时推进”。单 CPU 上靠时间片交替，多 CPU 上才可能真的同一时刻运行。}
\exam{常考并发与并行区别、单处理机下最多几个进程运行、并发导致为什么需要同步互斥。}


\concept{共享性}
\begin{points}
  \item 共享指系统资源可供多个并发执行的进程共同使用。
  \item \term{互斥共享}：一段时间内只允许一个进程访问，例如打印机、临界区。
  \item \term{同时共享}：宏观上多个进程可同时访问，例如只读文件、共享代码段。
  \item 并发与共享互为条件：没有并发，就谈不上资源竞争；没有共享管理，并发程序也难以正确运行。
\end{points}
\plain{多个进程都想用同一批资源，所以必须规定能不能一起用、谁先用、用完怎么释放。}
\exam{常考互斥共享/同时共享举例，或者问并发和共享为什么是 OS 最基本特征。}


\concept{虚拟性}
\begin{points}
  \item 虚拟性指通过某种技术把一个物理实体变成若干逻辑上的对应物。
  \item 典型例子：虚拟处理机、虚拟存储器、虚拟设备。
  \item 虚拟不是“凭空变多”，而是通过复用、映射、换入换出等机制，让用户看到更方便的逻辑资源。
\end{points}
\plain{不是资源真的变多，而是 OS 用映射、换入换出、排队等办法，让用户感觉自己有一份独立资源。}
\exam{常考虚拟处理机、虚拟存储器、虚拟设备的例子，以及虚拟性和资源复用的关系。}


\concept{异步性/不确定性}
\begin{points}
  \item 异步性指并发环境下，程序以不可预知的速度向前推进。
  \item 一个进程何时被调度、何时暂停、I/O 等多久、何时继续执行，不完全由程序自己决定。
  \item 异步性带来问题：若没有同步互斥机制，可能出现数据不一致、结果不可再现。
\end{points}
\plain{进程什么时候跑、什么时候停，不完全由程序自己决定，而是被调度、中断、I/O 共同影响。}
\exam{常考并发程序为什么不可再现、为什么会有与时间有关的错误。}


\section{操作系统的功能与接口}
\concept{功能模块}
\begin{points}
  \item \term{处理机管理/进程管理}：进程控制、进程同步、进程通信、调度。
  \item \term{存储器管理}：内存分配与回收、地址映射、保护共享、存储扩充。
  \item \term{设备管理}：I/O 控制、设备分配、缓冲、设备驱动、设备独立性。
  \item \term{文件管理}：文件存储、目录、文件共享、文件保护、外存空间管理。
  \item \term{作业管理/用户接口}：作业控制、命令接口、系统调用接口、图形接口。
\end{points}
\plain{五大管理就是 OS 的主线：CPU、内存、设备、文件、用户/作业。后面章节基本都从这里展开。}
\exam{常考某功能属于哪个模块，如进程调度属于处理机管理，文件目录属于文件管理。}


\concept{OS 提供的公共和基本服务}
\begin{points}
  \item 程序执行：把程序装入内存并启动运行。
  \item I/O 操作：统一管理设备，为程序提供输入输出服务。
  \item 信息保存：通过文件系统提供文件创建、读写、删除等服务。
  \item 通信服务：为进程之间交换信息提供支持。
  \item 错误检测与报告：发现程序或系统运行中的错误并报告给用户或操作员。
  \item 资源分配：为进程运行分配 CPU、内存、设备等资源。
  \item 统计/记账：记录用户使用资源的类型和数量。
  \item 保护与安全：控制对信息和资源的访问，防止非法使用。
\end{points}
\plain{这一节问的不是“OS 管哪几大模块”，而是“站在用户角度，OS 到底替你做了哪些通用服务”；考试里常把两组概念混着考。}
\exam{PPT 把这 8 项单独列成服务清单，选择题常考“资源分配/错误检测/记账/保护”是否属于 OS 基本服务。}


\concept{操作系统提供的接口}
\begin{points}
  \item \term{命令接口}：用户直接输入命令控制系统，如 \code{ls}、\code{cd}、\code{mkdir}。Shell 本质是命令解释器。
  \item 命令接口分为联机用户接口和脱机用户接口。联机接口用于交互式控制；脱机接口常见于批处理作业控制语言。
  \item \term{程序接口}：程序通过系统调用请求 OS 服务，如 \code{open()}、\code{read()}、\code{fork()}。
  \item \term{图形接口}：窗口、菜单、鼠标等 GUI 形式，本质上也是调用系统服务。
  \item 广义指令常作为系统调用的另一种说法，属于程序接口。
\end{points}
\plain{这节最容易混的是“系统调用”到底归哪类接口：它不是给人直接敲的命令，而是程序向 OS 提服务请求的程序接口。}
\exam{常考命令接口、程序接口、GUI 三类接口划分，尤其系统调用应归入程序接口而不是命令接口。}


\concept{脱机命令接口与联机命令接口}
\begin{points}
  \item \term{脱机命令接口}常由作业控制语言 JCL 组成，用户事先写好作业说明书，连同作业一起提交，由系统后续按说明书自动解释执行。
  \item 脱机控制方式下，用户不能在作业运行过程中直接交互干预。
  \item \term{联机命令接口}提供交互式命令，用户在终端输入命令后，由命令解释程序立即获取并执行下一条命令。
  \item 联机命令又常分为\term{内部命令}和\term{外部命令}：前者短小高频、常驻内存；后者功能较复杂，通常作为磁盘文件按需装入。
\end{points}
\plain{脱机接口像“提前写好任务单交给系统批量办”，联机接口像“边对话边下命令”；这两类题最容易和批处理/分时环境一起混考。}
\exam{PPT 专门把 JCL、内部命令、外部命令分开讲过，选择题常考谁常驻内存、谁按需从磁盘调入。}


\concept{图形用户接口 GUI}
\begin{points}
  \item GUI 以窗口、菜单、图标和鼠标操作作为主要交互方式。
  \item 它通过把命令与屏幕对象绑定，让用户通过点击、拖动等方式完成控制。
  \item 从本质上看，图形接口仍可视为命令接口的图形化形态。
\end{points}
\plain{GUI 不是第四类全新接口，本质上还是在用更直观的方式把原本要敲的命令和服务请求包装起来。}
\exam{判断题常把 GUI 和命令接口对立起来说；更稳的理解是 GUI 属于图形化的人机接口形式。}


\concept{系统调用分类}
\begin{points}
  \item 设备管理类：完成设备请求、释放、启动等功能。
  \item 文件管理类：完成文件创建、读写、删除等功能。
  \item 进程控制类：完成进程创建、撤销、阻塞、唤醒等功能。
  \item 进程通信类：完成消息传递、信号传递等功能。
  \item 内存管理类：完成内存分配、回收以及相关地址空间管理功能。
\end{points}
\plain{系统调用分类题本质是在问“程序想向 OS 申请哪类服务”，不要把它答成“命令接口/程序接口/图形接口”那一层。}
\exam{这类选择题经常给出一个调用场景，让你判断它属于设备管理、文件管理还是进程控制类系统调用。}

\concept{系统调用与特权操作}
\begin{points}
  \item I/O 指令、修改中断开关、修改页表、设置定时器等通常是特权指令，只能在内核态执行。
  \item 普通用户程序不能直接访问设备或随意操作内存，否则可能破坏其他进程或系统本身。
  \item 用户程序通过系统调用陷入内核，切换到内核态，由操作系统执行受保护的操作，再返回用户态。
\end{points}
\plain{系统调用这一节的关键不是接口分类本身，而是“为什么必须陷入内核去做”：因为 I/O、页表、中断等操作一旦放给用户程序会直接破坏系统保护。}
\exam{常考为什么 I/O 指令、改页表、关中断等要通过系统调用或内核态完成，以及系统调用会不会触发模式切换。}

\concept{系统调用的实现要点}
\begin{points}
  \item 内核中要为每类系统调用编写对应的处理程序。
  \item 系统通常维护一张\term{系统调用入口地址表}，按系统调用号找到对应处理程序入口；有的系统还会记录参数个数等信息。
  \item 陷入处理机制要能保护现场，保存发生系统调用时的处理器状态，便于返回后继续执行用户程序。
\end{points}
\plain{实现要点可以概括成“三件套”：有处理程序、有入口表、能保护现场；少了任何一个，系统调用都无法从“编号”真正落到内核代码。}
\exam{简答题常把“入口地址表”和“现场保护”作为关键词，答题时不要只写“切到内核态执行”这种空话。}

\concept{系统调用的调用过程}
\begin{points}
  \item \term{准备阶段}：保存现场，并把系统调用号和参数放入约定的寄存器或存储单元。
  \item \term{执行阶段}：通过陷入机构进入内核，按系统调用号查入口地址表，转去执行相应处理程序。
  \item \term{结束阶段}：处理程序完成后恢复现场，并把返回值或返回参数送回指定位置，再返回用户态。
\end{points}
\plain{调用过程本质就是“先准备参数，再按编号跳到内核处理，最后带着结果回来”；考试画流程图时通常按这三段展开就够了。}
\exam{流程题常考参数何时传入、入口地址如何确定、返回前为什么要恢复现场。}

\concept{系统调用与过程调用的区别}
\begin{points}
  \item \term{进入方式不同}：系统调用通过陷入/中断机构进入内核；普通过程调用直接按调用指令转入子程序。
  \item \term{运行状态不同}：系统调用在内核态执行；普通过程调用通常仍在用户态执行。
  \item \term{功能对象不同}：系统调用请求 OS 提供受保护的系统服务；过程调用主要是在同一程序内部复用已有功能。
\end{points}
\plain{不要把系统调用理解成“功能更强的函数调用”；它和普通过程调用最本质的差别，是会借助陷入机制跨过用户态/内核态边界。}
\exam{选择题常直接比较两者是否引起状态切换、是否通过中断机构进入、是否能直接执行特权操作。}


\section{操作系统结构}
\concept{常见结构}
\begin{points}
  \item \term{简单结构}：系统各部分边界不清，效率可能高，但维护困难。
  \item \term{层次结构}：从底层硬件到上层服务逐层构造，结构清晰，但层次划分和跨层调用可能带来开销。
  \item \term{微内核结构}：内核只保留最基本功能，如进程通信、低级地址空间管理、基本调度；文件系统、设备驱动等尽量放到用户态服务中。
  \item \term{模块化结构}：核心内核加可加载模块，兼顾性能与可扩展性。
  \item \term{虚拟机结构}：虚拟机监控器在一台物理机上抽象出多台虚拟计算机，每台虚拟机可运行自己的 OS。
\end{points}
\plain{OS 结构是在“好维护”和“高效率”之间权衡：简单结构快但乱，层次清楚但可能慢，微内核可靠但通信开销大。}
\exam{常考宏内核/微内核/模块化优缺点。}


\section{典型操作系统}
\concept{典型操作系统类别}
\begin{points}
  \item \term{批处理操作系统}强调吞吐量和资源利用率，适合大量无交互作业集中处理。
  \item \term{分时操作系统}强调多用户交互和较快响应，使多个终端用户感觉自己独占机器。
  \item \term{实时操作系统}强调在截止期限内完成响应，追求可靠性和可预测性。
  \item \term{个人计算机操作系统}强调易用性、图形交互和通用应用支持，如 Windows、macOS。
  \item \term{网络、分布式、嵌入式操作系统}分别强调联网资源共享、跨节点协同和专用场景控制。
\end{points}
\plain{这一节不是让你背品牌，而是分清不同 OS 的主目标到底是吞吐量、交互性、实时性还是专用控制。}
\exam{常考“某类操作系统最看重什么”，或给场景判断更适合哪类操作系统。}

\chapter{相关硬件基础}

\section{中央处理器}
\concept{CPU 的基本组成与工作流程}
\begin{points}
  \item CPU 的主要部件包括控制器、运算器、寄存器和高速缓存。
  \item CPU 的任务是执行指令，基本流程为：取指 $\rightarrow$ 译码 $\rightarrow$ 取操作数 $\rightarrow$ 执行指令。
  \item 影响 CPU 性能的因素包括主频、字长、指令集、高速缓存、核心数量、总线速度和处理技术。
\end{points}
\plain{这里先抓 CPU 怎样一条条执行指令，后面理解中断、系统调用、上下文切换都要靠这条执行主线。}
\exam{常考取指、译码、执行的顺序，以及 CPU 性能相关因素的辨析。}


\concept{单处理器与多处理器}
\begin{points}
  \item 单处理器系统只有一个运算处理器，同一时刻只能有一个进程在 CPU 上运行。
  \item 单 CPU 中，进程与进程只能并发不能并行；处理机与设备、设备与设备在硬件支持下可以并行工作。
  \item 多处理器系统有多个运算处理器，可分为对称多处理 SMP 和非对称多处理 AMP。
  \item SMP 中各处理器地位相同，共享内存和 I/O 系统；AMP 中主从处理器角色不同，主处理器负责调度与管理。
\end{points}
\plain{单处理器重点分清“并发但不并行”，多处理器才可能多个进程真的同时在不同 CPU 上跑。}
\exam{常考单处理机和多处理机各能实现什么并行关系，以及 SMP、AMP 的区别。}


\concept{Flynn 分类}
\begin{points}
  \item SISD：单指令流单数据流，传统顺序单处理器。
  \item SIMD：单指令流多数据流，一条指令同时处理多个数据，如向量处理、GPU 中的某些模式。
  \item MISD：多指令流单数据流，实际很少见。
  \item MIMD：多指令流多数据流，多处理器/多核系统常见。
\end{points}
\plain{Flynn 分类本质是在看“几条指令流、几条数据流”同时工作，不要和单核多任务混在一起。}
\exam{常考 SISD、SIMD、MIMD 的定义和典型代表。}


\concept{多处理器、多核与超线程}
\begin{points}
  \item 多处理器通常指在一个系统中放置多个 CPU 芯片；多核通常指在同一块 CPU 芯片内集成多个执行核。
  \item 多核结构通常比多 CPU 更紧凑，成本更低、功耗更小。
  \item 超线程是在一个物理处理器内部提供多个逻辑处理器/线程执行环境，共享执行部件和缓存，让处理器更充分利用流水线和功能部件。
  \item 直觉：多核更像“真的多套核心”；超线程更像“一套核心尽量同时喂多条线程”。
\end{points}
\plain{题目若把多核和超线程放一起，抓“多核是真多执行核，超线程是共享执行部件的多个逻辑执行环境”就够了。}
\exam{常考多处理器、多核、超线程的区别，以及谁是真正增加了完整执行核心。}


\concept{操作系统与多核处理器}
\begin{points}
  \item 多核处理器要求操作系统在处理器通信、数据共享、存储层次管理、同步、调度优化和能耗管理等方面提供支持。
  \item 核数增加后，系统不只是“多几个 CPU”，还要处理缓存一致性、线程同步、负载均衡等问题。
  \item 因此 OS 需要改进调度策略、并行执行模型和共享数据管理机制。
\end{points}
\plain{多核不是硬件自己多加几个核就结束了，OS 还得负责把线程分好、数据管好、同步做好。}
\exam{常考操作系统为何要针对多核做调度、同步和共享支持。}


\section{寄存器、特权指令与处理器状态}
\concept{寄存器}
\begin{points}
  \item 寄存器是 CPU 内部高速存储单元，用于保存指令、地址、数据和状态信息。
  \item 常见寄存器包括程序计数器 PC、指令寄存器 IR、通用寄存器、栈指针 SP、程序状态字 PSW。
  \item 上下文切换时，操作系统需要保存当前进程的寄存器现场，并恢复下一个进程的现场。
\end{points}
\plain{寄存器是最直接的执行现场；只要发生中断或进程切换，OS 首先得保存和恢复这些内容。}
\exam{常考 PC、IR、SP、PSW 等寄存器的作用，以及哪些内容属于处理机现场。}


\concept{寄存器的功能分类}
\begin{points}
  \item 以 x86 为例，寄存器可分为通用寄存器、指针及变址寄存器、段寄存器、指令指针与标志寄存器、控制寄存器等。
  \item 通用寄存器常用于普通算术逻辑和数据暂存；栈指针、基址寄存器等常用于过程调用和栈管理。
  \item 段寄存器、控制寄存器、标志寄存器更偏向地址变换、保护和处理器控制。
  \item 外设控制器也可能带有设备专用寄存器，但它们不等同于 CPU 内部通用寄存器。
\end{points}
\plain{这类题通常不是让你死背所有寄存器名，而是分清哪些寄存器管普通数据，哪些管控制、状态、地址。}
\exam{常考按功能给寄存器分类，或判断某寄存器更偏现场信息还是控制信息。}


\concept{特权指令与非特权指令}
\begin{points}
  \item 非特权指令可由用户程序直接执行，如普通算术、逻辑、数据移动指令。
  \item 特权指令只能由内核执行，如 I/O 指令、关中断、修改内存保护寄存器、修改页表、启动 DMA、设置定时器。
  \item 判断标准：凡是可能影响系统安全、影响其他进程、直接控制硬件资源的操作，通常属于特权操作。
\end{points}
\plain{普通程序不能直接碰全局危险资源；一旦需要 I/O、页表、中断等操作，就要进内核态。}
\exam{常考哪些指令是特权指令、系统调用是否引起模式切换、模式切换和进程切换区别。}


\concept{特权指令典型操作}
\begin{points}
  \item 典型特权操作包括设置定时器、读取时钟、清除内存、发起陷入、关中断、修改设备状态、切换用户态与内核态、访问 I/O 设备等。
  \item 这类操作一旦被普通用户程序任意执行，就可能破坏调度、公平性、保护机制和设备控制。
  \item 所以题目判断时，只要操作能直接改变系统控制权或全局硬件状态，通常就应归入特权操作。
\end{points}
\plain{判断特权指令时别死记列表，直接想：这事如果让普通程序自己干，会不会把整个系统搞乱。}
\exam{常考举例判断某操作是否属于特权操作，并说明原因。}


\concept{处理器状态：用户态与内核态}
\begin{points}
  \item 用户态下执行普通用户程序，只能执行非特权指令，不能直接访问受保护资源。
  \item 内核态下执行操作系统内核，可执行特权指令，访问核心数据结构。
  \item 从用户态到内核态的常见原因：系统调用、异常、外部中断。
  \item 从内核态返回用户态：中断/系统调用处理完成后，恢复用户进程现场并执行返回指令。
\end{points}
\plain{用户态和内核态本质是硬件给 OS 划出的权限边界，危险操作必须越过这条边界让内核代办。}
\exam{常考哪些事件会从用户态进入内核态，以及模式切换是否一定伴随进程切换。}


\concept{保护级别与 Ring}
\begin{points}
  \item 一些处理器不只区分两态，还可支持多个保护级别。
  \item 例如 Intel Pentium 体系结构有 4 个保护级别，Ring 0 权限最高，Ring 3 权限最低。
  \item 内核通常运行在 Ring 0，普通用户程序通常运行在 Ring 3。
  \item 这本质上仍是在做权限隔离，只是比“用户态/内核态”描述得更细。
\end{points}
\plain{考试里看到 Ring，直接把它理解成“更细粒度的权限圈层”；Ring0 最强，Ring3 最弱。}
\exam{常考 Ring0 和 Ring3 分别对应什么，以及它和用户态/内核态的关系。}


\concept{处理器状态切换步骤}
\begin{points}
  \item 用户态进入内核态的原因主要有两类：程序主动发起系统调用，或程序运行中发生中断/异常。
  \item 常见切换步骤包括：保存现场、根据中断号或陷入入口设置程序计数器、交换或修改 PSW、转入对应处理程序。
  \item 处理完成后，再恢复现场并返回原程序或转入调度程序。
\end{points}
\plain{把状态切换看成“先保现场，再找入口，再换权限去处理”，答题就不容易漏步骤。}
\exam{常考用户态切到内核态的触发原因，以及状态切换的基本顺序。}


\concept{PSW 程序状态字}
\begin{points}
  \item PSW 保存 CPU 的运行状态信息，包括条件码、中断屏蔽位、处理器模式位等。
  \item 模式位用于区分用户态和内核态，是硬件支持操作系统保护机制的重要基础。
  \item 中断发生时，硬件通常会保存 PC/PSW 等关键信息，以便中断处理完成后回到原程序继续执行。
\end{points}
\plain{PSW 可以看成 CPU 当前“运行身份和状态”的摘要，模式位、中断位都在这里体现。}
\exam{常考 PSW 的作用，以及它为什么是中断返回和现场恢复的重要内容。}


\concept{x86 总体架构中的关键寄存器}
\begin{points}
  \item 指令指针寄存器 EIP 保存下一条待执行指令的地址，常被 \code{CALL}、\code{RET}、\code{JMP} 等控制转移指令修改。
  \item 控制寄存器中常见位有保护模式位和分页相关位，它们决定处理器是否启用保护和分页机制。
  \item 与内存管理相关的寄存器还包括 GDTR、IDTR、TR 等，分别和全局描述符、中断描述符、任务状态保存等机制相关。
  \item 这些细节本章只需知道“它们服务于保护、分页、中断和任务切换”，不必按体系结构课深度去背位级细节。
\end{points}
\plain{x86 这部分不是要你变成体系结构专家，而是知道哪些寄存器在支撑 OS 的保护、中断和任务切换。}
\exam{常考 EIP、控制寄存器、GDTR/IDTR/TR 大概各管什么。}


\section{中断处理}
\concept{中断的基本概念}
\begin{points}
  \item 中断是 CPU 暂停当前程序执行，转而执行相应处理程序，处理完成后再返回原程序的机制。
  \item 中断使 CPU 不必忙等 I/O 完成，是实现多道程序设计和 CPU/I/O 并行工作的关键。
  \item 中断来源包括外部设备中断、定时器中断、异常和系统调用陷入。
\end{points}
\plain{答题时把步骤串起来：提出请求、保存现场、查中断向量、执行服务程序、恢复现场。}
\exam{常考中断响应流程排序题，或简答中断处理的基本步骤。}

\concept{中断分类}
\begin{points}
  \item \term{异步中断}通常来自 CPU 外部，如定时器中断、I/O 设备中断，与当前正在执行的指令本身无直接因果关系。
  \item \term{同步中断}通常来自 CPU 内部，也常称异常，由当前指令执行直接引起。
  \item 在不少体系结构表述中，系统调用也通过一种特殊异常/陷入机制实现。
  \item 因而教材里有时把“中断”狭义地只指异步外部中断，有时又把中断与异常统称为广义中断。
\end{points}
\plain{这块最容易混的是“中断”这个词有狭义和广义两层用法；做题时先看它是在按来源分，还是在把异常一起总称。}
\exam{常考定时器/I/O 属于异步中断，除零/缺页属于同步异常，以及系统调用在实现上通常归入哪一类。}


\concept{中断处理过程}
\begin{points}
  \item 设备或内部事件提出中断请求。
  \item CPU 在适当时刻响应中断，保存断点和现场，切换到内核态。
  \item 根据中断向量找到对应的中断服务程序。
  \item 中断服务程序处理中断事件，例如读取设备状态、唤醒等待进程。
  \item 恢复被中断程序现场，返回原程序或转入调度程序。
\end{points}
\plain{中断就是硬件或软件打断 CPU 当前执行，让 OS 先处理更紧急或必须处理的事件。}
\exam{常考中断处理步骤、中断与异常/系统调用区别、为什么中断能提高 CPU 与 I/O 并行度。}

\concept{中断装置与中断处理程序}
\begin{points}
  \item 发现中断源、仲裁优先级并引导 CPU 转入相应服务程序的硬件机制，可统称为\term{中断装置}。
  \item 当多个中断同时请求时，中断装置要先识别中断源并按优先级选择先响应者。
  \item \term{中断处理程序}负责具体处理中断事件，如读取设备状态、完成数据收尾、唤醒等待进程，并在结束后恢复现场。
  \item 因而可以把中断装置理解为“把 CPU 正确带到入口”，把中断处理程序理解为“真正把事情办完”。
\end{points}
\plain{考试里不要把硬件发现中断源和软件处理中断事件混成一步：前者偏“找到谁在叫、跳到哪去”，后者偏“具体怎么收尾”。}
\exam{若简答题问中断装置或中断处理程序的作用，答题时最好把“识别中断源/优先级选择/转入入口”和“处理事件/恢复现场”分开写。}


\concept{中断与异常、系统调用}
\begin{points}
  \item 外部中断通常来自 CPU 外部设备，如磁盘、网卡、键盘。
  \item 异常来自当前指令执行过程，如除零、非法指令、缺页异常。
  \item 系统调用是用户程序有意执行陷入指令，请求内核服务。
  \item 三者都会导致从用户态进入内核态，但触发原因不同。
\end{points}
\plain{三者都会陷入内核，但来源不同：外部设备来的是中断，当前指令出错是异常，程序主动请求服务是系统调用。}
\exam{常考中断、异常、系统调用的来源、同步异步性和是否由程序主动发起。}


\concept{中断优先级与多重中断}
\begin{points}
  \item 当多个中断请求同时到来时，系统通常按预先规定的优先级决定先响应谁。
  \item 高优先级中断可在一定条件下打断低优先级中断服务程序，这就形成多重中断。
  \item 是否允许中断嵌套，与处理中断时是否重新开中断、硬件优先级控制和现场保护能力有关。
  \item 多重中断提高了系统对紧急事件的响应能力，但也增加了中断处理和现场保存恢复的复杂度。
\end{points}
\plain{多重中断可以理解成“处理中断时又被更紧急的中断插队”，前提是系统能分清优先级并正确保存多层现场。}
\exam{PPT 单独列了这节，常考为什么需要中断优先级、什么情况下会出现中断嵌套，以及它带来的代价。}


\concept{同步与原子性}
\begin{points}
  \item \term{同步}强调多个事件或操作之间存在先后配合关系，必须按某种时序推进。
  \item \term{原子性}强调一个操作执行期间不可被打断，要么全部完成，要么完全不做。
  \item 在操作系统里，关中断、硬件原语、原子读改写指令等常用于保证关键操作的原子性。
  \item 原子性是后续实现互斥和同步机制的重要硬件基础，但“同步”和“原子”不是同一个概念。
\end{points}
\plain{同步管“先后顺序要配合”，原子性管“中间不能被拆开”；两个词经常一起出现，但问的不是一回事。}
\exam{这节虽然在硬件基础里，但会直接服务后面进程同步章节，常考同步与原子性的区别以及为什么原子指令对互斥实现重要。}


\section{操作系统的硬件支持}
\concept{内存保护}
\begin{points}
  \item 为防止进程访问非法内存，硬件需要支持地址变换和地址合法性检查。
  \item 简单系统中可用基址寄存器和界限寄存器检查地址是否越界。
  \item 分页/分段系统中由 MMU 根据页表或段表完成地址映射和保护检查。
\end{points}
\plain{内存保护的核心是“一个进程不能随便碰另一个进程或内核的地址空间”，否则多道程序根本不安全。}
\exam{常考基址/界限寄存器或页表、段表如何支持地址检查与保护。}


\concept{定时器}
\begin{points}
  \item 定时器用于防止用户程序长期独占 CPU。
  \item 操作系统设置时间片，到时后定时器中断发生，CPU 转入内核执行调度程序。
  \item 分时系统和抢占式调度离不开定时器支持。
\end{points}
\plain{定时器是抢占式 OS 的底盘，没有时钟中断，系统就很难防止用户程序长期霸占 CPU。}
\exam{常考定时器为什么是分时系统和抢占式调度的硬件基础。}


\concept{I/O 保护}
\begin{points}
  \item 用户程序不能直接执行 I/O 指令，必须通过系统调用请求内核完成。
  \item 这样可以统一设备分配、访问权限、缓冲、错误处理，防止多个进程直接抢设备。
  \item 因此 I/O 指令属于典型特权指令。
\end{points}
\plain{I/O 保护强调的是“普通进程不能直接发设备命令”，否则设备分配、权限和安全性都会失控。}
\exam{常考为什么 I/O 指令是特权指令，以及用户程序怎样合法发起 I/O。}


\section{计算机启动与虚拟机}
\concept{x86 实模式与保护模式}
\begin{points}
  \item 8086 是 16 位处理器，但地址总线为 20 位，因此通过“段基址 + 偏移”完成寻址，实模式地址可写为 $\text{physical addr}=16\times \text{segment}+\text{offset}$。
  \item 80386 及后续处理器为兼容早期体系结构，通常先从实模式启动，再切换到保护模式。
  \item 保护模式支持更强的地址保护、权限管理和更大的寻址能力，是现代操作系统正常运行的基础。
  \item 直觉：实模式更像早期裸奔方式，保护模式才真正支持 OS 要的保护与管理。
\end{points}
\plain{实模式重点看“段基址+偏移”的老式寻址；保护模式重点看“支持现代 OS 的保护机制”。}
\exam{常考为什么 x86 要先实模式启动再切到保护模式，以及实模式的基本寻址公式。}


\concept{计算机启动过程}
\begin{points}
  \item 上电后 CPU 从固定地址开始执行固件中的引导程序。
  \item 固件进行硬件自检和基本初始化，找到可启动设备。
  \item 引导程序把操作系统内核装入内存。
  \item 内核初始化数据结构、设备驱动、内存管理、进程管理，创建第一个用户态进程或系统服务。
\end{points}
\plain{开机不是直接进 OS，而是固件先跑，找到引导程序，再把 OS 内核装入内存。}
\exam{常考 BIOS/UEFI、引导程序、内核加载的大致顺序。}


\concept{虚拟机}
\begin{points}
  \item 虚拟机把底层物理机器抽象成多台逻辑机器，使多个 OS 能在同一物理平台上运行。
  \item 虚拟机监控器负责 CPU、内存、设备等资源的虚拟化与隔离。
  \item 直觉：普通 OS 把裸机虚拟成“适合应用使用的机器”；虚拟机监控器把裸机虚拟成“多台可运行 OS 的机器”。
\end{points}
\plain{虚拟机和“OS 的虚拟资源”不是一回事，这里是再加一层软件把一台物理机伪装成多台可装 OS 的机器。}
\exam{常考虚拟机监控器和普通操作系统的区别，以及虚拟机的基本作用。}

\chapter{进程}

\section{进程引入}
\concept{前趋图}
\begin{points}
  \item 前趋图是有向无循环图，用于描述程序段、语句或进程之间的先后执行关系。
  \item 若 $P_i \rightarrow P_j$，表示 $P_i$ 必须在 $P_j$ 开始之前完成。
  \item 没有前趋的结点称为初始结点，没有后继的结点称为终止结点。
  \item 做题时根据箭头判断哪些操作必须先完成，哪些操作可以并发。
\end{points}
\plain{用有向无环图画出“谁必须在谁之前完成”，目的是识别哪些操作能并发。}
\exam{常考给前趋关系画图、数可并发的语句、写同步信号量约束。}


\concept{程序的顺序执行}
\begin{points}
  \item 顺序执行指一个程序中的多个操作严格按照规定顺序执行，例如输入 $I \rightarrow$ 计算 $C \rightarrow$ 输出 $P$。
  \item 特征一：顺序性，每一步必须在下一步开始前完成。
  \item 特征二：封闭性，程序运行结果不受外界因素影响。
  \item 特征三：可再现性，只要初始条件和执行环境相同，每次运行结果相同。
\end{points}
\plain{一个程序自己按固定步骤走，外界不插手，所以结果可重复。}
\exam{常考顺序性、封闭性、可再现性三个特征。}


\concept{程序的并发执行}
\begin{points}
  \item 并发执行指多个程序或程序段在时间上重叠推进，一个程序尚未结束，另一个程序已经开始。
  \item 特征一：间断性，程序表现为“执行 - 暂停 - 执行”。
  \item 特征二：失去封闭性，多个程序共享资源，资源状态可能被他人改变。
  \item 特征三：不可再现性，在初始条件相同的情况下，执行顺序不同也可能导致结果不同。
\end{points}
\plain{并发执行这节重点不在和“并行”抠定义，而在看出它为什么会带来间断性、失去封闭性和结果不可再现。}
\exam{常考顺序执行和并发执行的三组特征对比，尤其为什么并发环境下会出现与时间有关的错误。}

\section{进程定义与描述}
\concept{进程的定义}
\begin{points}
  \item 进程是程序的一次执行过程，是系统进行资源分配和调度的基本单位之一。
  \item 进程不仅包含程序代码，还包含数据、堆栈、打开文件、寄存器现场、程序计数器和 PCB 等运行信息。
  \item 进程具有动态性、并发性、独立性、异步性和结构性。
\end{points}
\plain{这一定义里最关键的是“执行过程”和“资源分配调度单位”两层，不要把进程只理解成正在跑的代码片段。}
\exam{常考程序与进程区别、进程实体由程序段/数据段/PCB 构成、进程是资源分配基本单位。}


\concept{进程的特征}
\begin{points}
  \item \term{动态性}：进程是程序的一次执行过程，有创建、运行、暂停、撤销等生命周期；程序本身是静态实体。
  \item \term{并发性}：多个进程可同时存在于内存中，在一段时间内并发推进。
  \item \term{独立性}：进程是能独立运行、独立获得资源、独立接受调度的基本单位。
  \item \term{异步性}：进程按各自独立、不可预知的速度推进，因此执行结果可能受调度次序影响。
  \item \term{结构性}：进程实体通常由程序段、数据段和 PCB 等部分组成，也就是进程映像。
\end{points}
\plain{这五个特征里最容易混的是“动态性”和“并发性”：前者强调有生命周期，后者强调多个进程能重叠推进。}
\exam{简答题经常直接问“进程有哪些特征”，最好按这五点完整列出，而不是只答动态性和并发性。}


\concept{程序与进程的区别}
\begin{points}
  \item 程序是静态的代码文件，进程是动态的执行活动。
  \item 一个程序可以对应多个进程，例如多个用户同时打开同一个编辑器。
  \item 一个进程在运行中可能执行多个程序片段，例如通过 \code{exec} 装入新程序。
  \item 程序没有生命周期状态，进程有创建、就绪、运行、阻塞、终止等状态。
\end{points}
\plain{区分程序和进程时，最稳的是按“静态/动态、是否拥有生命周期、是否携带资源和状态”三条线来答。}
\exam{常考判断题：同一程序能否对应多个进程、进程能否在运行中装入新程序、程序本身是否具有状态。}


\concept{PCB 进程控制块}
\begin{points}
  \item PCB 是操作系统感知和管理进程的唯一标志。
  \item PCB 保存进程标识信息、处理机状态、调度信息、进程控制信息、内存管理信息、I/O 状态信息等。
  \item 上下文切换本质上是保存当前进程 PCB 中的现场，并从下一个进程 PCB 中恢复现场。
  \item 选择题直觉：问“进程存在的唯一标志”通常答 PCB，不是程序、PID 或页表。
\end{points}
\plain{PCB 就是 OS 给每个进程建的档案，没有 PCB，OS 就没法调度和恢复这个进程。}
\exam{常考 PCB 内容、PCB 的作用、进程存在的唯一标志。}


\concept{PCB 的组织方式}
\begin{points}
  \item 为了便于创建、撤销、阻塞、唤醒和调度进程，操作系统需要把大量 PCB 组织起来。
  \item \term{线性方式}：把 PCB 顺序存放在连续内存中，实现简单，但插入删除不够灵活。
  \item \term{链接方式}：把同一状态的 PCB 按链表连接，便于按状态组织就绪队列、阻塞队列。
  \item \term{索引方式}：建立索引表，再由索引项指向对应 PCB，便于按标识快速定位。
\end{points}
\plain{PCB 不只是一个个单独存在，OS 还得把它们按状态串起来或索引起来，否则调度器没法高效找到“下一个该处理谁”。}
\exam{PPT 专门讲了 PCB 组织方式，常考线性/链接/索引三种方法的适用特点，尤其链表为什么适合状态队列。}


\concept{Bernstein 条件}
\begin{points}
  \item 为判断两个程序段能否并发执行且不产生与时间有关的错误，可引入读集 $R(S_i)$ 和写集 $W(S_i)$。
  \item 若满足 $R(S_i)\cap W(S_j)=\varnothing$，$R(S_j)\cap W(S_i)=\varnothing$，$W(S_i)\cap W(S_j)=\varnothing$，则 $S_i$ 与 $S_j$ 可安全并发。
  \item 前两个条件避免“读到被对方修改中的数据”，第三个条件避免“两个程序段写同一数据导致结果丢失”。
  \item 满足这三个条件时，并发执行仍可保持封闭性和可再现性。
\end{points}
\plain{Bernstein 条件本质是在问：两段程序会不会互相读坏、写坏对方要用的数据。}
\exam{常考给出若干语句的读集写集，判断哪几条可以并发执行。}


\concept{\code{cobegin} 与 \code{coend} 描述并发}
\begin{points}
  \item 并发语句常写作 \code{cobegin S1; S2; ...; Sn; coend}。
  \item 其含义是多个语句段在逻辑上可并发执行，但 \code{coend} 之后的语句必须等这些并发段都结束后才能继续。
  \item 它常与前趋图一起使用，用于描述程序并发结构。
\end{points}
\plain{\code{cobegin} 像“这些分支一起干”，\code{coend} 像“所有分支会合后再往下走”。}
\exam{常考把前趋图改写成 \code{cobegin/coend} 形式，或反过来画前趋关系。}

\concept{为什么引入进程}
\begin{points}
  \item 引入进程是为了刻画程序执行过程的动态性，而不再只把程序看成静态指令集合。
  \item 引入进程有助于发挥系统并发性，提高资源利用率。
  \item 引入进程还用于处理共享性问题，正确描述程序执行状态及系统对运行实体的管理。
\end{points}
\plain{只有“程序”这个静态概念时，很难描述并发推进、资源共享和状态变化，所以 OS 需要再引入“进程”这个动态执行实体。}
\exam{若问“操作系统为什么引入进程”，答题可按“刻画动态性、支持并发、处理共享与状态管理”三点展开。}


\concept{进程映像}
\begin{points}
  \item 进程映像是某时刻进程的内容及其状态的集合，也可理解为进程在系统中的完整存在形态。
  \item 通常包括 PCB，以及进程在虚拟地址空间中的代码段、数据段、栈、堆等内容。
  \item 某些教材还会强调核心栈，用于保存进程在核心态工作时的现场保护。
\end{points}
\plain{进程映像回答的是“这个进程完整长什么样”，不只是代码，还包括数据和管理信息。}
\exam{常考进程映像由哪些部分组成，以及它和程序本身的区别。}


\concept{进程上下文}
\begin{points}
  \item 进程物理实体及其支持运行的环境合称为进程上下文，进程运行可看作在其上下文中执行。
  \item 进程映像强调“全部内容集合”；上下文更强调“当前运行现场和支撑环境”。
  \item 当系统调度新进程占有处理器时，就会发生上下文切换。
  \item 上下文可分为用户级上下文、寄存器级上下文和系统级上下文。
\end{points}
\plain{进程映像偏“全家福”，进程上下文偏“此刻跑到哪、现场是什么样”。}
\exam{常考进程映像和进程上下文的区别，以及上下文切换时保存恢复什么。}


\concept{进程上下文的组成}
\begin{points}
  \item \term{用户级上下文}：程序正文、数据、共享存储区和用户栈等，属于进程虚拟地址空间中的用户态部分。
  \item \term{寄存器级上下文}：程序计数器、程序状态字、栈指针、通用寄存器、控制寄存器等现场信息。
  \item \term{系统级上下文}：PCB、内核栈、页表等主存管理信息、I/O 状态、文件描述符表等内核管理信息。
  \item 上下文切换时，并不是只换一组寄存器，而是按需要切换这些和运行现场直接相关的内容。
\end{points}
\plain{用户级上下文偏“进程自己的地址空间内容”，寄存器级上下文偏“CPU 此刻现场”，系统级上下文偏“内核为这个进程维护的管理环境”。}
\exam{这类题常考“上下文切换保存什么”，按用户级/寄存器级/系统级三层去答最稳。}


\section{进程状态与转换}
\concept{三状态模型}
\begin{points}
  \item \term{就绪态}：进程已具备运行条件，只差 CPU。
  \item \term{运行态}：进程正在 CPU 上执行。
  \item \term{阻塞态/等待态}：进程因等待某事件不能运行，即使 CPU 空闲也不能执行。
  \item 状态转换：就绪 $\rightarrow$ 运行由调度引起；运行 $\rightarrow$ 就绪由时间片到或被抢占引起；运行 $\rightarrow$ 阻塞由等待 I/O/资源引起；阻塞 $\rightarrow$ 就绪由等待事件完成引起。
\end{points}
\plain{三态模型先抓最核心的三种活动状态：能不能运行、正在不在运行、是不是在等事件。}
\exam{常考一个处理器上运行态进程最多几个，以及三态之间的转换条件。}


\concept{五状态模型}
\begin{points}
  \item 在三状态基础上加入新建态和终止态。
  \item 新建态表示进程正在创建，尚未进入就绪队列。
  \item 终止态表示进程执行结束或被撤销，等待系统回收资源。
  \item 创建进程通常涉及分配 PCB、建立地址空间、装入程序、初始化资源、进入就绪队列。
\end{points}
\plain{五态就是在三态前后各补一个“进来”和“出去”的阶段，把完整生命周期补齐。}
\exam{常考五态模型比三态多了哪两个状态，以及创建、撤销分别对应什么阶段。}


\concept{挂起状态}
\begin{points}
  \item 当内存紧张或系统需要调节多道程序度时，可把进程换出到外存，形成挂起状态。
  \item 就绪挂起：进程不在内存，但一旦调入内存即可运行。
  \item 阻塞挂起：进程不在内存，且还在等待某事件。
  \item 中级调度/交换调度常与挂起状态相关。
\end{points}
\plain{挂起强调“暂时不在内存里”，它和阻塞不是一回事；一个进程既可能阻塞，也可能同时被挂起。}
\exam{常考为什么引入挂起状态，以及就绪挂起、阻塞挂起的区别。}

\concept{挂起进程的特征}
\begin{points}
  \item 挂起进程可视为暂时不在主存中的进程，在被重新调入主存前不参与处理机调度。
  \item 挂起进程不能立即执行，即使它原来已经具备运行条件也不行。
  \item 挂起进程可能仍在等待某个事件，但该事件结束并不自动解除挂起状态。
  \item 进入挂起常由操作系统、父进程或进程自身发起；结束挂起通常需由操作系统或父进程发出激活命令。
\end{points}
\plain{判断“挂起”时，关键看它是否已经被换出主存、被暂时禁止参与调度，而不只是看它有没有在等事件。}
\exam{区分“等待”和“挂起”时，要答出：等待针对事件，挂起针对是否驻留主存以及是否允许参与调度。}


\section{进程组织、控制与管理}
\concept{进程队列}
\begin{points}
  \item 就绪队列：所有驻留内存且处于就绪状态、等待 CPU 的进程集合。
  \item 设备队列：等待某个 I/O 设备的进程集合。
  \item 作业队列：外存上等待进入内存的作业集合。
  \item 进程在生命周期中会在多个队列之间迁移，调度程序就是从这些队列中选择合适对象。
\end{points}
\plain{队列是 OS 组织进程的方式：就绪的排就绪队列，等设备的排设备队列。}
\exam{常考进程在生命周期中如何在不同队列迁移。}


\concept{进程控制原语}
\begin{points}
  \item 进程控制的常见功能包括创建、撤销、阻塞、唤醒、挂起和激活。
  \item 这些功能通常由内核原语实现，原语由若干机器指令构成，执行期间不可分割。
  \item 原语必须保证原子性，否则在修改 PCB、状态和队列时可能被中断，导致系统状态不一致。
\end{points}
\plain{进程控制题先抓住“原语”二字：它是内核里不可分割的基本操作，专门负责改 PCB、改状态、改队列。}
\exam{常考什么叫原语、为什么原语必须原子执行，以及哪些进程管理操作属于控制原语。}


\concept{创建原语}
\begin{points}
  \item 创建进程需要申请唯一标识、分配 PCB、分配必要资源、初始化上下文、插入就绪队列。
  \item 引起创建的常见原因包括系统初始化、用户请求和父进程通过系统调用创建子进程。
  \item 创建后的父子进程可以并发执行，也可以由父进程等待子进程结束后再继续。
  \item 为描述父子关系，可用进程树/进程图表示“谁创建了谁”。
  \item Unix/Linux 中常见系统调用：\code{fork()} 创建子进程，\code{exec()} 装入新程序，\code{wait()} 等待子进程，\code{exit()} 退出进程。
  \item \code{fork()} 返回两次：父进程中返回子进程 PID，子进程中返回 0；失败返回负值。
\end{points}
\plain{创建原语的标准流程就是“要 PCB、配资源、填初值、进就绪队列”，代码题再把它和 \code{fork/exec/wait} 对起来。}
\exam{常考创建原语步骤、进程树、fork/exec/wait/exit，以及多次 \code{fork()} 后的进程数和输出行数。}


\concept{撤销原语}
\begin{points}
  \item 进程撤销的原因可分为正常结束、异常结束和外界干预。
  \item 异常结束常见于越界、非法指令、I/O 故障、内存不足等；外界干预常见于父进程或系统强制终止。
  \item 撤销原语的主要工作是找到对应 PCB，必要时先停止其执行，再回收该进程占有的资源和 PCB。
  \item 题目还可能区分两种策略：只撤销指定进程，或连同其全部子孙进程一起撤销。
\end{points}
\plain{撤销原语不是简单“结束程序”，而是把进程从系统里彻底清掉，包括停执行、收资源、删 PCB。}
\exam{常考进程终止原因分类、撤销原语步骤，以及“撤销某进程是否连带其子孙进程”这类策略题。}


\concept{阻塞原语与唤醒原语}
\begin{points}
  \item 进程因等待系统服务、I/O 完成、合作进程配合或某种事件到来时，可由执行态转为阻塞态。
  \item 阻塞原语的核心动作是停止当前进程执行、保存 CPU 现场、把状态改为阻塞，并插入相应等待队列。
  \item 当等待事件发生后，唤醒原语将该进程移出等待队列，改为就绪态，再插入就绪队列。
  \item 阻塞通常由进程自己触发；唤醒通常由发现事件发生的其他进程或内核触发。
\end{points}
\plain{阻塞和唤醒一定成对理解：自己因为等事件而睡下去，别人或内核在事件到了之后把你放回就绪队列。}
\exam{常考阻塞/唤醒的触发条件、两类原语各自修改哪种队列，以及“阻塞态不能直接变运行态”。}


\concept{挂起原语与激活原语}
\begin{points}
  \item 挂起强调把进程暂时从内存换出到外存，常用于缓解内存压力或调整多道程序度。
  \item 挂起时，活动就绪可变为静止就绪，活动阻塞可变为静止阻塞；若进程正在执行，还需先保护现场并重新调度。
  \item 激活原语的作用与之相反：把挂起进程重新调入内存，使其回到就绪或阻塞的活动状态。
  \item 挂起/激活既可针对指定进程，也可能连同其子孙进程一起处理。
\end{points}
\plain{挂起和阻塞别混：阻塞是“等事件”，挂起是“先不放在内存里”；一个进程可以既阻塞又挂起。}
\exam{常考活动就绪/静止就绪、活动阻塞/静止阻塞之间如何转换，以及挂起与激活对状态和调度的影响。}


\concept{进程切换与模式切换}
\begin{points}
  \item 模式切换是用户态和内核态之间的切换，不一定更换正在运行的进程。
  \item 进程切换是 CPU 从一个进程切换到另一个进程，一定需要内核介入，通常伴随上下文保存和恢复。
  \item 系统调用会导致模式切换；是否导致进程切换，要看是否阻塞或调度。
  \item I/O 请求通常先从用户态陷入内核态，若需要等待 I/O 完成，则当前进程阻塞并发生进程切换。
\end{points}
\plain{模式切换只是换权限；进程切换还要保存/恢复另一个进程的现场，开销更大。}
\exam{常考系统调用是否必然进程切换、进程切换需要保存哪些上下文。}


\section{进程通信}
\concept{进程协作的原因}
\begin{points}
  \item 信息共享：多个进程需要交换或共享感兴趣的数据。
  \item 提高运算速度：把任务分解后并发执行，以缩短总完成时间。
  \item 系统模块化：把复杂系统拆成多个相互协作的模块或进程。
\end{points}
\plain{进程通信存在的意义，不只是“把消息传过去”，而是为了共享信息、并发加速以及支撑系统模块化设计。}
\exam{若问“为什么需要进程通信”，不要只答数据交换，还要补上并发加速和系统模块化。}

\concept{低级通信与高级通信}
\begin{points}
  \item 低级通信主要通过信号量等同步工具传递少量控制信息。
  \item 高级通信用于传递大量数据或结构化消息，常见方式包括共享存储、消息传递、管道。
\end{points}
\plain{共享存储的优势是“数据直接放一块公共内存里”，所以传大块数据最快；代价是同步互斥必须自己另外补上。}
\exam{常考为什么共享存储速度快、为什么它又最容易和信号量/互斥锁联动出题。}

\concept{进程通信的基本形态}
\begin{points}
  \item 最基本的两类 IPC 模型是消息传递模型和共享内存模型。
  \item 消息传递模型通过发送和接收消息交换信息，进程之间不必直接共享地址空间。
  \item 共享内存模型让多个进程直接访问同一片内存区域，数据交换效率高，但必须额外处理同步互斥。
\end{points}
\plain{IPC 题先抓总类：要么靠 OS 传消息，要么大家直接看同一块共享内存。邮箱、消息队列、管道、共享存储，基本都能往这两类里归。}
\exam{若题目先问“IPC 的两种基本模型是什么”，就先答消息传递和共享内存，再展开各自优缺点。}


\concept{共享存储}
\begin{points}
  \item 多个进程共享同一块内存区域，通过读写共享区交换数据。
  \item 优点：速度快，适合大量数据交换。
  \item 缺点：必须另行解决同步互斥问题，防止同时读写导致不一致。
\end{points}
\plain{共享存储传数据最快，但同步问题也最直接暴露出来，所以它常和信号量、互斥一起考。}
\exam{常考共享存储为什么速度快、为什么还必须配合同步机制。}


\concept{消息传递}
\begin{points}
  \item 进程通过发送和接收消息通信，可由操作系统提供消息队列、邮箱等机制。
  \item 直接通信：发送者明确指定接收者。
  \item 间接通信：通过邮箱或消息队列进行。
  \item 优点：同步控制相对清晰，适合分布式系统；缺点是消息复制和内核介入可能带来开销。
\end{points}
\plain{消息传递的核心是由操作系统充当中介，发送方和接收方不必直接共享地址空间，但要付出内核介入和消息复制的代价。}
\exam{常考直接通信 vs 间接通信、阻塞 vs 非阻塞接收发送，以及消息传递为什么更适合分布式环境。}


\concept{管道}
\begin{points}
  \item 管道是一种半双工通信机制，常用于具有亲缘关系的进程之间。
  \item 匿名管道常用于父子进程；命名管道可用于无亲缘关系的进程。
  \item Shell 中的 \code{|} 就是把前一个命令的输出作为后一个命令的输入。
\end{points}
\plain{管道更像一根有容量限制的字节流通道，题目里除了“半双工”外，还常顺带考互斥、同步和写满/读空时谁该阻塞。}
\exam{常考匿名管道和命名管道区别、为什么匿名管道常用于亲缘进程、以及双向通信为何通常要两根管道。}

\chapter{线程}

\section{线程定义}
\concept{为什么引入线程}
\begin{points}
  \item 引入进程是为了描述和实现多道程序并发执行。
  \item 引入线程是为了减少并发执行的时空开销，让同一应用内部的多个任务能更轻量地并发。
  \item 早期进程同时承担“资源拥有单位”和“调度分派单位”两个角色，创建和切换开销较大。
  \item 线程把二者分开：进程更多作为资源拥有单位，线程作为 CPU 调度的基本执行单位。
\end{points}
\plain{线程是进程里面更轻的执行流，同进程线程共享资源，但各自有栈、寄存器和 PC。}
\exam{老师明确点过线程代码题和“为什么引入线程”，常考把进程的两个属性拆开来答。}


\concept{线程的定义}
\begin{points}
  \item 线程是进程内的一个执行单元，是进程内的可调度实体。
  \item 一个线程由线程 ID、程序计数器、寄存器集合和栈组成。
  \item 同一进程内的线程共享代码段、数据段、打开文件、地址空间等资源。
  \item 线程是“轻量级进程”，但不是独立资源拥有者。
\end{points}
\plain{定义题抓两半写：线程自己只带执行现场，资源主体仍挂在所属进程名下。}
\exam{常考线程由哪些部分组成，以及哪些资源是线程私有、哪些是进程内共享。}


\concept{线程私有资源与共享资源}
\begin{points}
  \item 线程私有的典型内容包括线程 ID、程序计数器、寄存器集合和执行栈。
  \item 同一进程内线程共享代码段、数据段、堆、打开文件及多数操作系统资源。
  \item 正因为共享地址空间，线程间通信方便，但一个线程写坏共享数据也更容易影响同进程其他线程。
\end{points}
\plain{记忆法：凡是“执行到哪一步”的现场通常归线程，凡是“这整个程序拥有什么”的资源通常归进程。}
\exam{常考判断题：栈、寄存器、PC 是否共享；代码段、全局数据、打开文件是否共享。}


\concept{进程与线程区别}
\begin{points}
  \item 进程是资源分配的基本单位，线程是 CPU 调度的基本单位。
  \item 同一进程内线程共享地址空间，线程之间通信方便但相互影响更大。
  \item 不同进程地址空间隔离，通信开销较大但保护性更好。
  \item 线程创建、撤销、切换通常比进程轻量，因为不需要切换完整地址空间和大量资源。
\end{points}
\plain{这题别答成“线程更小的进程”就结束，必须把“资源归进程、调度归线程、共享地址空间”三点写出来。}
\exam{常考进程和线程在调度单位、资源拥有、通信开销、切换开销上的对比。}


\concept{多线程优点}
\begin{points}
  \item 响应性更好：例如界面线程响应用户，后台线程执行耗时任务。
  \item 资源共享方便：同一进程内线程天然共享内存和文件。
  \item 经济性更好：创建和切换开销较小。
  \item 可伸缩性更好：多核系统中多个线程可并行执行。
\end{points}
\plain{PPT 这里强调四个关键词：响应度、共享、经济性、可伸缩性，答简答题按这四点展开最稳。}
\exam{常考“多线程有哪些优势”，最好按响应性/共享/经济性/可伸缩性分点写。}


\concept{并发与并行}
\begin{points}
  \item 并发 concurrency 指多个任务在同一时间段内交替推进。
  \item 并行 parallelism 指多个任务在同一时刻真正同时执行，通常依赖多核处理器。
  \item 单核系统通常体现为并发；多核系统既能并发，也能并行。
\end{points}
\plain{并发强调“都在推进”，并行强调“真的同时跑”。}
\exam{常考辨析题：单核只能并发不能真正并行，多核才可能并行。}


\concept{Amdahl 定律}
\begin{points}
  \item 设程序中串行部分比例为 $S$，并行部分比例为 $1-S$，处理器核心数为 $N$。
  \item 理想情况下，总执行时间满足：$T_{new}=S\times T_{old}+(1-S)\times \dfrac{T_{old}}{N}$。
  \item 当 $N$ 趋于无穷大时，加速比上限接近 $1/S$。
  \item 因此串行部分越大，多核带来的性能提升上限越低。
\end{points}
\plain{再多的核也救不了那段必须串行跑的代码，串行部分像木桶最短板。}
\exam{常考给定串行比例和核数算理论加速比，或说明为什么加核数后收益会递减。}


\concept{多核编程的挑战}
\begin{points}
  \item 任务分解：识别哪些子任务能够独立并发执行。
  \item 负载均衡：尽量让多个核心工作量接近，避免一核忙一核闲。
  \item 数据划分：让不同核心处理不同数据块，减少争抢。
  \item 数据依赖与同步：正确处理共享数据和先后约束。
  \item 测试调试：并发程序的错误具有时序相关性，更难复现。
\end{points}
\plain{多核难点不在“开几个线程”，而在“怎么拆、怎么同步、怎么别互相拖后腿”。}
\exam{常考简答：为什么多核程序比单线程程序更难设计和调试。}


\concept{数据并行与任务并行}
\begin{points}
  \item \term{数据并行}：把同一种操作分配到不同数据块上，由多个核心并行处理。
  \item \term{任务并行}：把不同功能的任务分给不同执行流，各自执行不同操作。
  \item 数据并行更强调“同算法处理不同数据”；任务并行更强调“不同任务分工协作”。
\end{points}
\plain{看题目先分辨是在“把一堆数据切开一起算”，还是“把不同工作拆给不同线程做”，前者是数据并行，后者是任务并行。}
\exam{PPT 在多核编程挑战后专门给了这两个概念，常考两者区别或让你判断一个场景属于哪一类并行。}


\section{线程实现}
\concept{用户级线程}
\begin{points}
  \item 用户级线程由用户态线程库管理，内核不知道线程存在。
  \item 优点：创建和切换无需进入内核，速度快；调度策略可由应用定制。
  \item 缺点：一个线程执行阻塞系统调用，可能导致整个进程阻塞；无法真正利用多核并行。
\end{points}
\plain{用户级线程的优势几乎都来自“不进内核”，它的主要短板也正因为“内核根本看不见它”。}
\exam{常考用户级线程为什么快、为什么一阻塞可能整个进程都受影响。}


\concept{内核级线程}
\begin{points}
  \item 内核级线程由操作系统内核管理，调度器直接调度线程。
  \item 优点：一个线程阻塞不影响同进程其他线程；多线程可在多核上并行。
  \item 缺点：线程创建、切换、同步需要内核参与，开销比用户级线程大。
\end{points}
\plain{内核级线程的关键点是“内核真能看见并分别调度每个线程”，所以一个线程卡住时别的线程还能继续。}
\exam{常考为什么内核级线程的阻塞不会拖死同进程其他线程，以及它为什么更能利用多核。}


\concept{为什么内核级线程阻塞不影响其他线程}
\begin{points}
  \item 在内核级线程实现中，处理机分配对象是线程而不是整个进程。
  \item 某个线程因系统调用或 I/O 阻塞时，内核仍能把 CPU 分给同进程的其他就绪线程。
  \item 因此“一个线程阻塞，整个进程全停住”的现象通常发生在用户级线程或多对一模型中，而不是内核级线程中。
\end{points}
\plain{因为内核调度表里登记的是“线程”，不是“这个进程整体”，所以它能跳过被阻塞的那个线程。}
\exam{常考 Why 题：为什么一个内核级线程阻塞后，同进程其他线程仍可运行。}


\concept{用户级线程为什么阻塞整个进程}
\begin{points}
  \item 操作系统看不见用户级线程，内核调度对象仍是整个进程。
  \item 若某个用户级线程执行阻塞系统调用，内核只知道“这个进程阻塞了”。
  \item 因而同一进程中的其他用户级线程也失去运行机会，直到该系统调用返回。
\end{points}
\plain{内核眼里根本没有这些用户线程，所以一个线程卡住时，内核会把整个进程一起挂起。}
\exam{常考用户级线程的典型缺点，尤其与多对一模型一起考。}


\concept{纯用户级线程的调度与抢占限制}
\begin{points}
  \item 纯用户级线程通常由线程库在用户空间完成调度，内核调度对象依然是进程。
  \item 若没有专门机制配合，线程之间通常不具备像内核调度那样的强制抢占能力。
  \item 因此纯用户级线程常依赖线程主动让出处理器，例如执行结束、调用调度函数，或显式调用 \code{thread\_yield()}。
\end{points}
\plain{纯用户级线程更像“进程自己在内部排班”，不主动交棒时，别的用户线程往往很难硬抢过来。}
\exam{常考讨论题：纯用户级线程是否支持抢占、如何切换、为何常需要主动让出 CPU。}


\concept{多线程模型}
\begin{points}
  \item 多对一模型：多个用户线程映射到一个内核线程，管理轻量但阻塞问题明显。
  \item 一对一模型：每个用户线程映射到一个内核线程，并行能力强但内核线程数量受限。
  \item 多对多模型：多个用户线程映射到多个内核线程，兼顾灵活性与并行能力，但实现复杂。
\end{points}
\plain{模型题核心不是背名字，而是看“用户线程和内核线程怎么映射”，再推导阻塞、并行和开销结果。}
\exam{常考多对一/一对一/多对多各自的阻塞特性、能否利用多核、实现代价。}


\concept{多对一模型}
\begin{points}
  \item 多个用户级线程映射到一个内核线程，线程管理主要由用户空间线程库完成。
  \item 优点是创建和管理轻量。
  \item 缺点是一个线程一旦执行阻塞系统调用，整个进程都会阻塞；同时任一时刻只有一个线程能进入内核，无法真正利用多核。
  \item 典型例子包括早期 Solaris Green Threads、GNU Portable Threads。
\end{points}
\plain{多对一最像“很多演员共用一个内核身份”，谁一旦在内核这层卡住，整组人都得等。}
\exam{常考多对一模型为什么阻塞问题最严重，以及为什么它无法发挥多核优势。}


\concept{一对一模型}
\begin{points}
  \item 每个用户线程映射到一个内核线程。
  \item 一个线程阻塞时，其他线程仍可被内核继续调度。
  \item 它支持多线程在多核系统上并发甚至并行运行。
  \item 代价是大量创建内核线程会带来更高系统开销，因此系统通常会限制线程数量。
\end{points}
\plain{一对一的好处最直接，但代价也最直接: 每多一个用户线程，内核里就真多一个调度实体。}
\exam{常考一对一模型相对多对一为什么更适合现代多核系统，以及它的主要成本。}


\concept{多对多模型}
\begin{points}
  \item 多个用户线程映射到少量或相等数量的内核线程。
  \item 用户可创建较多线程，而内核线程可在多处理器系统上并发运行。
  \item 当某个线程执行阻塞系统调用时，内核仍可调度其他内核线程继续推进。
  \item 其变种二层模型还允许部分线程和内核线程做一对一绑定。
\end{points}
\plain{多对多是在用户灵活性和内核并行能力之间折中: 用户线程很多，但不必每个都绑一个内核线程。}
\exam{常考多对多模型如何兼顾用户级线程灵活性和内核级线程并行能力。}


\concept{线程库}
\begin{points}
  \item 线程库向程序员提供创建、管理、同步线程的 API。
  \item 常见线程库包括 POSIX Pthreads、Windows threads、Java threads。
  \item 线程库可以完全在用户态实现，也可以由操作系统提供内核级线程支持。
\end{points}
\plain{线程库是程序员看到的 API 外壳，真正要区分的是“库函数只在用户态跑完”，还是“调用后会进内核”。}
\exam{常考线程库的作用，以及它和用户级线程/内核级线程实现之间的关系。}


\concept{线程库的两种实现方式}
\begin{points}
  \item \term{用户空间线程库}：代码和数据结构都在用户空间，调用 API 后通常只是过程调用，不触发系统调用。
  \item \term{内核支持的线程库}：线程控制依赖操作系统内核，调用 API 后往往会触发系统调用，由内核完成创建、调度和管理。
  \item 前者快但内核不可见，后者功能强且支持真正并行，但模式切换开销更大。
\end{points}
\plain{同样叫“线程库”，实现差别可能很大，考试常让你顺着“是否进内核”推出阻塞、并行能力和开销。}
\exam{常考线程库的两种实现方式，以及它们和用户态/内核态支持的关系。}


\concept{常见线程库与 API}
\begin{points}
  \item Pthreads 是 POSIX 线程标准，常见接口包括 \code{pthread\_create}、\code{pthread\_join}、\code{pthread\_exit}、\code{pthread\_cancel}。
  \item Win32 线程库常见接口包括 \code{CreateThread}、\code{ExitThread}、\code{WaitForSingleObject}。
  \item 这类 API 是代码题高频背景，要能看懂“创建几个线程、谁等待谁、哪些资源共享”。
\end{points}
\plain{考试不要求你背完整手册，但看到这些函数名要立刻知道它们在干什么。}
\exam{老师录音明确点过线程代码题，这里常和输出分析、共享变量分析一起出大题。}


\concept{\code{pthread\_create} 与 \code{pthread\_join}}
\begin{points}
  \item \code{pthread\_create} 用于创建一个新线程，指定线程标识、属性、线程函数和传入参数。
  \item 新线程与创建者属于同一进程，因此通常共享代码段、全局数据、堆和打开文件等资源。
  \item \code{pthread\_join} 用于等待指定线程结束，常用于保证主线程在子线程执行完后再继续。
  \item 代码题里若子线程修改全局变量，主线程在 \code{pthread\_join} 之后通常能看到修改结果；这与 \code{fork()} 后父子进程普通变量互不影响形成对比。
\end{points}
\plain{\code{pthread\_create} 重点看“新线程共享哪些资源”，\code{pthread\_join} 重点看“主线程是不是等它干完再往下执行”。}
\exam{线程代码题高频考点就是创建了几个线程、是否等待、共享变量最终值是多少，以及和 \code{fork()} 代码的差别。}


\section{操作系统中的线程例子}
\concept{Windows 线程}
\begin{points}
  \item Windows 采用一对一模型，每个用户线程对应一个内核线程。
  \item 一个线程通常包含线程 ID、寄存器组、用户栈、内核栈和私有存储区。
  \item 线程相关数据结构常见有 ETHREAD、KTHREAD、TEB。
\end{points}
\plain{Windows 线程是“每建一个用户线程，内核里真有一个对应实体”。}
\exam{常考操作系统例子时，重点不是背缩写全称，而是知道 Windows 基本上是一对一模型。}


\concept{Linux 线程}
\begin{points}
  \item Linux 中常把线程和进程都抽象成 task，线程创建常通过 \code{clone()} 实现。
  \item \code{fork()} 创建的新执行流通常不共享地址空间和大多数资源。
  \item \code{clone()} 可按标志位决定是否共享地址空间、文件表等资源，因此可用来实现线程。
  \item 在类 Unix 系统中，Pthreads 通常基于内核级线程的一对一模型实现。
  \item 区分 \code{fork} 和 \code{clone} 的关键在于“是否共享资源与地址空间”。
\end{points}
\plain{Linux 不是单独造一个完全不同的“线程生物种类”，而是靠共享资源程度不同来区分。}
\exam{常考代码题或概念题：为什么 \code{fork} 后普通变量互不影响，而线程共享全局数据。}


\concept{可重入代码}
\begin{points}
  \item 可重入代码指在被多个执行流同时调用或被中断后再次进入时，仍能得到正确结果的代码。
  \item 它通常不依赖不可保护的全局变量，不在函数内部保留共享静态状态，或对共享状态进行同步保护。
  \item 多线程、信号处理和中断处理程序常要求代码具备可重入性。
  \item 直觉：可重入代码像“每个人带自己的草稿纸”，不会抢同一张草稿纸写中间结果。
\end{points}
\plain{可重入代码像“公共只读函数”，多个线程同时进来也不会互相污染。}
\exam{常考可重入代码特征：不修改自身代码、不依赖全局可变状态或对其加保护。}

\chapter{进程调度}

\section{调度基本概念}
\concept{调度队列}
\begin{points}
  \item 作业队列：外存上等待进入系统的作业集合。
  \item 就绪队列：内存中已就绪、等待 CPU 的进程集合。
  \item 设备队列：等待某个 I/O 设备的进程集合。
  \item 进程在整个生命周期中会在这些队列之间迁移。
\end{points}
\plain{这三个队列对应三种典型等待对象：等进入系统、等 CPU、等设备。只要先看“在等谁”，队列就不容易混。}
\exam{选择题很喜欢把“外存后备作业”“内存就绪进程”“等待 I/O 的阻塞进程”三类对象打乱让你配对。}


\concept{三级调度}
\begin{points}
  \item 高级调度又称作业调度/长程调度，决定哪些作业从外存进入内存，影响多道程序度。
  \item 中级调度又称交换调度，决定哪些进程换入、换出内存，用于缓解内存紧张。
  \item 低级调度又称进程调度/短程调度，决定就绪队列中哪个进程获得 CPU，运行频率最高。
  \item 记忆：高级管“进不进系统”，中级管“在不在内存”，低级管“谁上 CPU”。
\end{points}
\plain{三级调度最容易记混的是“管到哪一层”：高级决定进系统，中级决定留不留内存，低级决定此刻谁拿 CPU。}
\exam{常考三个调度层次的别名和作用范围，尤其中级调度 = 交换调度，别漏。}


\concept{调度要解决的三个问题}
\begin{points}
  \item \term{WHAT}：按什么原则选择下一个要执行的进程，即调度算法本身。
  \item \term{WHEN}：何时进行调度，即调度时机，和抢占式/非抢占式密切相关。
  \item \term{HOW}：怎样把被选中的进程真正放上 CPU，即分派和上下文切换过程。
\end{points}
\plain{这页 PPT 的主线很适合简答题：先说选谁，再说什么时候选，最后说怎么切过去。}
\exam{常考把调度问题拆成 WHAT/WHEN/HOW，或据此解释算法、时机和切换过程的区别。}


\concept{CPU 密集型与 I/O 密集型}
\begin{points}
  \item CPU 密集型作业计算时间多，CPU burst 长，I/O 次数相对少。
  \item I/O 密集型作业 I/O 操作多，CPU burst 短。
  \item 作业调度应合理搭配两类作业：CPU 密集型过多会导致 I/O 设备空闲；I/O 密集型过多会导致就绪队列空、CPU 空闲。
\end{points}
\plain{调度不仅看“谁先跑”，还要看作业类型怎么搭配，才能让 CPU 和 I/O 设备都别闲着。}
\exam{常考 CPU 密集型和 I/O 密集型区别，以及为什么作业调度要混合搭配这两类作业。}


\concept{调度程序与分派程序}
\begin{points}
  \item 调度程序负责从就绪队列中选择下一个运行进程。
  \item 分派程序负责完成上下文切换、切换到用户态、跳转到被选进程的正确位置。
  \item 分派延迟是停止一个进程并启动另一个进程所需时间，应尽可能短。
\end{points}
\plain{调度器负责“选人”，分派器负责“把 CPU 真交过去”。}
\exam{常考 scheduler 和 dispatcher 的分工，以及什么叫分派延迟。}


\concept{上下文切换}
\begin{points}
  \item 上下文切换是保存当前进程状态并恢复另一个进程状态的过程。
  \item 上下文信息通常保存在 PCB 中。
  \item 上下文切换时间属于纯系统开销，不直接创造用户计算结果。
  \item 因此时间片过短、抢占过于频繁时，调度开销会明显上升。
\end{points}
\plain{切换本身不干活，只是在“换人上机”，所以切得太勤会把 CPU 时间浪费在管理上。}
\exam{常考上下文切换保存到哪里、为什么它属于系统开销、与时间片大小有什么关系。}


\section{调度目标与评价指标}
\concept{常见评价指标}
\begin{points}
  \item CPU 利用率：CPU 忙碌时间比例，越高越好。
  \item 吞吐量：单位时间完成的进程数量，越高越好。
  \item 周转时间：进程从提交到完成的总时间，越短越好。
  \item 等待时间：进程在就绪队列中等待 CPU 的总时间，越短越好。
  \item 响应时间：从提交请求到首次响应的时间，分时系统尤其关注。
  \item 公平性：不同用户或进程应获得合理的 CPU 份额，避免长期饥饿；不同系统对“公平”的定义并不完全相同。
\end{points}
\plain{这些指标不是并列硬背，真正的源头通常只有调度时间轴；时间轴一错，后面周转、等待、响应会全错。}
\exam{计算题通常先让你画甘特图，再顺着求完成时间和各指标；只写结果不写过程很容易丢步骤分。}


\concept{周转时间公式}
\begin{points}
  \item 完成时间 $C_i$，到达时间 $A_i$，服务时间 $S_i$。
  \item 周转时间 $T_i=C_i-A_i$。
  \item 带权周转时间 $W_i=\dfrac{T_i}{S_i}$。
  \item 平均周转时间和平均带权周转时间就是对所有进程取平均。
\end{points}
\plain{公式本身不难，难点是先把 $C_i$ 算对；周转和带权周转基本都是从完成时间往后顺推。}
\exam{常见陷阱是把服务时间当周转时间分母之外的别的量，或把到达时间不是 0 的进程直接套错公式。}


\concept{等待时间与响应时间}
\begin{points}
  \item 等待时间是进程在就绪队列中等待 CPU 的总时间。
  \item 对非抢占式算法，等待时间常可由“周转时间减服务时间”得到。
  \item 响应时间是从提交请求到第一次得到响应的时间，交互系统特别关心。
  \item 时间片轮转等分时算法常优先优化响应时间，而不一定使平均周转时间最小。
\end{points}
\plain{批处理更关心“多久做完”，交互系统更关心“多久先有反应”。}
\exam{常考等待时间、周转时间、响应时间的区别，不要把“第一次响应”和“最终完成”混掉。}


\concept{调度题作答顺序}
\begin{points}
  \item 先按题目给出的到达时间、服务时间、优先级、时间片等规则画出调度顺序或甘特图。
  \item 再列每个进程的完成时间。
  \item 然后依次计算周转时间、等待时间、带权周转时间等指标。
  \item 只写最终数字而不写过程，考试里很容易丢步骤分。
\end{points}
\plain{老师录音特别强调这类题要把过程写出来，因为后面所有指标都依赖前面的调度顺序。}
\exam{计算题中通常既看结果也看过程，至少要把调度序列和每个指标的来源写清楚。}


\section{调度方式}
\concept{抢占式与非抢占式}
\begin{points}
  \item 非抢占式调度：进程一旦获得 CPU，除非主动放弃、阻塞或结束，否则不会被强行剥夺。
  \item 抢占式调度：操作系统可因时间片到、更高优先级进程到达等原因剥夺 CPU。
  \item 分时系统通常需要抢占式调度；批处理系统可以使用非抢占式调度。
\end{points}
\plain{抢占就是正在运行的进程也可能被更高优先级/时间片到期打断。}
\exam{常考抢占式与非抢占式适用场景。}


\concept{调度发生的时机}
\begin{points}
  \item 非抢占式调度通常在进程终止，或进程由运行态转入阻塞态时发生。
  \item 抢占式调度除上述时机外，还可在时间片到、运行态转就绪态、高优先级进程就绪、新进程进入就绪队列等情况下发生。
  \item 记忆关键：抢占式比非抢占式多了“即使当前进程还没结束，也可能因为外界条件变化被换下”。
\end{points}
\plain{调度时机题本质是在问：当前正在跑的进程，什么时候必须交棒，什么时候可能被强行收回 CPU。}
\exam{常考非抢占和抢占条件对比，尤其时间片到和高优先级进程到达是否触发调度。}


\section{常用调度算法}
\concept{先来先服务 FCFS}
\begin{points}
  \item 按进程到达就绪队列的先后顺序分配 CPU。
  \item 优点：简单、公平、不会饥饿。
  \item 缺点：短作业可能排在长作业后面，平均等待时间可能很长，存在“护航效应”。
  \item 适合批处理场景，不适合交互式短任务较多的场景。
\end{points}
\plain{谁先来谁先上，最朴素但容易让短作业排在长作业后面。}
\exam{常考护航效应和平均等待时间计算。}


\concept{短作业优先 SJF 与最短剩余时间优先 SRTF}
\begin{points}
  \item SJF 选择预计运行时间最短的作业先执行，是非抢占式算法。
  \item SRTF 是抢占式版本，总是选择剩余运行时间最短的进程。
  \item 优点：理论上可使平均等待时间最短。
  \item 缺点：需要预测运行时间；长作业可能长期等待，产生饥饿。
\end{points}
\plain{优先做短的，平均等待时间好看，但长作业可能一直等。}
\exam{常考非抢占 SJF 和抢占式 SRTF 的甘特图区别。}


\concept{SJF 为什么平均等待时间最短}
\begin{points}
  \item 在已知运行时间且都同时到达的理想条件下，SJF 可使平均等待时间达到最小。
  \item 直观原因是先完成短任务，能减少后续所有任务累计等待。
  \item 但现实系统通常无法准确预知服务时间，因此常需基于历史行为估计。
\end{points}
\plain{短任务先做，后面排队的人整体少等一些，所以平均值最好看。}
\exam{常考“哪种算法平均等待时间理论最优”，答案通常是 SJF/SRTF。}


\concept{优先级调度}
\begin{points}
  \item 为每个进程赋予优先级，优先级高者先运行。
  \item 可分为抢占式和非抢占式。
  \item 静态优先级创建后不变，动态优先级可根据等待时间、资源使用情况改变。
  \item 缺点是低优先级进程可能饥饿；常用老化 aging 提高等待过久进程的优先级。
\end{points}
\plain{给进程分等级，重要的先跑；问题是低优先级可能饥饿。}
\exam{常考静态/动态优先级、抢占/非抢占、老化技术防饥饿。}


\concept{时间片轮转 RR}
\begin{points}
  \item 每个就绪进程轮流获得一个时间片，时间片用完仍未完成则回到就绪队列尾部。
  \item 适合分时系统，强调响应时间和公平性。
  \item 时间片太大，退化为 FCFS；时间片太小，上下文切换开销过大。
  \item 做题时严格画时间轴，注意新到达进程插入队列的位置和时间片剩余规则。
\end{points}
\plain{每人轮流跑一小段，适合分时系统；时间片太大像 FCFS，太小切换开销大。}
\exam{常考时间片轮转甘特图，以及时间片大小对响应时间和开销的影响。}


\concept{高响应比优先 HRRN}
\begin{points}
  \item 设等待时间为 $W$、服务时间为 $S$，则响应比公式为 $R=\dfrac{W+S}{S}=1+\dfrac{W}{S}$。
  \item 等待时间越长，响应比越大；服务时间越短，响应比也越大。
  \item 该算法兼顾短作业优先和防止长作业饥饿。
  \item 通常是非抢占式调度，用在作业调度中较常见。
\end{points}
\plain{它像给短作业加分，同时又让等太久的长作业慢慢把分数追上来。}
\exam{常考响应比公式计算，以及 HRRN 为什么能缓解 SJF 的饥饿问题。}


\concept{多级队列调度}
\begin{points}
  \item 将就绪队列划分为多个队列，不同队列服务不同类型进程，如系统进程、交互进程、批处理进程。
  \item 队列之间有固定优先级或时间比例分配，队列内部可使用不同调度算法。
  \item 固定多级队列的问题是低优先级队列可能长期得不到服务。
\end{points}
\plain{先把不同类型任务分到不同窗口，再决定窗口之间谁优先。}
\exam{常考多级队列和多级反馈队列的区别：前者队列固定，后者进程可在队列间移动。}


\concept{多级反馈队列 MLFQ}
\begin{points}
  \item 多级反馈队列允许进程在不同优先级队列之间移动。
  \item 新进程通常先进入高优先级队列，时间片较短，若用完时间片还未完成则降级。
  \item 等待过久的低优先级进程可被提升，避免饥饿。
  \item 直觉：先假设新来的任务可能是短交互任务，给它快速响应；若它表现得像长计算任务，就逐步降到低优先级。
\end{points}
\plain{这是老师点名的重点算法：先快照顾短交互任务，再把长任务慢慢往低队列放。}
\exam{常考调度大题：写清抢占规则、降级规则、时间片大小、低优先级进程是否会老化提升。}


\section{实时调度、多处理器调度与线程调度}
\concept{实时调度基本概念}
\begin{points}
  \item 实时系统中，正确结果如果迟到，也可能等同于错误结果。
  \item 周期性任务常用处理时间 $t_i$、截止期限 $d_i$、周期 $p_i$ 描述，通常 $0\le t_i\le d_i\le p_i$。
  \item 必要的可调度条件之一：$\sum_{i=1}^{m}\dfrac{t_i}{p_i}\le 1$，表示 CPU 总利用需求不能超过 100\%。
\end{points}
\plain{这里的利用率条件是在做“算术上装不装得下”的第一层判断：若总需求都超过 100\%，后面再好的算法也救不了。}
\exam{常考给出若干周期任务的 $t_i,p_i$，先判断总利用率是否至少满足可调度必要条件。}


\concept{硬实时、软实时与两类延迟}
\begin{points}
  \item 硬实时系统存在必须满足的时间限制，错过截止期限通常不可接受。
  \item 软实时系统偶尔超时可以容忍，但仍希望尽量及时。
  \item \term{中断延迟}是从 CPU 收到中断到中断处理程序开始的时间。
  \item \term{分派延迟}是调度程序从停止一个进程到切换到另一个进程的时间。
  \item 实时系统要求这两类延迟都尽量短，且上界可预测。
\end{points}
\plain{实时系统不只看算法名字，往往先考“什么叫中断延迟、什么叫分派延迟，它们为什么要可预测”。}
\exam{常考硬实时/软实时区别，以及中断延迟和分派延迟分别指什么。}


\concept{速率单调 RMS 与最早截止期限 EDF}
\begin{points}
  \item RMS 是静态优先级算法，周期越短，优先级越高。
  \item EDF 是动态优先级算法，截止期限越近，优先级越高。
  \item EDF 在单处理器、理想条件下可达到更高利用率，但实现和开销更复杂。
\end{points}
\plain{RMS 和 EDF 的本质区别是“优先级按周期固定下来”还是“优先级随当前截止期限动态变化”。}
\exam{常考比较题：谁是静态优先级、谁是动态优先级、谁通常能支持更高利用率。}


\concept{RMS 可调度上界}
\begin{points}
  \item RMS 采用抢占式静态优先级策略，更短周期的任务优先级更高。
  \item 虽然总利用率 $\sum t_i/p_i \le 1$ 是必要条件，但对 RMS 来说还不总是充分。
  \item 调度 $N$ 个周期任务时，最坏情况下 CPU 利用率上界为 $N(2^{1/N}-1)$。
  \item 因此有些“总利用率不到 100\%”的任务集，按 RMS 仍可能不可调度。
\end{points}
\plain{RMS 不是只看总利用率小于 1 就一定行，它还有更严格的保守上界。}
\exam{常考为什么 RMS 下利用率小于 1 也可能错过截止期限，以及两任务时上界约为 83\%。}


\concept{比例分享调度与 POSIX 实时调度}
\begin{points}
  \item \term{比例分享调度}强调按预先约定的比例在多个任务、用户或任务组之间分配 CPU 份额，而不只是按先来先服务或固定优先级抢占。
  \item 公平分享调度 fair-share scheduling 可看作比例分享思想的一种实现，更关注“每个用户/用户组分到多少 CPU”，而不只是单个进程。
  \item \term{POSIX 实时调度}常见策略包括 \code{SCHED\_FIFO} 和 \code{SCHED\_RR}：前者按固定优先级先来先服务，后者在同优先级任务间再配合时间片轮转。
  \item 这类实时策略通常服务于实时优先级任务，强调高优先级先运行，并允许更高优先级任务抢占当前任务。
\end{points}
\plain{比例分享关心“CPU 份额怎么分”，POSIX 实时调度关心“同一实时优先级下是一直跑到底，还是轮流给时间片”。}
\exam{PPT 明确点了 \code{SCHED\_FIFO} 和 \code{SCHED\_RR}，选择题里常考两者区别：一个同优先级不轮转，一个同优先级按时间片轮转。}


\concept{多处理器调度}
\begin{points}
  \item 非对称多处理调度中，一个主处理器负责调度，其他处理器执行任务。
  \item 对称多处理调度中，各处理器都可参与调度，常见挑战是负载均衡和缓存亲和性。
  \item 处理器亲和性指尽量让进程继续在原 CPU 上运行，以利用缓存中的数据。
\end{points}
\plain{多核下不只要挑进程，还要挑“放到哪颗 CPU 上跑”才划算。}
\exam{常考非对称/对称多处理、亲和性与负载均衡之间的矛盾。}


\concept{亲和性与负载均衡}
\begin{points}
  \item \term{处理器亲和性}强调尽量让进程留在原 CPU 上运行，以减少缓存失效、TLB 失效等迁移代价。
  \item 亲和性可分为 soft affinity 和 hard affinity。
  \item \term{负载均衡}强调让各 CPU 工作负载尽量均匀，常见实现有 push migration 和 pull migration。
  \item 亲和性与负载均衡常有矛盾：留在原 CPU 更省缓存代价，但迁移到空闲 CPU 更能避免忙闲不均。
\end{points}
\plain{这类题别只背名词，核心矛盾就是“少迁移更省缓存”与“适当迁移更均衡负载”之间的取舍。}
\exam{常考 soft/hard affinity、push/pull migration 的含义，以及亲和性和负载均衡为什么可能冲突。}


\concept{多处理器调度算法}
\begin{points}
  \item \term{负载共享}：系统维护全局就绪队列，空闲处理器从中取线程运行，优点是处理器不易闲置，缺点是公共队列可能成为瓶颈。
  \item \term{群调度}：把一组相关线程同时调度到一组处理器上运行，适合同步紧密的并行任务。
  \item \term{专用处理器分配}：给一个应用固定分配一组处理器直到其运行结束，追求高并行度但可能牺牲单个处理器利用率。
  \item \term{动态调度}：OS 负责在应用间分配处理器，应用自己决定其获得的处理器上哪些线程运行，灵活但实现复杂。
\end{points}
\plain{多处理器算法题重点看“用一个全局池分配”，还是“按线程组一起分配”，还是“直接包核给应用”。}
\exam{常考几种多处理器调度算法的基本思想、优缺点和适用场景。}


\concept{线程调度}
\begin{points}
  \item 用户级线程由线程库调度，内核只看到进程。
  \item 内核级线程由内核调度，可以在多个处理器上并行。
  \item 线程调度更强调轻量切换和同一进程内线程之间的同步关系。
\end{points}
\plain{线程调度这节重点不在重新定义线程，而在分清“用户级线程谁在调度、内核级线程谁在调度”，以及这会怎样影响阻塞和并行能力。}
\exam{常考进程/线程区别、用户级线程和内核级线程优缺点、多对一/一对一/多对多模型。}


\concept{典型系统调度实例}
\begin{points}
  \item \term{UNIX System V}：采用动态优先级多级队列调度，CPU 用得越久优先级越低，阻塞等待后返回的交互/I/O 型进程优先级相对更高。
  \item \term{Solaris}：基于优先级的线程调度，不同调度类（如实时、分时、公平分享）可采用不同策略。
  \item \term{Linux}：现代普通任务默认使用 CFS，实时任务常采用符合 POSIX 的 \code{SCHED\_FIFO}/\code{SCHED\_RR}。
  \item \term{Windows}：以线程为调度对象，采用动态优先级、抢占式、多级反馈队列思想，并对交互/I/O 型线程做短暂优先级提升。
\end{points}
\plain{这部分不要求把各系统实现细节全背下来，更重要的是抓住“UNIX 动态优先级、Linux CFS、Windows 线程调度、Solaris 多调度类”这几个标签。}
\exam{若简答题问“现代系统调度器有什么共同趋势”，常可答动态优先级、偏向交互任务、支持实时类和多处理器环境。}


\concept{调度算法的评估方法}
\begin{points}
  \item \term{确定性模型}：给定一组确定的进程到达时间和服务时间，直接计算不同算法下的等待时间、周转时间等指标，适合课堂示例和手工比较。
  \item \term{排队模型}：把进程到达和服务时间看成概率分布，用排队论分析 CPU 利用率、平均等待时间、平均队列长度等期望指标。
  \item \term{仿真法}：用程序模拟调度过程，可用随机数据或真实系统跟踪数据作为输入，结果更接近实际，但实现和统计成本更高。
  \item 三者大致对应“最简单但最理想化”“数学分析更强但依赖分布假设”“最贴近现实但成本最高”的递进关系。
\end{points}
\plain{这一节不是新调度算法，而是在讲“怎么比较算法好坏”：小题目手算常用确定性模型，系统级研究更常见排队分析和仿真。}
\exam{PPT 在章末单独列了评估方法，选择题常考三者各自优缺点：确定性模型快但依赖输入，排队模型依赖分布假设，仿真法更真实但成本高。}

\chapter{进程同步}

\section{同步与互斥的概念}
\concept{并发进程的制约关系}
\begin{points}
  \item 多道程序系统中，由于资源共享或进程合作，并发进程之间形成相互制约关系。
  \item 间接相互制约：多个进程竞争同一资源，例如多个进程争用打印机。
  \item 直接相互制约：进程之间存在合作顺序，例如生产者生产后消费者才能消费。
  \item 进程同步的任务是协调这些关系，使并发进程能正确共享资源和相互合作，保证结果可再现。
\end{points}
\plain{同步这一章关心的不是“有没有并发”，而是“并发以后怎么别把共享数据搞坏”。}
\exam{常考直接/间接制约、为什么并发会导致执行结果不可再现。}

\concept{同步与互斥}
\begin{points}
  \item \term{同步}针对合作关系：多个进程在某些关键点上需要互相等待或按次序推进。
  \item \term{互斥}针对竞争关系：某个进程使用临界资源时，其他进程必须等待它释放资源。
  \item 同步强调“先后次序正确”；互斥强调“同一时刻只能一个进入”。
  \item 两者常同时出现：例如生产者-消费者问题中，访问缓冲区要互斥，放入/取出产品又要满足同步关系。
\end{points}
\plain{区分这两个词最稳的方法是看问题本质：如果是在管“谁先谁后”，偏同步；如果是在管“能不能同时进”，偏互斥。}
\exam{简答题经常直接问同步和互斥的区别，也可能给场景让你判断哪些关系属于同步、哪些属于互斥。}


\concept{与时间有关的错误/race condition}
\begin{points}
  \item 多个进程并发读写共享变量，执行次序不同可能导致不同结果，这种错误称为与时间有关的错误或竞争条件。
  \item 例子：订票系统中 A、B 同时读取剩余票数 $x=1$，都判断有票并出票，最终可能超卖。
  \item 根本原因：对共享变量的复合操作没有互斥保护。
  \item 解决思路：把读-改-写共享变量的代码段放入临界区，一次只允许一个进程进入。
\end{points}
\plain{竞争条件的关键不是“同时访问”这四个字，而是多个进程把一次完整更新拆成了几步，步骤交叉后就会把共享结果写乱。}
\exam{常考给一段并发读改写代码，让你解释为什么可能超卖、少减或结果不唯一。}


\concept{与时间有关的另一类错误：永远等待}
\begin{points}
  \item 与时间有关的错误不只会造成“结果不唯一”，还可能造成“该等的人没被正确唤醒”，从而永远等待。
  \item 典型情形是：进程刚判断完条件、还未来得及把自己放入等待队列，另一个进程就先释放资源并完成唤醒操作。
  \item 结果是释放动作已经发生，但等待进程还没真正登记到等待队列中，后续再也没人唤醒它。
  \item 这说明同步操作必须把“检查条件、入队等待、释放资源、唤醒他人”这些关键动作作为受控整体来设计。
\end{points}
\plain{有的并发错误不是把结果算错，而是把唤醒时机错过去，最后进程一直睡着没人再叫。}
\exam{常考“为什么简单 if/while + 等待队列操作也会出错”，答案要点是检查条件与入队/唤醒不是原子完成。}


\concept{临界资源与临界区}
\begin{points}
  \item 临界资源是一段时间内只允许一个进程使用的资源。
  \item 临界区是进程中访问临界资源的代码段。
  \item 进入区负责申请进入临界区；退出区负责释放临界区；剩余区是其他代码。
  \item 互斥要求：任意时刻最多一个进程在相关临界区内执行。
\end{points}
\plain{临界资源是“对象”，临界区是“访问这个对象的那段代码”；题目里真正要保护的是代码段，不是抽象名词本身。}
\exam{选择题很爱把临界资源和临界区混着考，答题时最好直接举“打印机/共享变量”和“访问它的代码段”来区分。}


\concept{临界区访问原则}
\begin{points}
  \item \key{空闲让进}：临界区空闲时，应允许一个请求进入的进程立即进入。
  \item \key{忙则等待}：临界区已有进程时，其他请求进程必须等待。
  \item \key{有限等待}：等待进入临界区的进程不能无限期等待，避免饥饿。
  \item \key{让权等待}：不能进入临界区的进程应释放处理机，不应长时间忙等浪费 CPU。
\end{points}
\plain{四条规则里最容易被忽略的是“让权等待”：进不去就该睡，不应一直空转耗 CPU，这也是后面记录型信号量优于忙等法的原因。}
\exam{常考把某个互斥方案和四条准则对照，判断它是否满足有限等待或让权等待。}


\concept{课本中的三原则表述}
\begin{points}
  \item 教材中常把临界区要求概括为：\key{互斥原则}、\key{推进原则}、\key{有限等待原则}。
  \item 互斥原则对应“同一时刻最多一个进程在临界区”。
  \item 推进原则强调当临界区空闲时，不能无限拖延某个请求进入的进程。
  \item 有限等待原则强调不能让某个进程无限期等下去。
\end{points}
\plain{不同资料有“四条规则”和“三条原则”两种说法，考试里看到哪种表述都要能对上同一组含义。}
\exam{常考把“空闲让进、忙则等待、有限等待、让权等待”和“互斥/推进/有限等待”对应起来。}

\section{互斥的实现方法}
\concept{软件方法}
\begin{points}
  \item 单标志法、双标志先检查法、双标志后检查法、Peterson 算法都是早期软件互斥方法。
  \item 软件方法的关键是用共享变量表达“谁想进”和“轮到谁进”。
  \item Peterson 算法能满足互斥、空闲让进和有限等待，但主要适合两个进程的理论模型。
  \item 软件方法缺点是复杂、容易出错、依赖忙等，难以推广。
\end{points}
\plain{软件法这节重点不在会背所有算法，而在看出它们共同缺陷：逻辑绕、推广差，而且很多方案本质上还在忙等。}
\exam{常考 Peterson 为什么理论上正确、却不适合现代通用系统实现；也常问软件法共同短板。}

\concept{Peterson 算法}
\begin{points}
  \item Peterson 算法面向两个进程，结合了\term{意愿标志} \code{flag[i]} 与\term{轮转变量} \code{turn}。
  \item 进程进入前先声明“我想进”（置 \code{flag[i]=true}），再把优先权让给对方（设置 \code{turn}）。
  \item 若对方也想进入，且当前 \code{turn} 指向对方，则本进程在 while 条件中等待；否则进入临界区。
  \item 退出临界区时把自己的 \code{flag[i]} 复位为 \code{false}。
\end{points}
\plain{Peterson 的直觉是“先表明意愿，再主动谦让一次”；如果对方此时也想进，就按 \code{turn} 决定谁先走，因此既避免同时进入，也避免两边永久互让。}
\exam{代码题常考 \code{flag} 和 \code{turn} 各表示什么，或问它为什么能满足互斥与有限等待。}

\concept{Dekker 算法}
\begin{points}
  \item Dekker 算法也是面向两个进程的软件互斥算法，结合了意愿标志 \code{flag[]} 和轮转变量 \code{turn}。
  \item 它的核心思想是：若双方都想进入临界区，则按 \code{turn} 决定谁暂时后退，让另一方先进入。
  \item 因而 Dekker 算法修正了“双方都置位后互相谦让到谁也进不去”的问题，是一个正确的软件互斥方案。
  \item 但其逻辑较复杂，通常更多作为理论过渡，后续常用更简洁的 Peterson 算法来讲解。
\end{points}
\plain{Dekker 可以理解成 Peterson 之前“更绕但能用”的版本：它已经把“想进”和“轮到谁”两类信息结合起来了，只是代码结构更复杂。}
\exam{若题目问 Dekker 和 Peterson 的共同思想，抓住“意愿标志 + 轮转变量”这条主线即可。}


\concept{硬件方法}
\begin{points}
  \item 关中断：单处理器中进入临界区前关中断，退出后开中断，可防止进程被中断切换。
  \item Test-and-Set 指令：原子地测试并设置锁变量。
  \item Swap 指令：原子地交换两个变量的值。
  \item 硬件方法可保证互斥，但若采用自旋锁，会造成忙等，不满足让权等待。
\end{points}
\plain{硬件法的价值是“提供原子性”，短板是很多实现仍让进程原地自旋，所以它解决了互斥，不一定解决了让权等待。}
\exam{常考为什么关中断只适合单处理器、为什么 Test-and-Set 能保互斥却仍可能浪费 CPU。}

\concept{关中断方法}
\begin{points}
  \item 在单处理器系统中，进程切换通常依赖时钟中断或 I/O 中断；进入临界区前先关中断，可避免当前进程被切走。
  \item 因而可用“关中断 $\rightarrow$ 执行临界区 $\rightarrow$ 开中断”的方式实现互斥。
  \item 它的主要缺点是会推迟中断响应、降低系统并发性，而且对多处理器无效。
  \item 由于若用户随意关中断可能导致系统失控，所以这种方法通常只适合内核中的短临界区。
\end{points}
\plain{关中断法的本质不是“上锁”，而是直接暂时禁止调度发生；所以它虽然简单，但代价是把系统响应能力一起压住了。}
\exam{比较题常问为什么关中断不适合多处理器，以及为什么不能把关中断权限交给普通用户进程。}

\concept{Test-and-Set 与 Swap 原子指令}
\begin{points}
  \item \term{Test-and-Set} 指令能原子地读出锁旧值并把锁置为占用，常写成 \code{while TS(\&lock);} 后进入临界区。
  \item \term{Swap} 指令能原子地交换共享锁变量与本地变量的值，也可据此反复尝试获取锁。
  \item 这类原子指令的共同点是：把“检查锁是否空闲”和“把锁改成占用”合并成不可分割的硬件操作。
  \item 它们适用于多处理器环境，但若失败后一直循环重试，就会形成忙等，不满足让权等待，且可能导致饥饿。
\end{points}
\plain{原子指令法真正解决的是“检查”和“上锁”之间不能被别人插进来；它修复了软件法最难保证的原子性问题，但没自动解决忙等。}
\exam{常考 TS/Swap 为什么能保证互斥，以及它们为什么仍可能浪费 CPU、不能保证有限等待。}


\concept{互斥锁机制}
\begin{points}
  \item 互斥锁 mutex lock 是操作系统提供的互斥访问工具，通常借助前面的原子硬件指令实现。
  \item 锁变量常用两种状态表示资源：0 表示资源可用（开锁），1 表示资源已被占用（关锁）。
  \item 进入临界区前执行 \code{lock}：若锁空闲则把锁置为 1 并进入；若锁已被占用，则继续等待。
  \item 退出临界区后执行 \code{unlock}，把锁恢复为 0，使其他等待者有机会进入。
  \item 若等待时一直反复检测锁状态，这种实现就是自旋锁；它有忙等开销，但在临界区很短、多处理器环境下常较实用。
\end{points}
\plain{互斥锁可以看成把“测试并上锁”“离开时解锁”包装成统一接口；真正的考点通常是它为什么简单好用、又为什么可能忙等浪费 CPU。}
\exam{PPT 把 mutex 和 spinlock 连着讲，选择题常考 0/1 两种状态、\code{lock}/\code{unlock} 的作用，以及自旋锁更适合“锁持有时间短”的场景。}


\section{信号量与 P/V 操作}
\concept{信号量定义}
\begin{points}
  \item 信号量 semaphore 由整数值 $s$ 和等待队列 $q$ 组成。
  \item 除初始化外，信号量的值只能通过 P 操作和 V 操作改变。
  \item 信号量 $s>0$ 时，表示可用资源数；$s=0$ 表示无可用资源但无人等待；$s<0$ 时，$|s|$ 表示等待队列中的进程数。
  \item P/V 操作必须是原子操作，执行过程中不能被其他进程打断。
\end{points}
\plain{定义题真正要抓的是“整数值 + 等待队列 + 只能被 P/V 改变 + 原子操作”，别只把它背成一个计数器。}
\exam{常考记录型信号量的组成，以及为什么 P/V 必须原子执行。}


\concept{P 操作}
\begin{points}
  \item P(S)：先执行 $s=s-1$。
  \item 若减完后 $s<0$，说明资源不够，当前进程阻塞，进入该信号量等待队列，并转调度程序。
  \item 若减完后 $s\ge 0$，说明成功获得资源，进程继续执行。
  \item 直觉：P 是“申请/占用一个资源”，先登记需求，再判断是否需要等待。
\end{points}
\plain{P 操作最容易写错的地方是“先减再判断”，不是先看够不够再减；这个顺序正是统一资源计数和等待队列语义的关键。}
\exam{常考手写 P 原语伪代码，或问为什么减完后小于 0 就表示当前进程必须阻塞。}


\concept{V 操作}
\begin{points}
  \item V(S)：先执行 $s=s+1$。
  \item 若加完后 $s\le 0$，说明仍有进程在等待队列中，唤醒一个等待进程并插入就绪队列。
  \item 若加完后 $s>0$，说明没有进程等待，只是增加了一个可用资源。
  \item 直觉：V 是“释放一个资源”。即使加完后 $s\le 0$，也要唤醒，因为负数表示有等待者；释放的资源应交给其中一个等待进程。
\end{points}
\plain{V 操作不是简单“把数加回去”，而是把释放出来的资源真正交还给等待队列里的某个进程，所以加完后若仍不大于 0 就得唤醒别人。}
\exam{常考为什么 V 后条件写成 `s <= 0` 而不是 `s < 0`，答题要结合“负数表示仍有人排队”等待来解释。}


\concept{互斥信号量取值范围}
\begin{points}
  \item 若 $N$ 个进程共享一个临界资源，互斥信号量初值通常设为 1。
  \item 当没有进程占用该资源时，信号量值可为 1。
  \item 当已有 1 个进程在临界区、其余部分进程等待时，信号量值会落到 0、-1、-2、...。
  \item 因此对一个临界资源的互斥信号量，其取值范围可理解为 $1,0,-1,\dots,-(N-1)$。
\end{points}
\plain{互斥信号量只有一份资源，所以初值是 1；后面一旦出现负数，本质上就是在数有多少人排队。}
\exam{常考“一个临界资源被 N 个进程共享时，互斥信号量初值和取值范围是什么”。}

\concept{整型信号量与记录型信号量}
\begin{points}
  \item 整型信号量只用整数表示资源数，P 操作中可能忙等。
  \item 记录型信号量增加等待队列，资源不足时让进程阻塞，满足让权等待。
  \item 考试中常默认 P/V 是记录型信号量，阻塞和唤醒由 OS 完成。
\end{points}
\plain{整型信号量和记录型信号量的根本差别不在“一个是 int 一个是 struct”，而在资源不足时是忙等还是阻塞让权。}
\exam{常考为什么记录型信号量比整型信号量更符合临界区的让权等待原则。}


\section{经典同步问题}
\concept{生产者-消费者问题}
\begin{points}
  \item 问题：多个生产者向有界缓冲区放产品，多个消费者从缓冲区取产品。
  \item 需要三个信号量：\code{mutex=1} 保护缓冲区互斥访问，\code{empty=n} 表示空缓冲区数，\code{full=0} 表示满缓冲区数。
  \item 生产者顺序：生产产品 $\rightarrow$ P(empty) $\rightarrow$ P(mutex) $\rightarrow$ 放入缓冲区 $\rightarrow$ V(mutex) $\rightarrow$ V(full)。
  \item 消费者顺序：P(full) $\rightarrow$ P(mutex) $\rightarrow$ 取出产品 $\rightarrow$ V(mutex) $\rightarrow$ V(empty) $\rightarrow$ 消费产品。
  \item 注意：先申请同步资源 empty/full，再申请互斥 mutex，避免拿着锁等待条件导致其他进程无法改变条件。
\end{points}
\plain{生产者消费者题最关键的顺序是先 P(empty/full) 再 P(mutex)，否则容易出现“拿着锁等条件”，把能改变条件的人也堵死。}
\exam{常考把 mutex 和 empty/full 的次序打乱后会不会死锁，或让你改正错误伪代码。}


\concept{P/V 题的规范写法}
\begin{points}
  \item 先定义每个信号量、初值以及它表示的含义。
  \item 伪代码中把每类进程的关键步骤写完整，不要只写零散几句 P/V。
  \item 写完后检查是否“有 P 就有对应的 V”，以及是否有人拿着互斥锁等待条件。
\end{points}
\plain{老师录音里反复强调：不会全写对也要规范写，定义清楚、关系写清楚，就能拿过程分。}
\exam{大题常按步骤给分，哪怕只分析出部分互斥/同步关系，也比空着强。}


\concept{读者-写者问题}
\begin{points}
  \item 多个读者可同时读共享文件，但写者必须互斥写；写者与读者也互斥。
  \item 读者优先方案中，第一个读者负责阻止写者，最后一个读者负责释放写者。
  \item 需要计数变量 \code{readcount} 记录当前读者数，并用互斥信号量保护它。
  \item 易错点：读者之间不是互斥读文件，而是互斥修改 \code{readcount}。
\end{points}
\plain{读者写者题最容易误判的是“读者共享读文件，但不共享改 `readcount` 这一步”；真正要互斥保护的是读者计数变量。}
\exam{常考第一个/最后一个读者分别负责什么，或者为什么读者之间不是完全没有互斥。}


\concept{前趋关系}
\begin{points}
  \item 若语句 $S_1$ 必须在 $S_2$ 之前完成，可设置初值为 0 的同步信号量 $S$。
  \item 前驱进程在完成 $S_1$ 后执行 V(S)。
  \item 后继进程在执行 $S_2$ 前先执行 P(S)。
  \item 前驱图较复杂时，可按每条前趋边设置同步信号量。
\end{points}
\plain{谁必须先做完，就让前面的人 V 一下，后面的人先 P 住等信号。}
\exam{常考把前驱图翻译成 P/V 代码。}


\concept{复杂前趋图的信号量设置方法}
\begin{points}
  \item 对复杂前趋图，常可按“每一条前趋边对应一个初值为 0 的同步信号量”来建模。
  \item 某进程若有多个直接前驱，则通常需要在开始前对多个同步信号量分别执行 P 操作。
  \item 某进程若有多个直接后继，则通常在结束后对多个同步信号量分别执行 V 操作。
  \item 这种方法虽然信号量较多，但结构清楚，考试中最稳妥。
\end{points}
\plain{多前驱就多次 P，多后继就多次 V，把图上的箭头一条条翻译掉，最不容易漏关系。}
\exam{常考六进程前趋图这类题，评分点通常就在“信号量初值设 0、前驱后继逐条翻译”上。}


\concept{哲学家进餐问题}
\begin{points}
  \item 五个哲学家围桌，每人左右各一只筷子，必须同时拿到两只筷子才能进餐。
  \item 若所有哲学家都先拿左筷子再等右筷子，可能形成死锁。
  \item 解决方法包括：最多允许四人同时取筷子；奇偶哲学家取筷子顺序相反；一次性申请两只筷子；使用管程。
  \item 核心直觉：避免每个进程“占有一个资源，同时等待另一个资源”。
\end{points}
\plain{哲学家进餐题不是单纯考代码套路，而是在追问“怎样打破占有一部分资源再等待另一部分资源”的死锁结构。}
\exam{常考判断某种取筷子策略是否会死锁，并说明它破坏了哪个死锁必要条件。}


\section{信号量集机制}
\concept{AND 型信号量}
\begin{points}
  \item AND 型信号量适用于进程同时需要多类资源、每类一个的情况。
  \item 基本思想：进程需要的多类资源一次性全部分配；只要有一个资源不满足，其他资源也不分配。
  \item 这样可以避免进程占有部分资源再等待其他资源，从而降低死锁风险。
  \item P 原语称为 SP/Swait，V 原语称为 SV/Ssignal。
\end{points}
\plain{它相当于“要么一起批，要么一个也别先拿”，避免拿到一半卡死。}
\exam{常考它为什么能缓解哲学家进餐这类“先占一部分再等另一部分”的死锁风险。}


\concept{一般信号量集}
\begin{points}
  \item 一般信号量集是 AND 型的扩展：一次可申请多类资源，每类资源可申请多个。
  \item 当某类资源低于下限或不能满足申请量时，本次原语不分配任何资源。
  \item 它把“多个 P 操作”合并成一个原子判断，避免中间状态。
\end{points}
\plain{关键不是新符号多，而是“多个资源请求合成一次原子判断”。}
\exam{常考它和基本记录型信号量、AND 型信号量的关系。}


\concept{信号量集的特殊退化形式}
\begin{points}
  \item 当信号量集写成 \code{SP(S,d,d)} 时，可理解为一次申请同类资源的 $d$ 个实例。
  \item 当写成 \code{SP(S,1,1)} 时，基本退化为普通记录型信号量。
  \item 当写成 \code{SP(S,1,0)} 时，可把它看成一种“开关型”控制：$S\ge 1$ 时允许进入，$S=0$ 时禁止进入。
\end{points}
\plain{这些特殊形式的价值在于帮你看出：信号量集并不是全新体系，而是把普通 P/V 和批量申请统一到一套写法里。}
\exam{常考信号量集和普通信号量的退化关系，尤其 \code{SP(S,1,1)} 为什么等价于基本信号量。}


\section{管程}
\concept{管程是什么}
\begin{points}
  \item 管程是一种高级同步机制，把共享数据、对共享数据的操作、同步条件封装在一个模块中。
  \item 管程保证任意时刻只有一个进程/线程在管程内部执行，因此天然提供互斥。
  \item 条件变量用于表达“某个条件不满足时等待，条件满足时被唤醒”。
  \item 典型操作包括 \code{wait} 和 \code{signal}。
\end{points}
\plain{这里先抓管程的“整体观”：共享数据、对数据的操作、同步条件被打包在一起，所以互斥不再靠程序员手动到处插 P/V。}
\exam{常考管程定义题，答题时最好同时写出“封装共享数据”和“任一时刻只允许一个进程进入”两层。}


\concept{管程的构成}
\begin{points}
  \item 管程通常由四部分组成：管程名、局部共享数据结构、对共享数据进行操作的一组过程、初始化共享数据的语句。
  \item 共享数据和相关操作被封装在同一抽象模块内，外部进程不能随意直接改动内部数据。
  \item 管程中的过程常可分为外部可调用过程和仅供管程内部使用的内部过程。
\end{points}
\plain{这节要抓“数据和操作被打包到一起”，它不像信号量那样只是给你一个原语，而更像带规则的抽象数据类型。}
\exam{PPT 单独列了管程构成，简答题里最好把“四部分”按顺序写出来，而不是只笼统说“有共享数据和过程”。}


\concept{管程的基本特性}
\begin{points}
  \item \term{安全性}：管程内的局部数据只能被管程内的过程访问。
  \item \term{共享性}：多个进程可通过调用同一个管程过程来共享受保护资源。
  \item \term{互斥性}：任意时刻至多一个进程真正进入管程内部执行。
  \item \term{透明性}：互斥进入等机制由语言/系统自动保证，程序员不必像写 P/V 那样手工到处安插同步操作。
\end{points}
\plain{这四条里最有管程味道的是“透明性”：程序员看到的是过程调用，互斥细节被语言和运行时隐藏起来。}
\exam{常考“为什么说管程比信号量更不容易写错”，答案通常要落到互斥自动化和透明性。}


\concept{入口等待队列与条件队列}
\begin{points}
  \item 由于管程一次只允许一个进程进入，所以当管程已被占用时，后来想进入的进程要先在\term{入口等待队列}排队。
  \item 若进程已经进入管程，但因为某个条件不满足而不能继续执行，则它会在对应的\term{条件变量等待队列}上睡眠。
  \item 入口等待队列反映的是“还没进门就在等互斥入口”；条件队列反映的是“已经进门，但因资源条件不满足而等待”。
  \item \code{wait} 会释放管程互斥权并转入相应条件队列；之后别的进程才能从入口队列进入管程继续执行。
\end{points}
\plain{这两个队列最容易混：一个是在门外排队等“进去”，一个是在屋里发现条件不够、主动让出管程后按原因排队等“继续”。}
\exam{PPT 把入口等待队列和条件变量分开讲了，简答题里若问管程内部等待机制，最好把这两层队列都写出来。}


\concept{管程与 P/V 的区别}
\begin{points}
  \item P/V 是低级同步原语，程序员需要自己安排每个 P、V 的位置，容易写错。
  \item 管程把共享变量和互斥规则封装起来，进入管程自动互斥，程序员主要表达条件等待。
  \item P/V 更像“手动锁门开门”；管程更像“把房间设计成一次只能进一个人，并提供排队等待条件”。
  \item 管程更结构化，适合高级语言；信号量更基础，适合描述底层同步。
\end{points}
\plain{这一节真正要比较的是同步风格：P/V 靠程序员手工排布原语，管程则把互斥和等待条件封进模块接口里，结构更清晰。}
\exam{常考比较题：为什么说管程比信号量更高级、更不容易出错，但底层实现仍可依赖 P/V。}


\concept{条件变量的 \code{wait}/\code{signal} 与 P/V 的区别}
\begin{points}
  \item 条件变量用于区分不同等待原因，每个条件变量对应一个条件队列。
  \item \code{x.wait} 会使当前进程在条件变量 $x$ 上等待，并释放管程互斥权。
  \item \code{x.signal} 唤醒 $x$ 队列中的一个等待进程；若队列为空，则相当于空操作。
  \item 与信号量不同，条件变量不做资源计数，没有累计值；它只维护等待队列。
\end{points}
\plain{条件变量更像“按原因排队”，不是“拿一个会累计加减的资源计数器”。}
\exam{常考为什么条件变量的 \code{signal} 在没人等待时不会“记账”，以及它和 V 操作的本质区别。}


\chapter{死锁}

\section{死锁概念}
\concept{死锁定义}
\begin{points}
  \item 死锁是多个进程因竞争系统资源而造成的一种僵局，若无外力作用，这些进程都将永远不能向前推进。
  \item 若进程集合中的每个进程都在等待只能由该集合中其他进程引发的事件，则该集合发生死锁。
  \item 典型例子：P1 占有打印机等待读卡机，P2 占有读卡机等待打印机，双方都无法继续。
\end{points}
\plain{定义题要抓“相互等待 + 无外力就永远推进不了”这两层，重点不是普通阻塞，而是整个等待环自己解不开。}
\exam{名词解释或判断题常拿“暂时阻塞”和“死锁”混淆，答案里最好点出“若无外力作用将永远不能推进”。}


\concept{死锁与饥饿、死循环}
\begin{points}
  \item 死锁：一组进程相互等待，若无外力永远不能推进。
  \item 饥饿：某进程长期得不到所需资源，但系统整体仍在推进。
  \item 死循环：一个进程自身逻辑错误导致反复执行，不一定涉及资源等待。
  \item 活锁：进程没有阻塞，但不断改变状态、相互让步，导致没有实际进展。
\end{points}
\plain{区分题的关键是看“系统整体还动不动”：饥饿是别人还能跑，死循环是自己瞎转，死锁则是一组进程彼此卡死。}
\exam{很适合出选择/判断：给一个现象，让你判断是死锁、饥饿、活锁还是死循环。}


\section{资源分类与死锁必要条件}
\concept{资源分类}
\begin{points}
  \item 可剥夺资源：可从进程手中强行取回且不破坏计算，例如 CPU、内存页框。
  \item 不可剥夺资源：一旦分配，不能强行剥夺，否则会造成错误，例如打印机、部分外设。
  \item 可重用资源：使用完可再次分配，如设备、内存。
  \item 消耗性资源：使用后消失，如消息、信号。
  \item 竞争永久性资源和消耗性资源都可能产生死锁；竞争可剥夺资源通常不会形成典型死锁。
\end{points}
\plain{这里别把死锁只和“打印机这类永久资源”绑死，PPT 明确强调消息这类消耗性资源也能形成循环等待。}
\exam{常考可剥夺/不可剥夺、永久性/消耗性资源的分类，以及哪些竞争可能引发死锁。}


\concept{死锁产生的四个必要条件}
\begin{points}
  \item \key{互斥条件}：资源一次只能被一个进程使用。
  \item \key{请求和保持条件}：进程已保持至少一个资源，同时又请求新的资源。
  \item \key{不可剥夺条件}：资源不能被强行夺走，只能由占有进程主动释放。
  \item \key{循环等待条件}：存在进程资源等待环，$P_1$ 等 $P_2$，$P_2$ 等 $P_3$，$\ldots$，$P_n$ 等 $P_1$。
\end{points}
\plain{四个条件不是“死锁特征列表”，而是判断能否从制度上预防死锁的抓手: 只要破坏其中一个，死锁就不成立。}
\exam{简答题常要求逐条解释四个必要条件，或反问某种措施到底破坏了哪一个条件。}

\concept{死锁的根本原因与诱因}
\begin{points}
  \item 根本原因是竞争资源，而资源又不能同时满足所有进程的请求。
  \item 诱因是进程推进顺序不当，即资源申请和释放次序安排得不好。
  \item 只有“竞争资源”还不一定死锁；还要配合不合适的推进顺序才会卡死。
\end{points}
\plain{资源紧张只是埋雷，申请顺序乱了才真的把雷踩爆。}
\exam{常考简答：死锁为什么不是“有竞争就必死锁”，而是资源条件和推进顺序共同作用。}

\section{处理死锁的基本方法}
\concept{四类方法总览}
\begin{points}
  \item 死锁预防：通过制度性限制破坏死锁必要条件之一。
  \item 死锁避免：运行时根据资源需求判断分配是否安全，典型算法是银行家算法。
  \item 死锁检测与解除：允许死锁发生，定期检测，发生后恢复。
  \item 鸵鸟策略：忽略死锁问题，适用于死锁概率很低且处理成本过高的系统。
\end{points}
\plain{这题最容易混的是“预防、避免、检测”三个层次：预防是事前禁规则，避免是分配前看安全性，检测是先允许出事再查。}
\exam{常考四种处理思路的区别，尤其银行家算法属于避免，不属于检测。}


\concept{死锁预防}
\begin{points}
  \item 破坏互斥条件：把独占设备改造成共享设备，例如 SPOOLing 把打印机虚拟成共享设备。但并非所有资源都能共享。
  \item 破坏请求和保持条件：进程运行前一次性申请所有资源，或申请新资源前释放已占资源。缺点是资源利用率低、可能饥饿。
  \item 破坏不可剥夺条件：资源申请失败时剥夺其已占资源。适用于 CPU、内存等可剥夺资源，不适合打印机等。
  \item 破坏循环等待条件：对资源统一编号，进程按编号递增顺序申请资源。
\end{points}
\plain{预防题不是背四条措施，而是要能明确说出“这条办法破坏了哪个必要条件，代价是什么”。}
\exam{常考静态分配法/有序资源分配法为什么能预防死锁，答案里最好点名破坏的条件。}


\concept{有序资源分配法为什么能防止死锁}
\begin{points}
  \item 给所有资源类型统一编号，要求进程只能按编号递增顺序申请资源。
  \item 这样若进程已占有编号较小的资源，就不能再回头去申请编号更小或相同层次的资源。
  \item 因而不可能形成“资源编号沿环路严格递增、最后又回到起点”的循环等待。
  \item 所以该方法本质上是通过破坏循环等待条件来预防死锁。
\end{points}
\plain{考试答“为什么有效”时不要只写一句“按顺序申请”，要把“因此不可能形成环”这层逻辑补出来。}
\exam{常考证明式简答：为什么有序资源分配法可以防止死锁。}


\concept{安全状态与不安全状态}
\begin{points}
  \item 安全状态指系统能找到一个安全序列，使所有进程都能按某种顺序获得所需资源并完成。
  \item 不安全状态不一定已经死锁，但可能发展为死锁。
  \item 死锁状态一定是不安全状态；不安全状态不一定是死锁状态。
\end{points}
\plain{死锁避免的核心不是“现在卡没卡死”，而是“这一步分配后系统还留不留得下安全退路”。}
\exam{常考安全、不安全、死锁三者关系，尤其“不安全不等于已死锁”。}


\concept{死锁、 安全状态、不安全状态的关系}
\begin{points}
  \item 安全状态一定不会死锁，因为至少存在一个安全序列。
  \item 不安全状态不等于已经死锁，但系统有可能进一步走向死锁。
  \item 死锁状态一定是不安全状态。
\end{points}
\plain{安全像“还有退路”，不安全像“退路没法保证了”，死锁则是“已经彻底堵死了”。}
\exam{常考判断题：不安全状态不一定死锁，这句话必须会。}


\concept{银行家算法}
\begin{points}
  \item 银行家算法用于死锁避免，假设系统知道每个进程的最大资源需求。
  \item 关键数据结构：Available 可用资源，Max 最大需求，Allocation 已分配，Need 尚需，且 $Need=Max-Allocation$。
  \item 资源请求满足 $Request_i\le Need_i$ 且 $Request_i\le Available$ 后，先试探性分配。
  \item 再执行安全性检查：寻找一个进程满足 $Need_i\le Work$，假设其完成并释放资源，更新 $Work=Work+Allocation_i$，直到所有进程完成或找不到下一个。
  \item 若试分配后仍安全，则真正分配；否则让进程等待。
\end{points}
\plain{银行家算法题的核心不是算表本身，而是证明“试分配之后还能排出一条让所有进程完成的安全序列”。}
\exam{常考 Available/Max/Allocation/Need 的含义、试分配步骤，以及安全序列的完整书写。}


\concept{银行家算法做题顺序}
\begin{points}
  \item 先检查 $Request_i\le Need_i$，否则请求本身非法。
  \item 再检查 $Request_i\le Available$，否则资源不足，进程等待。
  \item 只有前两步都通过，才允许试分配并修改 Available、Allocation、Need。
  \item 最后做安全性检查，若找不到安全序列，则撤销试分配。
\end{points}
\plain{老师在录音里反复强调的就是：别一上来就算安全序列，先做两道合法性检查。}
\exam{大题最容易扣分的点就是漏写前两步检查，或只写“安全”却不写安全序列。}


\concept{银行家算法与死锁检测的区别}
\begin{points}
  \item 银行家算法发生在资源分配之前，目标是避免系统进入不安全状态。
  \item 死锁检测发生在资源已经按通常方式分配之后，目标是判断当前是否已经发生死锁。
  \item 前者看到“不安全”就拒绝这次分配；后者若检测到死锁，还要继续指出死锁进程并考虑解除方案。
\end{points}
\plain{看题先分清是在问“这次能不能批”还是“现在是不是已经卡死”，这一步错了后面全会答偏。}
\exam{老师录音特别提醒要区分银行家算法和死锁检测，题型边界要先判断清楚。}


\concept{资源分配图}
\begin{points}
  \item 资源分配图用圆表示进程，用方框表示资源类型，方框内点表示资源实例。
  \item 请求边：进程 $\rightarrow$ 资源；分配边：资源实例 $\rightarrow$ 进程。
  \item 若每类资源只有一个实例，图中有环等价于发生死锁。
  \item 若资源类型有多个实例，有环只是死锁的必要条件，不一定死锁。
\end{points}
\plain{资源分配图题最易错的一句就是“有环是否一定死锁”，关键要先看资源类型是不是单实例。}
\exam{常考资源分配图中环的判定条件，尤其单实例和多实例资源时结论不同。}


\concept{死锁检测题的书写要求}
\begin{points}
  \item 若题目考死锁检测，要分清它不是银行家算法，不要求“避免进入不安全状态”。
  \item 检测出存在死锁时，答案不能只写“有死锁”，还应指出哪些进程陷入死锁。
  \item 当资源种类和实例较多时，可直接按检测过程判断，不一定非画资源分配图。
\end{points}
\plain{检测题关注“现在是不是已经卡死了”，而不是“这次要不要批资源”。}
\exam{老师录音明确强调：若存在死锁，一定写出发生死锁的进程集合。}


\concept{死锁检测的算法}
\begin{points}
  \item 检测算法与银行家算法中的安全性检查相似，但它处理的是当前已分配后的系统状态，不涉及 \code{Max} 和 \code{Need}，而是看当前 \code{Request} 能否被满足。
  \item 常用数据结构包括：\term{Available} 表示当前可用资源向量，\term{Allocation} 表示各进程已占资源矩阵，\term{Request} 表示各进程当前尚在请求的资源矩阵。
  \item 初始化时令 \code{Work = Available}；若某进程满足 $Allocation_i=0$，可先令 \code{Finish[i]=true}，否则 \code{Finish[i]=false}。
  \item 反复寻找满足 \code{Finish[i]=false} 且 $Request_i \le Work$ 的进程，若找到，就假设其完成并释放资源：\code{Work = Work + Allocation\_i}，再置 \code{Finish[i]=true}。
  \item 若最终仍存在 \code{Finish[i]=false} 的进程，则这些进程无法被当前系统资源链条推进，可判定它们陷入死锁。
\end{points}
\plain{检测算法的直觉是：不断试着找“现在就能被喂饱并完成”的进程，把它完成后释放出来的资源再拿去喂别人；最后一直喂不动的那批，就是被卡死的。}
\exam{常考它和银行家算法的区别：银行家算法看“试分配后安不安全”，检测算法看“在当前分配结果下还有没有进程能继续推进”。}


\concept{死锁检测与解除}
\begin{points}
  \item 死锁检测允许资源正常分配，周期性检查是否存在死锁。
  \item 检测后可通过撤销进程、回滚进程、剥夺资源等方式解除死锁。
  \item 选择牺牲进程时可考虑优先级、已运行时间、已占资源、还需资源、回滚代价等。
  \item 缺点是检测和恢复都有代价，且可能造成进程饥饿。
\end{points}
\plain{检测题要答到“查出来以后怎么办”，否则就只答了半题。}
\exam{常考撤销进程、资源剥夺、回滚恢复三类解除思路及其代价。}


\concept{综合方法}
\begin{points}
  \item 实际系统常把预防、避免、检测和人工干预结合使用，而不是全盘只用一种策略。
  \item 对关键且代价高的资源可用较严格策略，对一般资源可放宽限制，提高整体利用率。
\end{points}
\plain{现实系统不是教材单选题，往往是“重要资源严管，普通资源宽管”。}
\exam{常考综合方法的思想：在安全和效率之间折中。}

\chapter{存储器管理}

\section{存储器管理基本概念}
\concept{管理对象与目标}
\begin{points}
  \item 存储管理的对象是主存储器，即内存。
  \item 主存通常分为系统区和用户区：系统区存放内核程序与数据结构，用户区存放用户进程代码和数据。
  \item 目标：为用户提供方便、安全、充分大的存储空间，支持大型应用和系统程序运行。
\end{points}
\plain{这一节先抓总目标：既要让进程“觉得内存够用又好用”，又要保证不同进程彼此隔离，还要让物理内存尽量被充分利用。}
\exam{简答题常直接问“存储器管理的目标是什么”，按“方便、安全、充分大”三个关键词展开即可。}


\concept{存储器管理四个功能}
\begin{points}
  \item 分配与回收：为进程分配内存，进程结束或换出时回收。
  \item 抽象与映射：让进程看到连续的逻辑地址空间，并把逻辑地址转换成物理地址。
  \item 保护与共享：防止进程非法访问他人空间，同时允许受控共享。
  \item 存储扩充：在逻辑上提供比实际物理内存更大的存储空间。
\end{points}
\plain{这四项功能基本对应后面整章内容：连续/非连续分配解决“怎么分”，地址变换解决“怎么映射”，保护共享解决“怎么不乱”，虚拟存储解决“怎么扩充”。}
\exam{选择题喜欢把某个机制和某个功能错配，比如把分页说成只负责保护不负责映射，要会对应。}


\concept{逻辑地址与物理地址}
\begin{points}
  \item 逻辑地址是 CPU 运行的进程所看到的地址，也叫虚地址或相对地址。
  \item 物理地址是内存中的实际地址，也叫实地址或绝对地址。
  \item 地址变换/地址映射/重定位是把逻辑地址转换为物理地址的过程。
  \item MMU 是常见的硬件支持，负责地址转换和合法性检查。
\end{points}
\plain{进程写代码时面对的是逻辑地址；真正访问内存前，硬件还要把它翻成物理地址。MMU 不只做转换，还会顺手检查“这个地址是不是该进程能访问的”。}
\exam{这部分既会单独考概念，也会作为后面分页、分段、请求分页地址题的前置基础。}


\concept{地址绑定时机}
\begin{points}
  \item 编译时绑定：若装入地址已知，编译器可直接生成绝对地址；位置改变则需重新编译。
  \item 装入时绑定：编译生成可重定位代码，装入时确定物理地址。
  \item 运行时绑定：进程运行中仍可移动，地址在执行时由硬件转换，现代分页/分段常用。
  \item 从源程序到可执行映像通常要经历编译、汇编、链接、装载几个阶段，地址在不同阶段的表示形式不同。
\end{points}
\plain{记忆关键不是死背三个名字，而是分清“什么时候最终确定物理地址”：编译时最早、装入时其次、运行时最灵活。}
\exam{常考“哪种方式需要重新编译”“哪种方式需要硬件支持”“哪种方式允许程序运行时移动”。}


\concept{静态重定位与动态重定位}
\begin{points}
  \item 静态重定位在程序装入内存时一次性完成地址修改，运行中不能随意移动。
  \item 动态重定位在程序运行时进行，每次访问内存时由重定位寄存器或 MMU 转换。
  \item 动态重定位支持进程换入换出、移动和虚拟存储，但需要硬件支持。
  \item 动态重定位的基本形式是：物理地址 $=$ 基址寄存器值 $+$ 逻辑地址。
\end{points}
\plain{静态重定位是“装入时一次改完地址”；动态重定位是“运行时每访问一次都现算一次地址”，所以后者才能支持进程搬家。}
\exam{简答题要能写出两者在“变换时机、是否可移动、是否需硬件支持”上的区别。}


\concept{绝对装入、静态重定位装入与动态重定位装入}
\begin{points}
  \item 绝对装入对应编译时绑定，目标代码中已经是确定物理地址；若装入位置变化，通常必须重新编译。
  \item 静态重定位装入对应装入时绑定，装入程序一次性把逻辑地址改成物理地址，装入后地址固定不变。
  \item 动态重定位装入对应运行时绑定，装入时不改逻辑地址，执行时靠重定位寄存器或 MMU 完成转换。
  \item 三者灵活性依次增强，但对硬件支持的要求也逐步提高。
\end{points}
\plain{这三个“装入方式”本质上就是前三种地址绑定时机在装入阶段的具体落地形式，题目换叫法也要能对应上。}
\exam{选择/简答常混着考“编译时绑定=绝对装入”“装入时绑定=静态重定位装入”“执行时绑定=动态重定位装入”。}


\concept{程序链接方式}
\begin{points}
  \item 静态链接：在程序运行前，把各目标模块和所需库函数链接成一个完整装入模块。
  \item 装入时动态链接：装入内存时再完成部分链接工作，便于多个程序共享同一组库模块。
  \item 运行时动态链接：程序执行过程中按需装入并链接目标模块或共享库，灵活性最高。
  \item 动态链接有利于节省内存和便于库升级，但运行时需要额外的链接支持机制。
\end{points}
\plain{链接解决的是“多个目标模块怎么拼成一个能运行的整体”，重定位解决的是“这个整体放到哪里、地址怎么变”。两者别混。}
\exam{简答题里容易把“动态链接”和“动态重定位”混写，前者处理模块装配，后者处理地址转换。}


\section{连续分配管理}
\concept{单一连续分配}
\begin{points}
  \item 内存用户区一次只装入一道用户程序，用户程序独占用户区。
  \item 实现简单，适合早期单道系统。
  \item 缺点是内存利用率低，不能支持多道程序并发。
\end{points}
\plain{这是最简单的连续分配方式，优点只有“简单”，代价是用户区大部分时间容易空着，吞吐量很低。}
\exam{多见于概念对比题，考它为什么只适合早期单道系统。}


\concept{固定分区分配}
\begin{points}
  \item 系统启动时把内存划分为若干固定大小分区，每个分区装入一个进程。
  \item 分区大小可相等，也可不等。
  \item 优点：实现简单，支持多道程序。
  \item 缺点：分区内部未用空间形成内部碎片；分区数限制多道程序道数。
\end{points}
\plain{固定分区把“分配决策”提前到开机时完成，运行时管理容易，但进程大小和分区大小很难总是正好匹配。}
\exam{常考它为什么产生内部碎片，以及为什么分区数会限制并发进程数。}


\concept{动态分区分配}
\begin{points}
  \item 不预先固定分区，进程装入时根据需要动态划分连续内存区。
  \item 优点：比固定分区灵活，内部碎片少。
  \item 缺点：随着分配和回收，会产生外部碎片，需要紧凑/拼接。
\end{points}
\plain{动态分区是按需切块，所以比固定分区更“省料”，但空闲区会越来越碎，最后常卡在“总量够、连续块不够”。}
\exam{算法题多从这里出，尤其是 FF/NF/BF/WF 的分配结果和外部碎片判断。}


\concept{内部碎片与外部碎片}
\begin{points}
  \item 内部碎片：已经分配给进程但进程用不完的空间，位于已分配区域内部。
  \item 外部碎片：空闲空间总量足够，但分散在各处，无法形成一个足够大的连续区域。
  \item 固定分区和分页容易有内部碎片；动态分区和分段容易有外部碎片。
  \item 直觉：内部碎片是“房间给你了但你没住满”；外部碎片是“空房很多但都太零散，放不下大团队”。
\end{points}
\plain{判断碎片类型时就看“浪费空间是在已分配块里面，还是散落在空闲块之间”；位置不同，处理办法也不同。}
\exam{选择题最爱反着说，尤其把“分页有外部碎片、动态分区有内部碎片”作为干扰项。}


\concept{动态分区分配算法}
\begin{points}
  \item 首次适应 FF：从低地址开始找到第一个足够大的空闲分区。速度较快，但低地址碎片较多。
  \item 循环首次适应 NF：从上次查找结束位置继续查找，分布较均匀。
  \item 最佳适应 BF：选择能满足要求的最小空闲分区，尽量减少剩余空间，但容易产生很多小碎片。
  \item 最坏适应 WF：选择最大的空闲分区，剩余空间仍较大，但可能破坏大空闲区。
\end{points}
\plain{做题别只背定义，要抓“扫描从哪开始、挑哪块、剩余碎片会变什么样”这三件事。}
\exam{算法题要把每次分配后的空闲分区表写出来，否则很容易在 NF 或 BF 上算错。}


\concept{动态分区回收与合并}
\begin{points}
  \item 进程释放分区后，需要把回收区按地址顺序插回空闲分区表或空闲分区链。
  \item 若回收区上方邻接空闲分区，则与上邻空闲分区合并；若下方邻接空闲分区，则与下邻空闲分区合并。
  \item 若上下两侧都邻接空闲分区，则三者合并成一个更大的连续空闲区。
  \item 若上下都不邻接空闲分区，则为该回收区单独建立新空闲表项。
  \item 因而回收后空闲分区个数可能不变、减少 1，或增加 1，要根据相邻情况判断。
\end{points}
\plain{动态分区回收题本质是在判“回收区左右有没有空闲邻居”，再决定是并到旧块里还是新开一个空闲块。}
\exam{PPT 专门分四种相邻情况讲回收过程，题目常要求更新空闲分区表，并判断回收后空闲分区数怎么变化。}


\concept{伙伴系统}
\begin{points}
  \item 伙伴系统按 $2^k$ 大小管理内存块，申请时分裂大块，释放时若相邻伙伴空闲则合并。
  \item 优点：分裂和合并规则简单，管理效率较高。
  \item 缺点：申请大小不一定正好是 $2^k$，可能产生内部碎片。
\end{points}
\plain{伙伴系统的关键不是“按 2 的幂分配”这句话本身，而是你能不能顺着分裂和回收过程把内存图画出来。}
\exam{老师明确提过会考画图题：给定总内存和若干申请/释放，按分裂与合并顺序画出每一步。}


\concept{伙伴地址公式}
\begin{points}
  \item 若某块大小为 $2^k$，起始地址为 $d$，则其伙伴块大小相同，起始地址与它只相差 $2^k$。
  \item 若以地址 0 为基准编号，可写成：当 $d \bmod 2^{k+1}=0$ 时，伙伴地址为 $d+2^k$；当 $d \bmod 2^{k+1}=2^k$ 时，伙伴地址为 $d-2^k$。
  \item 若整段可分配空间不是从 0 开始，而是从基址 $a$ 开始，则把上式中的 $d$ 改为相对地址 $d-a$ 再判断。
  \item 直觉上就是看当前块处在某个 $2^{k+1}$ 大块的前半段还是后半段：前半段找后半段，后半段找前半段。
\end{points}
\plain{伙伴地址题本质是在做“同一对兄弟块之间的切换”：块大小不变，只把起始地址在同一更大块的前后两半之间来回翻。}
\exam{PPT 单独给了地址公式和二进制地址例题，考试若出“块大小分别为 4 和 16 时伙伴地址是多少”，先确定 $k$，再按前半段/后半段判断即可。}


\section{分页存储管理}
\concept{分页基本思想}
\begin{points}
  \item 将逻辑地址空间划分为大小相等的页 page。
  \item 将物理内存划分为同样大小的页框/页帧 frame。
  \item 以页为单位分配内存，逻辑上连续，物理上可以分散。
  \item 分页取消了作业对物理连续内存的要求，避免外部碎片，但最后一页可能有内部碎片。
\end{points}
\plain{分页的出发点就是把“必须找一整块连续内存”改成“找若干个页框就行”，所以它主要解决的是连续分配下的外部碎片问题。}
\exam{本章重点在分页，至少要会说清楚它为什么能避免外部碎片、为什么又还会留下内部碎片。}


\concept{页表}
\begin{points}
  \item 系统为每个进程建立一张页表，记录逻辑页号到物理块号的映射。
  \item 每个页面对应一个页表项，页表项通常包含物理块号和标志位。
  \item 标志位可能包括存在位、修改位、引用位、保护位等。
  \item 进程切换时，页表基址和长度等信息会被装入页表寄存器。
\end{points}
\plain{页表就是分页系统里的“翻译词典”：逻辑页号是索引，页表项给出物理块号和访问控制信息。}
\exam{除地址计算外，还常考页表是不是系统共享一张，以及页表项里常见标志位的含义。}


\concept{分页地址结构与转换}
\begin{points}
  \item 逻辑地址由页号 $p$ 和页内偏移 $d$ 组成。
  \item 页大小为 $2^n$ 字节时，低 $n$ 位是页内偏移，高位是页号。
  \item 地址转换：用页号查页表得到物理块号 $f$，物理地址为 $f\times$ 页大小 $+d$。
  \item 若页表项无效或越界，则产生地址越界或缺页异常。
\end{points}
\plain{分页地址题就是三步：拆页号和偏移，查页表，拼物理地址。}
\exam{步骤题要把“先比页号和页表长度、越界则中断”也写出来，老师 PPT 专门强调过合法性检查。}

\concept{分页系统中的碎片}
\begin{points}
  \item 分页消除了外部碎片，因为页框可离散分配。
  \item 但最后一页往往装不满，因此仍可能存在内部碎片。
  \item 内部碎片通常不超过一页大小。
\end{points}
\plain{分页把“大块连续放不下”的问题解决了，但最后那一页角落里可能还有点空着。}
\exam{常考分页有没有外部碎片、为什么还有内部碎片。}

\concept{快表 TLB 与有效访问时间}
\begin{points}
  \item 快表 TLB 是页表项的高速缓存，用于减少访问页表的内存次数。
  \item 无快表时，访问一次数据通常需先访问页表，再访问数据，共两次内存访问。
  \item 有快表时，若命中，直接得到物理块号，只需访问数据；若未命中，需要访问页表再访问数据。
  \item 有效访问时间公式常写为：$EAT=\alpha(t_{TLB}+t_M)+(1-\alpha)(t_{TLB}+2t_M)$，其中 $\alpha$ 为命中率。
\end{points}
\plain{TLB 的作用就是把“每次都查内存里的页表”变成“尽量先查一个更快的小缓存”，所以 EAT 题本质上是在算命中和未命中的加权平均。}
\exam{EAT 是老师点名的重点，做题先分清命中路径和未命中路径各访问几次内存，再代公式。}


\concept{页表实现要点}
\begin{points}
  \item 进程创建时初始化页表，并把页表基址和页表长度保存到 PCB 中。
  \item 进程切换时，内核把当前进程的页表基址、长度等信息装入相应寄存器。
  \item 因而页表是“每个进程一张”，不是整个系统只共用一张。
\end{points}
\plain{页表本质上是“这个进程自己的地址翻译字典”，切换进程就得切换这本字典。}
\exam{常考页表属于谁、什么时候切换页表寄存器。}


\concept{多级页表}
\begin{points}
  \item 多级页表把页表本身分页，只为实际用到的页表页分配内存。
  \item 解决页表过大、必须连续存放的问题。
  \item 缺点：地址转换需要多次查表，若无 TLB 支持会增加访问开销。
  \item 做题时根据每级索引位数判断每级表项数，根据页大小和表项大小判断一页能装多少表项。
\end{points}
\plain{多级页表是“用时间换空间”：多查几次表，换来页表本身不必一次性整块常驻内存。}
\exam{计算题常把“页大小、表项大小、逻辑地址位数”混在一起考，要先算一页能容纳多少表项，再反推需要几级。}

\concept{分页存储管理的保护与共享}
\begin{points}
  \item 地址越界保护：将逻辑地址中的页号与页表长度比较，页号越界则产生中断。
  \item 存取控制保护：在页表项中设置保护位，限制页面的读、写、执行等访问方式。
  \item 数据共享：不同进程可把各自页表中的某些页号映射到同一物理页框，实现共享数据页。
  \item 代码共享：若共享代码页中含有逻辑地址，则不同进程通常需把该代码页放在各自逻辑空间中的相同页号处。
\end{points}
\plain{分页下的共享本质是“多张页表指向同一页框”；分页下的保护本质是“先防越界，再靠页表项权限位限制怎么访问”。}
\exam{简答题若问分页如何实现保护与共享，答题结构直接写“越界保护、存取控制保护、数据共享、代码共享条件”最稳。}


\section{分段与段页式存储管理}
\concept{分段基本思想}
\begin{points}
  \item 分段按程序的逻辑结构划分地址空间，如代码段、数据段、栈段、共享库段。
  \item 逻辑地址由段号和段内偏移组成。
  \item 段表项包含段基址、段长和保护信息。
  \item 地址转换时先检查段号和段内偏移是否合法，再由段基址加偏移得到物理地址。
\end{points}
\plain{分段这节最该抓的是“先看逻辑意义再谈地址”：代码段、数据段、栈段天然分开，所以共享和保护都更符合程序员直觉。}
\exam{常考分段地址题时先做越界判断，再做“段基址 + 段内偏移”；若一上来只会相加，通常会丢步骤分。}

\concept{分段保护和共享}
\begin{points}
  \item 地址越界保护包括两层：先比较段号与段表长度，再比较段内位移与该段段长。
  \item 存取控制保护：可按段统一设置读、写、执行权限，因为同一段内信息通常具有相同访问属性。
  \item 分段天然便于共享，因为段本身就是逻辑信息单位。
  \item 共享实现方式：让多个进程的段表项指向同一段在内存中的起始地址即可。
\end{points}
\plain{分段比分页更适合讲“共享”和“权限”，因为它按代码段、数据段、栈段这类逻辑单位划分，管理粒度更贴合程序结构。}
\exam{问“为什么分段比分页更便于实现共享和保护”时，要抓住“段是逻辑单位、可按段统一赋权、多个段表项可共指同一内存段”三点。}


\concept{分页与分段比较}
\begin{longtable}{p{0.20\textwidth}p{0.36\textwidth}p{0.36\textwidth}}
\toprule
比较项 & 分页 & 分段 \\
\midrule
划分依据 & 系统按固定大小机械划分 & 用户程序逻辑结构划分 \\
单位大小 & 页大小固定 & 段大小可变 \\
地址空间 & 一维地址空间 & 二维地址空间：段号+段内偏移 \\
碎片 & 无外部碎片，有内部碎片 & 无内部碎片或较少，有外部碎片 \\
共享保护 & 不如分段自然 & 以逻辑段为单位，便于共享与保护 \\
\bottomrule
\end{longtable}

\concept{段页式管理}
\begin{points}
  \item 段页式先把程序按逻辑分段，再把每个段分页。
  \item 地址结构通常为段号、段内页号、页内偏移。
  \item 每个段有段表项，段表项指向该段对应的页表。
  \item 优点：兼具分段的逻辑保护共享和分页的离散分配；缺点是地址转换更复杂。
\end{points}
\plain{段页式不要背成三个名词堆在一起，它本质是在说：先借“段”保留逻辑结构，再借“页”解决离散分配问题。}
\exam{常考段页式地址结构和查表顺序，答题时最好明确写出“段号 -> 段表 -> 页表 -> 物理块号 -> 页内偏移”。}

\concept{段页式地址题的做法}
\begin{points}
  \item 先从逻辑地址中拆出段号、段内页号、页内偏移。
  \item 用段号查段表，得到该段页表的起始位置并检查段是否合法。
  \item 再用段内页号查该段对应页表，得到物理块号。
  \item 最后由物理块号和页内偏移拼出物理地址。
\end{points}
\plain{段页式就是“先定位哪一段，再定位这一段里的哪一页”，顺序不能反。}
\exam{这类题老师录音特别提过，步骤一定要写完整：拆地址、查段表、查页表、拼地址。}

\chapter{虚拟存储器}

\section{虚拟存储器概念}
\concept{为什么需要虚拟存储器}
\begin{points}
  \item 传统物理存储管理要求作业运行前全部装入内存，运行期间常驻内存，限制了大作业和多作业运行。
  \item 物理扩充内存会增加成本，因此需要从逻辑上扩充内存。
  \item 虚拟存储器打破“作业必须全部装入内存才能运行”的限制。
\end{points}
\plain{虚拟存储器首先解决的是“内存装不下整个作业怎么办”这个问题：不再要求程序一开始全部进内存，而是先装一部分，边运行边调入。}
\exam{简答题常从“为什么需要虚拟存储器”切入，答案核心是：突破全部装入限制、逻辑扩充内存、降低物理扩容成本。}


\concept{局部性原理}
\begin{points}
  \item 程序执行时，CPU 访问的指令和数据倾向于聚集在较小范围内。
  \item 时间局部性：刚访问过的指令或数据，短时间内可能再次访问。
  \item 空间局部性：访问某地址后，附近地址很可能被访问。
  \item 顺序局部性：程序大多数情况下顺序执行，下一条指令通常紧接当前指令。
  \item 局部性原理是虚拟存储器的理论基础。
\end{points}
\plain{程序一段时间只爱访问附近的一小块代码和数据，所以没必要一开始全装入内存。}
\exam{除了时间/空间局部性定义，还要会说它为什么能支持请求分页和工作集模型。}


\concept{虚拟存储器定义与特征}
\begin{points}
  \item 虚拟存储器是具有请求调入和置换功能、能从逻辑上扩充内存容量的存储器系统。
  \item 它是一种以时间换空间的技术：用缺页中断、磁盘换入换出换取更大的逻辑空间。
  \item 特征：离散性、多次性、对换性、虚拟性。
  \item 容量受地址结构和外存容量限制，不是无限大。
\end{points}
\plain{“虚拟”不是凭空变出内存，而是把程序访问地址和真实物理内存解耦，再靠外存做后备仓库。}
\exam{四特征是高频名词解释；回答时最好逐个落到机制上：离散对应非连续分配，多次对应分批装入，对换对应换入换出，虚拟对应逻辑扩充。}


\concept{虚拟存储器的好处与实现基础}
\begin{points}
  \item 好处：可以运行更大的程序，因为不再要求程序整体常驻内存。
  \item 好处：可以同时运行更多程序，因为每个程序通常只需把活跃页面保留在内存中。
  \item 好处：程序启动更快，因为初始装入的页面更少。
  \item 物质基础：相当数量的外存、一定容量的内存，以及能动态完成地址转换的硬件机构。
\end{points}
\plain{这部分可以连起来记：外存负责“兜底”，内存负责“放当前活跃部分”，地址变换机构负责“把逻辑访问导向当前真正所在的位置”。}
\exam{PPT 明确提了“好处”和“实现基础”，简答题不要只背定义而漏掉条件。}


\section{请求分页存储管理}
\concept{请求分页的基本思想}
\begin{points}
  \item 程序运行前只装入部分页面即可启动。
  \item 运行中访问的页面若不在内存，产生缺页中断，由 OS 调入所需页面。
  \item 若内存无空闲页框，需要选择某个页面换出。
  \item 请求分页 = 基本分页 + 请求调页 + 页面置换。
\end{points}
\plain{请求分页和基本分页的根本差别在于：基本分页默认页已经在内存，请求分页允许“访问时才发现页不在”，然后现调。}
\exam{这正是老师点名的比较点，答题时一定写出“允许页面不在内存”和“缺页中断是正常机制”这两句。}


\concept{请求分页系统的组成}
\begin{points}
  \item 请求分页建立在分页系统基础上，还需要缺页中断机构。
  \item 需要扩充页表项，加入存在位、修改位、访问位、外存地址等信息。
  \item 还需要请求调页和页面置换的软件支持。
  \item 其硬件核心仍是 MMU，它负责地址转换、合法性检查，并与页表/TLB 配合工作。
\end{points}
\plain{请求分页不是只靠一个“缺页中断”词就结束了，它还要有页表扩展和调页/置换机制一起配合。}
\exam{常考“请求分页比基本分页多了什么支持机构”。}


\concept{MMU 在请求分页中的作用}
\begin{points}
  \item MMU 接受 CPU 给出的逻辑地址，完成页号与页内位移的分解，并把逻辑地址转换成物理地址。
  \item MMU 会先结合 TLB 和页表检查该页是否在主存中，以及访问是否合法。
  \item 若页表表明该页不在内存，MMU 触发缺页异常，由操作系统调页。
  \item 因而 MMU 既负责地址转换，也承担内存保护和虚拟内存支持的重要硬件职责。
\end{points}
\plain{可以把 MMU 看成 CPU 和内存之间的“翻译兼门卫”：先翻译地址，再判断能不能进。}
\exam{PPT 对 MMU 单独成节，说明它不只是第 8 章的旧概念复读，这里要会把 MMU 和缺页、TLB、保护联系起来。}


\concept{基本分页与请求分页区别}
\begin{points}
  \item 基本分页要求进程运行所需页面已在内存中，地址转换时页表项应有效。
  \item 请求分页允许部分页面不在内存，访问未驻留页面时产生缺页中断。
  \item 因此基本分页本身不以缺页中断为常规机制；请求分页才把缺页作为正常运行过程的一部分。
\end{points}
\plain{第 8 章基本分页讲的是“怎么按页映射”；第 9 章请求分页讲的是“页不在内存时怎么办”。}
\exam{比较题回答时别只写“一个高级一个低级”，要落到“是否允许页面暂不在内存、是否需要缺页中断和置换机制”。}


\concept{页表项扩展}
\begin{points}
  \item 存在位/有效位：表示该页是否在内存。
  \item 修改位/脏位：表示页面是否被写过，换出时若为脏页需写回外存。
  \item 引用位/访问位：表示页面近期是否被访问，用于页面置换算法。
  \item 外存地址：指出页面在外存中的位置。
\end{points}
\plain{请求分页下的页表不再只是“页号到块号”的映射表，而是还要告诉 OS 这页在不在内存、最近用没用、改没改、若不在内存该去磁盘哪找。}
\exam{填空/选择最爱考存在位、修改位、访问位、外存地址各自干什么，尤其修改位决定换出时是否必须回写。}


\concept{缺页中断处理过程}
\begin{points}
  \item CPU 访问页面，发现页表项存在位为 0，产生缺页中断。
  \item 操作系统检查访问是否合法；若非法，终止进程。
  \item 若合法，查找空闲页框；若无空闲页框，执行页面置换。
  \item 若被换出页被修改过，需要写回外存。
  \item 从外存读入所需页面，更新页表和快表，重新执行引起缺页的指令。
\end{points}
\plain{这类流程题的核心顺序是：发现不在内存 -> 判断是否合法 -> 找页框/选牺牲页 -> 必要时回写 -> 调入新页 -> 改页表/TLB -> 重执行原指令。}
\exam{老师强调过请求分页地址变换，答题时不要漏掉“重新执行产生缺页的指令”，这正是它和一般中断返回到下一条指令的重要区别。}


\concept{缺页中断与一般中断的区别}
\begin{points}
  \item 缺页中断在指令执行期间发生，而一般外部中断常在一条指令执行结束后响应。
  \item 一条指令执行过程中可能涉及多个地址访问，因此可能产生多次缺页中断。
  \item 缺页中断处理完成后，要重新执行引起缺页的那条指令；一般中断返回后通常执行下一条指令。
\end{points}
\plain{缺页中断不是“这条指令做完了再说”，而是这条指令做到一半发现缺页，必须先把缺的页补齐，回来继续把这条指令做完。}
\exam{这是 PPT 明确单列的区别点，容易出选择或简答。}


\concept{请求分页中的地址变换过程}
\begin{points}
  \item CPU 给出逻辑地址后，MMU 先拆分出页号和页内位移。
  \item 先查快表 TLB；若命中且表项有效，可直接形成物理地址访问内存。
  \item 若 TLB 未命中，则进一步查页表；若页在内存，形成物理地址并可把表项装入 TLB。
  \item 若页表显示该页不在内存，则产生缺页中断，调页成功后再更新页表/TLB 并重执行指令。
\end{points}
\plain{请求分页地址变换比基本分页只多了一条分支：查到“页不在内存”时，不是直接报错，而是进入缺页处理。}
\exam{步骤题建议写成“拆地址 -> 查 TLB -> 查页表 -> 缺页则中断 -> 调页 -> 更新 -> 重执行”，过程比一句结论更重要。}


\concept{TLB 与刷新问题}
\begin{points}
  \item TLB 是页表项的高速缓存，用于加快地址转换。
  \item 进程切换时，旧进程的地址映射可能失效，因此通常需要刷新 TLB 或切换到对应地址空间标识。
  \item 请求调页或页面置换后，如果被替换页的旧页表项仍留在 TLB 中，也必须使其失效或更新。
  \item 因而 TLB 命中并不等于“永远正确”，它必须与当前页表状态保持一致。
\end{points}
\plain{TLB 快是快，但前提是里面缓存的映射没过期；进程切换和换页都会让旧缓存失真，所以要刷新。}
\exam{PPT 专门提了 TLB 刷新，说明考试可能不只考“TLB 提高速度”，还会考“什么时候必须失效/刷新”。}


\concept{缺页次数、缺页率与缺页开销}
\begin{points}
  \item 缺页次数指页面引用过程中发生缺页的总次数。
  \item 缺页率通常指缺页次数占总页面访问次数的比例。
  \item 缺页开销远高于普通内存访问，因为往往需要磁盘 I/O。
\end{points}
\plain{一次缺页不是“多查一步表”这么简单，而是可能要从磁盘把整页搬进来，代价大得多。}
\exam{常考为什么页面置换算法和页框分配策略会显著影响系统性能。}


\concept{有效访问时间 EAT}
\begin{points}
  \item 若页在主存且 TLB 命中，访问时间约为查 TLB 时间加一次内存访问时间。
  \item 若页在主存但 TLB 未命中，则通常还需访问页表并更新 TLB，开销高于命中情况。
  \item 若页不在主存，还需加上缺页中断处理与调页时间，因此开销远大于普通内存访问。
  \item EAT 本质上是按 TLB 命中率、缺页率对几种访问路径求加权平均。
\end{points}
\plain{EAT 题不要背成死公式，先把路径画清楚：命中怎么走、未命中但不缺页怎么走、缺页怎么走，再做概率加权。}
\exam{老师在第 8 章强调过 EAT，第 9 章 PPT 又把“含缺页率的 EAT”展开了，计算题要注意把缺页处理时间单独算进去。}


\concept{请求分页的优缺点}
\begin{points}
  \item 优点：页可离散存放，减少移动和紧凑开销，基本消除外部碎片。
  \item 优点：主存利用率更高，有利于提高多道程序度和系统吞吐量。
  \item 缺点：需要硬件支持和缺页中断处理，系统实现更复杂。
  \item 缺点：若缺页频繁，磁盘 I/O 和中断处理开销会明显增加。
\end{points}
\plain{请求分页的收益来自“少装、晚装、按需装”，代价来自“查得更复杂、缺页更昂贵”。}
\exam{简答题适合按“优点两条、缺点两条”组织答案，别只写优点。}


\section{页面置换算法}
\concept{最佳置换 OPT}
\begin{points}
  \item OPT 选择未来最长时间不会被访问的页面淘汰。
  \item 它缺页率最低，但需要预知未来，实际无法实现。
  \item 常作为评价其他算法的理论下界。
\end{points}
\plain{OPT 是标准答案里的“理想参照物”：真实系统做不到，但做题时常用它来说明某个访问串的最少缺页次数。}
\exam{算法题里若同时要求比较多种置换算法，OPT 往往用来作为最优基准。}


\concept{先进先出 FIFO}
\begin{points}
  \item FIFO 淘汰最早进入内存的页面。
  \item 实现简单，只需维护队列。
  \item 缺点是不考虑页面使用频率，可能淘汰仍常用页面。
  \item FIFO 可能出现 Belady 异常：页框数增加，缺页次数反而增加。
\end{points}
\plain{FIFO 只看“谁来得早”，不看“谁最近还常用”，所以很容易把热点页错踢出去。}
\exam{Belady 异常基本就是 FIFO 的招牌考点，要会解释“页框更多却缺页更多”为什么可能发生。}


\concept{最近最久未使用 LRU}
\begin{points}
  \item LRU 淘汰最近最长时间没有被访问的页面。
  \item 基于局部性原理：最近很久未用的页面，未来短期内也可能不用。
  \item 实现可用计数器、栈或近似算法，但精确 LRU 开销较高。
  \item LRU 不会出现 Belady 异常。
\end{points}
\plain{LRU 比 FIFO 更贴合程序局部性，因为它默认“刚用过的页近期大概率还会再用”。}
\exam{题目若问为什么 LRU 通常优于 FIFO，关键词就是“利用局部性原理”。}


\concept{时钟 CLOCK/二次机会算法}
\begin{points}
  \item Clock 用访问位近似 LRU，把页面组成循环队列。
  \item 指针扫描页面：若访问位为 0，淘汰；若为 1，清 0 并给第二次机会。
  \item 改进型 Clock 还考虑修改位，优先淘汰未访问且未修改的页面。
\end{points}
\plain{Clock 的本质是“用一位访问位，低成本地近似 LRU”，牺牲一点精度换来更低实现开销。}
\exam{若题目写“二次机会算法”“简单时钟算法”，都要能和访问位、循环扫描联系起来。}

\concept{页面置换题的标准做法}
\begin{points}
  \item 按页面访问串从左到右逐次处理。
  \item 每一步先判断命中还是缺页，再决定是否置换。
  \item 若算法需要维护额外状态，如 LRU 次序或 Clock 访问位，要在每一步同步更新。
  \item 答题时最好写出每一步页框内容变化，不要只报最终缺页次数。
\end{points}
\plain{这类题最怕心算跳步，一跳步就很容易后面全错。}
\exam{老师录音强调过页面置换经常是大题，过程要完整写。}


\section{页面分配与置换策略}
\concept{页面装入与清除策略}
\begin{points}
  \item 页面装入策略决定何时把页面调入主存，常见有请求式调页和预调页两类。
  \item 请求式调页按需装入，只有访问到缺页时才调入，优点是初始装入少，缺点是运行中可能频繁触发磁盘 I/O。
  \item 预调页根据局部性原理和历史访问情况，预先把若干可能很快会访问的页面装入主存，目的是减少后续缺页次数。
  \item 页面清除策略决定何时把修改过的页面写回外存：请求式清除在页面被选中换出时才写回；预约式清除则提前成批写回脏页。
\end{points}
\plain{这里本质是在回答两个“何时”问题：何时把页装进来，何时把脏页写出去。请求式更省眼前工作，预调页和预约式清除则是想用“提前做一点”换取后续更少的阻塞。}
\exam{简答题若问页面管理策略，别只写置换算法，还要能区分“装入策略”和“清除策略”，并分别说出请求式/预调页、请求式清除/预约式清除。}

\concept{从何处调入页面}
\begin{points}
  \item 文件区用于存放普通文件，通常采用离散分配；对换区用于存放对换页面，通常采用连续分配，因此对换区 I/O 一般更快。
  \item 若系统有足够对换区，可把与进程有关的页面预先复制到对换区，缺页时直接从对换区调入。
  \item 若对换区不足，则不会被修改的页面可直接从文件区调入；可能被修改的页面换出后再写入对换区，后续优先从对换区调回。
  \item UNIX 方式中，未运行过的页面通常从文件区调入；曾运行后被换出的页面，下次一般从对换区调入；若共享页面已在内存，则可直接复用。
\end{points}
\plain{做题时别把“文件区”和“对换区”当成同一种外存位置。文件区更像原始来源地，对换区更像专门为换页准备的中转仓库。}
\exam{若考“缺页时一定从对换区调页吗”，答案不是绝对的，要看系统对换区是否充足、页面是否可修改，以及是否采用 UNIX 式处理。}

\concept{最小物理块数}
\begin{points}
  \item 最小物理块数是保证进程能够正常运行所需的最少页框数，少于该值时进程无法执行。
  \item 若是单地址指令且采用直接寻址，通常至少需要 2 个物理块。
  \item 若是单地址指令且允许间接寻址，通常至少需要 3 个物理块。
  \item 若机器较复杂，指令、源地址和目标地址都可能跨页，则所需最小物理块数会更大。
\end{points}
\plain{这个概念不是在问“运行得快不快”，而是在问“最少得有几块内存，指令和操作数访问才不至于根本做不下去”。}
\exam{选择题喜欢把“最小物理块数”和“分配越多越好”混在一起。前者是能否运行的下限，后者才是性能优化问题。}

\concept{驻留集与页面分配}
\begin{points}
  \item 驻留集指某进程当前驻留内存的页面集合。
  \item 固定分配：进程运行期间分得的页框数固定。
  \item 可变分配：根据缺页率、工作集等动态调整页框数。
\end{points}
\plain{驻留集就是“这个进程当前手里实际占着的页框集合”。}
\exam{常考固定/可变分配、局部/全局置换的组合关系。}


\concept{局部置换与全局置换}
\begin{points}
  \item 局部置换：缺页时只从本进程驻留集中选择页面淘汰。
  \item 全局置换：可从全系统所有可换页面中选择淘汰，进程页框数会动态变化。
  \item 固定分配只能配合局部置换；可变分配可配合局部置换或全局置换。
\end{points}
\plain{区分局部/全局置换时，关键看“缺页时能不能拿别的进程的页框来顶”。}
\exam{这类题是概念题，不是置换算法题，别答成 FIFO/LRU。}


\concept{不能组合的策略}
\begin{points}
  \item 固定分配 + 全局置换通常不能组合。
  \item 原因：固定分配要求每个进程页框数不变，而全局置换可能把其他进程页框拿走或让本进程获得额外页框，改变驻留集大小。
\end{points}
\plain{一个要求“页框数固定不变”，另一个却允许“从别人那儿拿页框”，两者自然冲突。}
\exam{常考为什么固定分配不能配全局置换。}


\section{工作集、抖动与缺页率}
\concept{工作集模型}
\begin{points}
  \item 工作集是在某段时间窗口内进程实际访问过的页面集合。
  \item 工作集反映进程当前局部性需求，系统应尽量给进程分配足够页框容纳其工作集。
  \item 若分配页框少于工作集需要，缺页率会显著上升。
\end{points}
\plain{工作集不是“这个进程所有页面”，而是“它最近这段时间真正活跃、离不开的页面集合”。}
\exam{老师特别强调了工作集的作用和意义，答题时要点出它是利用局部性来估计当前页面需求，从而降低缺页率、避免抖动。}


\concept{工作集算法}
\begin{points}
  \item 系统定期检查进程在窗口长度 $\Delta$ 内访问过的页面，得到当前工作集。
  \item 若某进程工作集装不下现有驻留集，就应为其增配页框；若系统无空闲页框，可暂停部分进程以释放内存。
  \item 核心思想是让每个进程的当前工作集尽量驻留内存，避免刚换出的页又立刻被访问。
\end{points}
\plain{工作集模型偏“描述当前需要多少页”，工作集算法偏“按这个估计结果去动态分配页框”。}
\exam{名词解释或简答题里，要点出时间窗口 $\Delta$、按工作集调整驻留集、必要时降低多道程序度。}


\concept{抖动 thrashing}
\begin{points}
  \item 抖动指系统花费大量时间进行页面换入换出，真正执行用户程序的时间很少。
  \item 现象：CPU 利用率低，磁盘交换区利用率很高，缺页频繁。
  \item 原因：多道程序度过高或每个进程分得页框不足，工作集无法驻留内存。
  \item 解决：降低多道程序度、增加页框、采用工作集策略、改善置换算法。
\end{points}
\plain{抖动不是简单的“缺页多一点”，而是系统大部分时间都耗在换页上，几乎不干正事。}
\exam{要会和 Belady 异常区分：Belady 异常是某些算法下页框增加反而缺页增多；抖动是整体页框紧张导致系统频繁换页。}


\concept{局部抖动与全局抖动}
\begin{points}
  \item 局部抖动：某个进程分到的页框太少，主要是该进程频繁缺页、反复换页。
  \item 全局抖动：系统中许多进程同时页框不足，彼此争抢内存，导致整个系统频繁换页。
  \item 局部问题若继续提高多道程序度或采用不当的全局置换，可能演变成全局抖动。
\end{points}
\plain{局部抖动是“一个进程转不动”，全局抖动是“整个系统都在疯狂换页”。}
\exam{判断题常考“抖动是单个进程现象还是系统现象”，作答时应区分局部和全局两个层次。}


\concept{利用缺页率判断抖动}
\begin{points}
  \item 缺页率持续偏高，说明进程当前驻留集可能小于其工作集需求。
  \item 可设置缺页率上、下界：高于上界时增配页框或暂停部分进程；低于下界时可适当回收页框。
  \item 若许多进程缺页率同时过高，通常表明系统已接近或进入全局抖动。
\end{points}
\plain{缺页率可以看成内存压力表；不是看某一次缺页，而是看一段时间内是不是一直高。}
\exam{简答题容易问“如何发现抖动”，答题时要写出“监视缺页率”而不只是笼统地说“看系统变慢”。}


\concept{影响缺页率的因素}
\begin{points}
  \item 页面大小：过小会增加页表和 I/O 次数，过大会增加内部碎片和无用调入。
  \item 分配页框数：一般页框越多缺页率越低，但 FIFO 可能有 Belady 异常。
  \item 页面置换算法：OPT 最优，LRU 通常较好，FIFO 简单但可能异常。
  \item 程序局部性：局部性越好，缺页率越低。
  \item 多道程序度：过高会导致每个进程页框不足，引发抖动。
\end{points}
\plain{这一节适合按“页面本身大小、给进程多少页框、怎么置换、程序局部性好不好、系统里进程是不是太多”五个维度来记。}
\exam{主要是概念理解题，要求能解释“为什么这个因素会让缺页率升高或降低”，而不是只会列名词。}

\chapter{文件系统}

\section{文件基本概念}
\concept{文件系统的组成}
\begin{points}
  \item 文件系统是操作系统中与管理文件有关的软件和数据的集合。
  \item 组成包括：被管理的文件集合、管理文件所需的数据结构、相应管理软件。
  \item 管理软件负责外存空间分配回收、目录管理、文件读写、共享与保护等。
\end{points}
\plain{文件系统这一章的主线是把“磁盘上的原始块”包装成“用户按名字访问的文件”，中间靠目录、FCB、分配方式、保护机制来衔接。}
\exam{简答题常问文件系统由什么组成，回答时最好对应“文件集合 + 数据结构 + 管理软件”三层。}


\concept{文件定义}
\begin{points}
  \item 文件是记录在外部存储上，具有名字，在逻辑上有完整意义的信息项序列。
  \item 信息项是构成文件内容的基本单位，可以是字节或记录。
  \item 从用户角度看，文件是逻辑外存的最小分配单元。
  \item 文件内容的意义由文件创建者和使用者解释。
\end{points}
\plain{用户看到的是“一个有名字、有内容的逻辑对象”，而不是磁盘块号；文件系统正是负责做这层抽象。}
\exam{名词解释题核心抓“外存上、有名字、逻辑完整、信息项序列”这几个关键词。}


\concept{文件属性}
\begin{points}
  \item 名称：便于用户识别的符号文件名。
  \item 标识符：系统内部唯一标识，常为数字形式。
  \item 类型：普通文件、目录文件、设备文件等。
  \item 位置：文件在设备上的位置指针。
  \item 大小：当前大小和最大允许大小。
  \item 保护：访问控制信息，规定谁能读、写、执行。
  \item 时间、日期、用户标识：创建、修改、访问等元数据。
\end{points}
\plain{这些属性本质上就是文件的元数据，大多保存在目录项/FCB 里，而不是文件正文里。}
\exam{选择题容易混淆“文件内容”和“文件属性”，例如把文件长度、权限、时间戳错当成正文数据。}


\concept{文件分类}
\begin{points}
  \item 按用途可分为系统文件、库文件、用户文件。
  \item 按保护级别可分为只读文件、读写文件、执行文件、不受保护文件。
  \item 按信息流向可分为输入文件、输出文件、输入输出文件。
  \item 按数据形式可分为源文件、目标文件、可执行文件。
\end{points}
\plain{分类题的关键是先看清题目按哪一个维度在分，不要把几种分类标准混成一类。}
\exam{常考选择题：问某文件属于哪一类，或者让你辨别“用途分类”和“保护级别分类”。}


\concept{文件操作}
\begin{points}
  \item Create：创建文件，为文件分配目录项和必要空间。
  \item Write：根据写指针向文件写入数据。
  \item Read：根据读指针从文件读取数据。
  \item Seek：移动文件当前读写位置。
  \item Delete：删除文件，释放空间，删除目录项。
  \item Truncate：删除文件内容但保留属性。
  \item Open/Close：打开文件时把文件控制信息装入内存打开文件表，关闭时释放相关表项。
\end{points}
\plain{`open` 的关键作用不是“开始读文件”这五个字，而是把后续读写反复要用到的控制信息先搬进内存，避免每次都重查磁盘目录。}
\exam{问 `open` 为什么重要时，重点答“减少目录搜索和磁盘访问，建立打开文件表项，保存读写状态”。}


\concept{文件类型}
\begin{points}
  \item 操作系统需要识别并支持不同文件类型，以便按合适方式处理文件。
  \item 常见类型包括普通文件、目录文件、可执行文件、设备文件、链接文件等。
  \item 文件类型信息既可能来自扩展名，也可能来自元数据或文件头。
\end{points}
\plain{文件“类型”不是纯粹给人看的标签，它会影响系统如何解释、打开和保护这个文件。}
\exam{PPT 单独列了文件类型，选择题里常把“分类标准”和“系统识别的文件类型”混在一起考。}


\section{文件访问方法与逻辑结构}
\concept{顺序访问}
\begin{points}
  \item 顺序访问按文件中信息项的先后顺序读写。
  \item 适合文本流、日志、磁带等场景。
  \item 若要访问第 $n$ 个记录，通常需要从前面依次读到它。
\end{points}
\plain{顺序访问的特点是“当前位置依赖前一步”，适合从头到尾处理文件，但不适合频繁跳着找。}
\exam{PPT 里有 \code{read\_next()}、\code{write\_next()}、\code{reset()} 这一类操作，简答时可以用来说明顺序访问接口。}


\concept{直接访问/随机访问}
\begin{points}
  \item 直接访问允许按记录号或块号直接定位到某个位置。
  \item 磁盘支持随机访问，因此文件系统可实现直接访问。
  \item 适合数据库、索引文件、需要频繁定位的数据文件。
\end{points}
\plain{直接访问的关键是能根据编号或偏移快速算出位置，所以定长记录文件尤其适合这种方式。}
\exam{PPT 特别强调了“定长记录便于直接访问，变长记录不便于直接访问”，这个点容易出判断题。}


\concept{索引访问}
\begin{points}
  \item 索引访问本质上建立在直接访问之上，通过索引表把关键字映射到记录号或偏移位置。
  \item 它适合按关键字快速查找记录，如数据库索引。
  \item 对变长记录文件，也可借助索引实现较高效定位。
\end{points}
\plain{索引访问不是顺着文件一条条扫，而是先查“目录/索引”，再跳到目标位置。}
\exam{若题目给出“关键字 -> 记录位置”的描述，通常考的就是索引访问而不是普通直接访问。}


\concept{文件逻辑结构}
\begin{points}
  \item 无结构文件/流式文件：文件被看作字节序列，如普通文本和二进制文件。
  \item 有结构文件/记录式文件：文件由一组记录组成，每条记录有内部字段结构。
  \item 顺序文件：记录按某种顺序排列，适合顺序处理。
  \item 索引文件：建立索引表，从键值映射到记录位置，提高随机查找效率。
  \item 索引顺序文件：主文件按键有序，同时建立分组索引，兼顾顺序处理和快速定位。
  \item 散列文件：通过散列函数直接计算记录位置，适合等值查找。
\end{points}
\plain{逻辑结构回答的是“用户看来文件内部怎么组织”；它和磁盘上具体放在哪些块里是两回事。}
\exam{题目若问“流式文件和记录式文件区别”，不要答成连续分配/链接分配，那已经变成物理结构了。}


\concept{文件的物理结构}
\begin{points}
  \item 文件物理结构又称存储结构，描述文件逻辑块如何映射到磁盘物理块。
  \item 它与存储设备特性和外存分配方式密切相关。
  \item 常见物理结构有顺序结构、链接结构、索引结构。
\end{points}
\plain{逻辑结构是“给用户看”的，物理结构是“给操作系统和磁盘打交道”的。}
\exam{文件逻辑结构和物理结构是经典对比点，答题时一定分清“与用户观点有关”还是“与外存组织有关”。}


\concept{顺序结构、链接结构与索引结构}
\begin{points}
  \item 顺序结构/连续结构：文件逻辑上相邻的块在磁盘上也连续，顺序访问快，但容易产生碎片，不利于动态扩充。
  \item 链接结构/串联结构：每块指向下一块，分配灵活、便于增长，但通常只适合顺序访问，且指针占空间。
  \item 索引结构：为文件建立索引表记录逻辑块与物理块映射，既支持顺序访问也支持随机访问，但有额外索引开销。
\end{points}
\plain{这三种结构本质上是在“访问效率、扩展灵活性、空间开销”三者之间做取舍。}
\exam{判断题喜欢把“链接结构便于随机访问”或“顺序结构不会产生碎片”写成陷阱。}


\section{目录结构}
\concept{磁盘结构的几个常识}
\begin{points}
  \item 磁盘通常按\key{盘块/物理块}为单位进行信息传输和空间分配，文件最终也要落实到这些磁盘块上。
  \item 一个磁盘块内部可以存放若干个逻辑记录；反过来，一个较大的逻辑记录也可能跨越多个物理块。
  \item 目录、FCB/inode、文件数据块等都只是磁盘上不同用途的数据块，只是内容和作用不同。
  \item 文件系统讨论“目录结构、物理结构、空闲块管理”时，默认背景就是这些数据最终都存放在外存块中。
\end{points}
\plain{这块不是重复大容量存储系统那章的磁头细节，而是先把文件系统这一章默认使用的“块”概念讲清：文件、目录、索引项最后都要落到磁盘块上。}
\exam{PPT 在目录结构前单独补了这层背景，选择题或简答题里若问“文件分配和传输的基本单位是什么”，通常就落在盘块/物理块。}

\concept{目录与目录项}
\begin{points}
  \item 目录用于组织系统中的文件并提供文件信息。
  \item 目录项通常包含文件名和文件控制块指针或 inode 号。
  \item 目录检索就是根据路径名逐级查找目录项，找到目标文件控制信息。
\end{points}
\plain{目录就是文件名到文件控制信息的映射表，路径就是沿目录树找到文件。}
\exam{常考一级/两级/树形目录优缺点、平均查找磁盘访问次数。}


\concept{目录管理的基本操作}
\begin{points}
  \item 目录上的基本操作通常包括：搜索文件、创建文件、删除文件、查看目录、重命名文件、遍历文件系统。
  \item 这些操作本质上都是对“文件名到目录项映射表”的查找、插入、删除或修改。
  \item 树形目录下删除目录时还要考虑删除策略：有的系统要求目录非空时不能删除，有的系统允许连同子目录和文件一起删除。
\end{points}
\plain{目录管理题不要只会说“目录能存文件名”，还要知道系统实际要支持查找、增删、改名、遍历这些操作，以及删除目录时的规则差异。}
\exam{PPT 在目录结构前单独列了目录基本操作和删除策略，选择题常考“非空目录能不能删”以及重命名/遍历属于哪类目录操作。}


\concept{分区、卷、FCB 与目录文件}
\begin{points}
  \item 分区是对物理硬盘空间的切分；原始分区需格式化生成文件系统后才能正常存放文件。
  \item 卷是已经安装好文件系统、可供读写的逻辑单元；常见场景下一个分区对应一个卷。
  \item 文件控制块 FCB 是保存文件属性和物理位置等信息的数据结构，文件与 FCB 一一对应，是文件存在的标志。
  \item 文件目录本质上是 FCB 的集合；目录文件则是把这些目录信息也作为文件保存在磁盘上。
\end{points}
\plain{这几个概念常连着考：分区/卷说明文件系统“长在什么上面”，FCB/目录说明文件系统“怎么记住每个文件”。}
\exam{PPT 对 FCB 和目录文件讲得很实，简答时可直接写“文件 = 文件体 + 文件说明（FCB）”。}


\concept{一级目录}
\begin{points}
  \item 系统只有一个目录，所有文件都放在其中。
  \item 优点：实现简单。
  \item 缺点：文件多时查找慢，不能解决重名问题，不便于多用户管理。
\end{points}
\plain{单级目录的最大问题不是“难看”，而是所有用户、所有文件都挤在一个表里，规模一大就很难管。}
\exam{答缺点时至少写出“查找慢、不能重名、不利于多用户”。}


\concept{两级目录}
\begin{points}
  \item 每个用户有自己的用户文件目录 UFD，系统有主文件目录 MFD。
  \item 解决不同用户之间文件重名问题。
  \item 用户内部文件多时仍不够灵活。
\end{points}
\plain{二级目录解决的是“不同用户同名文件冲突”，但对同一用户内部的大量分类管理帮助仍有限。}
\exam{PPT 还提到它不利于合作共享，这也是两级目录常见缺点。}


\concept{树形目录与路径}
\begin{points}
  \item 树形目录把目录组织成层次结构，便于分类管理。
  \item 绝对路径从根目录开始，如 \code{/home/a.txt}。
  \item 相对路径从当前目录开始。
  \item 优点：层次清晰、重名控制灵活、便于权限管理。
\end{points}
\plain{树形目录的关键是“路径”与“当前目录”两个概念：同一个文件既能用绝对路径定位，也能结合当前目录用相对路径定位。}
\exam{PPT 还提了特殊目录 `.` 和 `..`，以及绝对/相对路径的区别，这些很适合出选择题。}

\concept{当前目录与特殊目录项}
\begin{points}
  \item \term{当前目录}（working directory）是用户在某一段时间内默认工作的目录。
  \item 相对路径不是从根目录出发，而是从当前目录出发解释。
  \item 特殊目录项 \code{.} 表示当前目录本身，\code{..} 表示当前目录的父目录。
  \item 这两个特殊目录项使用户能在目录树中方便地进行相对定位与层次切换。
\end{points}
\plain{当前目录可以理解成“你现在站在哪个目录里”；有了它，很多路径就不必每次都从根目录写全。}
\exam{选择题常直接问 \code{.} 和 \code{..} 各表示什么，或问相对路径是相对于根目录还是相对于当前目录。}


\concept{无环图目录与共享}
\begin{points}
  \item 无环图目录允许多个目录项共享同一个文件或子目录，但不允许形成环。
  \item 共享文件需要引用计数，记录有多少目录项指向该文件。
  \item 引用计数为 0 时，文件数据和控制信息才能真正回收。
\end{points}
\plain{无环图目录是在树形目录基础上加入“共享同一个文件/目录”的能力，但又必须防止形成环导致路径遍历死循环。}
\exam{容易考“为什么要引用计数”“为什么不允许形成环”，不要答成并发互斥题。}


\concept{通用图目录}
\begin{points}
  \item 若在目录树中不断增加链接，可能形成一般图结构。
  \item 通用图目录允许更灵活的共享，但会带来循环搜索、删除回收更复杂等问题。
  \item 为避免无限遍历，系统通常需要限制搜索深度或在遍历时记录已访问结点。
\end{points}
\plain{图结构比无环图更自由，但自由的代价是目录遍历、删除和垃圾回收都更麻烦。}
\exam{PPT 单独点出它的风险是“循环搜索”和“删除复杂”，这是判断题常见考点。}


\section{文件系统安装、共享与保护}
\concept{文件系统安装 mount}
\begin{points}
  \item 文件系统必须挂载到某个目录位置后，用户才能通过统一目录树访问它。
  \item 挂载点是原有目录树中的一个目录。
  \item 挂载后，该目录下呈现被挂载文件系统的根内容。
\end{points}
\plain{挂载的本质是把“另一个卷上的文件系统根目录”接到当前目录树的某个节点下面。}
\exam{若题目给出 `mount` 场景，关键点是“统一命名空间”和“挂载点目录”。}


\concept{多用户共享中的所有者与组}
\begin{points}
  \item 多用户系统通常用所有者 owner 和组 group 来组织文件共享与保护。
  \item 所有者通常拥有修改文件属性和授予访问权限的最高控制权。
  \item 组表示一组用户，系统可按组统一授予读、写、执行等权限。
  \item 文件或目录会保存用户标识 UserID 和组标识 GroupID，访问时据此判断权限。
\end{points}
\plain{这套机制的核心是“先判你是不是文件主人，再判你是不是这个组的人”，不同身份走不同权限规则。}
\exam{PPT 在文件共享里专门先讲 owner/group，选择题常考 UserID、GroupID 和组权限分别起什么作用。}


\concept{硬链接与符号链接}
\begin{points}
  \item 硬链接是多个目录项指向同一个 inode/文件控制块，本质上是同一个文件有多个名字。
  \item 删除一个硬链接只减少引用计数；只有引用计数为 0 时文件才真正删除。
  \item 符号链接是一个特殊文件，内容是目标文件路径。
  \item 符号链接可以跨文件系统，也可以指向目录；但目标删除后可能变成悬空链接。
\end{points}
\plain{硬链接共享同一个 inode，符号链接是另一个文件，内容是目标路径。打开文件表保存运行时读写状态。}
\exam{常考硬/软链接区别、引用计数、系统级/进程级打开文件表。}


\concept{共享一致性语义}
\begin{points}
  \item 一致性语义描述多个用户同时访问共享文件时，一个用户的修改何时对其他用户可见。
  \item Unix 语义通常认为写操作对其他已打开同一文件的用户立即可见。
  \item AFS 会话语义通常在文件关闭后，修改才对之后新打开该文件的会话可见。
  \item 不可变共享文件语义则要求共享文件一旦发布后内容不可再修改。
\end{points}
\plain{共享不只是“能不能一起用”，还包括“我改完之后别人什么时候看到我的改动”。}
\exam{这类题容易被忽视，但 PPT 单独列了共享一致性语义，至少要会区分 Unix 立即可见和 AFS 会话可见。}


\concept{远程文件系统与故障模式}
\begin{points}
  \item 远程文件访问常见有三类方式：手工程序传输方式（如 FTP）、自动化无缝访问方式（如 NFS、SMB、HDFS 一类 DFS）、以及基于 Web 的半自动化访问方式。
  \item DFS 的特点是可把远端存储挂载到本地目录树中，程序可像访问本地文件那样透明访问远程文件。
  \item 纯 Web/HTTP 方式通常更适合获取和展示文件，不天然提供透明随机写；但配合 WebDAV 等扩展后也可接近挂载访问。
  \item 与本地文件系统相比，远程文件系统除硬件故障、元数据损坏、数据损坏外，还会增加网络故障和服务器故障。
  \item 故障恢复时常涉及请求状态管理，因此会出现无状态 DFS 与有状态 DFS 的区分。
\end{points}
\plain{远程文件系统把“磁盘问题”扩展成了“磁盘 + 网络 + 服务器”三类问题，所以故障模式更复杂。}
\exam{老师录音没展开，但 PPT 有这一块；若出选择题，常考远程文件系统新增的故障来源。}


\concept{文件锁}
\begin{points}
  \item 文件锁允许进程锁定文件，以防其他进程在不合适的时机访问它。
  \item 共享锁类似读锁，允许多个进程并发读文件。
  \item 独占锁类似写锁，同一时刻只允许一个进程持有。
  \item 强制锁由操作系统直接阻止冲突访问；建议锁只提供机制，要求程序自行遵守。
\end{points}
\plain{文件锁本质上就是把读者写者问题搬到了文件访问上。}
\exam{常考共享锁和独占锁的区别，或问强制锁/建议锁到底由谁保证生效。}


\concept{文件保护}
\begin{points}
  \item 文件保护用于控制用户对文件的读、写、执行、删除等权限。
  \item 文件保护既要防止系统崩溃带来的破坏，也要防止用户有意或无意的非法操作。
  \item 常见方法包括访问控制表 ACL、权限位、口令保护、能力表等。
  \item Unix 风格权限常分为所有者、同组用户、其他用户三类，每类有读写执行位。
\end{points}
\plain{文件保护一部分是“防人乱动”，另一部分是“防系统崩了把文件弄坏”；前者偏权限控制，后者偏备份容错。}
\exam{PPT 把这两类问题分开讲了，简答题不要只写访问控制，还可补“定时转储、多副本”等保护思路。}


\concept{文件保护的两类问题}
\begin{points}
  \item 第一类问题是防止系统崩溃或介质故障导致文件被破坏。
  \item 针对这类问题，常见办法有定时转储、建立多副本、备份关键目录和分配表等。
  \item 第二类问题是防止文件所有者或其他用户有意、无意地做出非法操作。
  \item 针对这类问题，核心机制是访问控制，即按“用户 - 对象 - 权限”决定能否访问。
\end{points}
\plain{文件保护题别只盯着权限位；PPT 的原始分法是先问“防系统坏”还是“防人乱动”，然后再分别给办法。}
\exam{简答题若问“文件保护包括哪些内容”，最好先按这两大类分开答，再补 ACL、权限位、备份、多副本等具体手段。}


\concept{打开文件表}
\begin{points}
  \item 打开文件后，系统把文件控制信息缓存在内存打开文件表中，避免每次读写都重新搜索目录。
  \item 系统级打开文件表记录文件位置、引用计数、磁盘位置等全局信息。
  \item 进程级打开文件表记录进程自己的文件描述符、访问模式、当前读写指针等。
  \item 若父子进程共享同一个打开文件表项，可能共享文件偏移；若分别打开同一文件，则偏移独立。
\end{points}
\plain{打开文件表可以理解成“把磁盘上的静态目录信息转成内存里的运行时状态”，这样后续读写才高效。}
\exam{PPT 明确讲了“局部表 + 全局表”，要会区分哪些信息按进程独有，哪些信息全系统共享。}

\chapter{文件系统实现}

\section{文件系统结构}
\concept{用户视图与系统视图}
\begin{points}
  \item 用户视图中，文件是带名字的逻辑数据集合，通过系统调用访问。
  \item 系统视图中，文件是磁盘数据块的集合，文件系统负责从逻辑文件偏移映射到物理磁盘块。
  \item 文件系统基本操作单位是数据块 block，而磁盘物理基本单位是扇区 sector。
  \item 操作系统通常一次读写一个块，块可由多个连续扇区组成。
\end{points}
\plain{同一个 `getc/putc` 看起来只动了 1 个字节，但系统层面往往还是先把整个数据块调入缓存，再在块内读写那个字节。}
\exam{判断题常考“文件系统的基本操作单位是块，不是字节也不是扇区”，这个层次关系要分清。}


\concept{分层结构}
\begin{points}
  \item 逻辑文件系统：处理文件名、目录、权限、FCB/inode 等。
  \item 文件组织模块：完成逻辑块到物理块的映射。
  \item 基本文件系统：向设备驱动发出读写块请求。
  \item I/O 控制层：设备驱动和中断处理程序控制硬件。
\end{points}
\plain{分层结构就是把“文件名解析、块映射、具体设备访问”拆开，各层各管一段。}
\exam{除职责分工外，PPT 还强调了代价：分层能降低复杂度和冗余，但会引入额外开销、降低性能。}


\concept{FCB 与 inode}
\begin{points}
  \item FCB 文件控制块保存文件属性和存储位置信息。
  \item inode 是类 Unix 系统中的索引节点，保存文件类型、权限、大小、时间、数据块地址等元数据。
  \item 目录项通常保存文件名和 inode 号，文件名不一定存放在 inode 内。
  \item inode 与文件物理分配方式有关，因为 inode 中需要保存指向数据块或索引块的指针。
\end{points}
\plain{FCB/inode 保存文件属性和物理位置，是文件系统找到文件内容的核心索引。}
\exam{常考目录项是否要放全部 FCB、inode 与文件名分离的好处。}

\concept{inode 号与按名查找过程}
\begin{points}
  \item 每个 inode 都有唯一的 \term{inode 号}，操作系统内部常用 inode 号而不是文件名来唯一标识文件。
  \item 目录项中通常保存“文件名 + inode 号”，因此目录的首要任务是把名字翻译成 inode 号。
  \item 按名查找文件时，系统先在目录中找到目标文件名对应的 inode 号，再据此读取 inode 信息，最后找到文件数据块。
  \item 因而文件名主要服务于用户使用，inode 号才是内核继续定位文件内容时真正依赖的内部标识。
\end{points}
\plain{可以把文件名看成给人看的“别名”，把 inode 号看成系统内部真正拿来追踪文件本体的“身份证号”。}
\exam{选择题常考“目录项里存的是什么”“系统内部按什么识别文件”“按名查找文件分几步完成”。}


\concept{盘上结构}
\begin{points}
  \item 引导控制块保存从该卷引导操作系统所需的信息。
  \item 卷控制块保存块总数、块大小、空闲块信息、空闲 FCB 或 inode 信息等。
  \item 目录结构负责组织文件名到 FCB/inode 的映射。
  \item FCB 或 inode 区保存每个文件的元数据和数据块地址。
\end{points}
\plain{盘上结构就是文件系统写在磁盘上的“管理骨架”，系统得先找到它们，后面才能找到文件。}
\exam{常考引导块、超级块、目录、inode 区分别干什么。}


\concept{内存结构}
\begin{points}
  \item 安装表 mount table 记录已挂载文件系统、挂载点和文件系统类型。
  \item 目录缓存保存最近访问过的目录信息，用于加速路径解析。
  \item 系统级打开文件表保存已打开文件的 FCB 副本、打开计数和磁盘位置等共享信息。
  \item 进程级打开文件表保存文件描述符、访问权限、当前读写位置等进程私有信息。
\end{points}
\plain{盘上结构负责“长期保存”，内存结构负责“当前加速访问”；其中 buffer 更偏传输缓冲，cache 更偏复用热点数据和元数据。}
\exam{PPT 明确区分了 buffer 和 cache，选择题里不要把两者完全当同义词。}


\concept{open 与 close 的实现}
\begin{points}
  \item open 先查系统打开文件表，若文件已被别的进程打开，则增加打开计数并在本进程表中建立指针。
  \item 若系统表中没有该文件，则搜索目录，把对应 FCB/inode 装入系统打开文件表，再返回文件句柄。
  \item close 删除进程级打开文件表项，并将系统级打开计数减一。
  \item 当打开计数降为 0 时，才把需要写回的元数据更新到磁盘并释放系统表项。
\end{points}
\plain{open 的真正价值是把后续读写反复要用的控制信息先搬进内存，不必每次都重新搜目录。}
\exam{常考 open 为什么能减少访盘，以及多个进程同时打开同一文件时哪些信息共享、哪些不共享。}


\section{目录实现}
\concept{线性列表目录}
\begin{points}
  \item 目录项按线性表保存，查找时顺序扫描。
  \item 优点：实现简单。
  \item 缺点：目录项多时查找慢，平均查找约需扫描一半目录项。
  \item 目录项访问次数题常用“目录项数/每块目录项数/平均扫描块数”计算。
\end{points}
\plain{线性目录查找慢的根因是要顺着比对名字，目录越大，平均扫描项数越多。}
\exam{常考线性目录为什么实现简单但查找代价高，以及平均查找长度估算。}


\concept{哈希目录}
\begin{points}
  \item 根据文件名计算哈希值，快速定位目录项。
  \item 优点：查找速度快。
  \item 缺点：需要处理哈希冲突，目录扩展和哈希表维护更复杂。
\end{points}
\plain{哈希目录的价值在于按名字快速定位，但冲突处理做不好，性能和管理复杂度都会出问题。}
\exam{常考哈希目录相对线性目录的优缺点，尤其冲突问题。}


\concept{inode 对目录检索的加速作用}
\begin{points}
  \item 若目录项中直接保存完整 FCB，则每个目录块能容纳的目录项数量较少。
  \item 把文件名和其他描述信息分离后，目录项只保留“文件名 + inode 号”，一个目录块能装下更多文件名。
  \item 检索时先按文件名找到 inode 号，再读取对应 inode。
  \item 因此按名查找目录时，平均访盘次数通常会下降。
\end{points}
\plain{inode 不是为了让结构更复杂，而是为了让“按文件名查找”这件事更省盘块访问。}
\exam{计算题常用平均访盘次数公式：平均扫描目录块数再加是否还要额外访问 inode 的次数，过程要写清。}


\section{文件存储空间分配方法}
\concept{文件分配的目标}
\begin{points}
  \item 文件存储空间分配首先要回答两个问题：如何为文件分配磁盘块，以及如何在高效利用磁盘空间的同时保证文件访问速度。
  \item 理想目标通常包括：提高磁盘空间利用率、支持文件动态增长、便于顺序访问与随机访问、减少外部碎片与管理开销。
  \item 不同分配方法本质上是在“访问性能、空间利用率、扩展方便性、实现复杂度”之间做权衡。
\end{points}
\plain{这一节不是一上来就背连续/链式/索引三种方法，而是先看它们在解决什么矛盾：既想省空间，又想访问快，还想文件好扩展。}
\exam{简答题若问“文件分配方法设计时要考虑什么”，可以按“空间利用率、访问效率、动态增长、管理开销”四点组织答案。}

\concept{连续分配}
\begin{points}
  \item 文件占用磁盘上一组连续块，目录项只需记录起始块号和长度。
  \item 优点：顺序访问速度快，随机访问也容易，寻道少。
  \item 缺点：需要预先知道文件大小，文件增长困难，会产生外部碎片。
  \item 适合大小较稳定、顺序访问多的文件。
\end{points}
\plain{连续分配答题抓三点：定位快、顺序访问好、增长困难且有外部碎片。}
\exam{常考连续分配的优缺点，以及为什么它不利于文件动态增长。}


\concept{链式分配}
\begin{points}
  \item 文件由多个磁盘块组成链表，每个块包含指向下一个块的指针。
  \item 优点：无外部碎片，文件容易增长。
  \item 缺点：随机访问差，访问第 $k$ 个块必须从第一个块顺链找；指针损坏会影响后续数据。
  \item FAT 是链式分配的改进，把链接指针集中存放在文件分配表中。
\end{points}
\plain{每个块指向下一个块，不怕文件增长，但找第 n 块必须顺着链走。}
\exam{PPT 还给了 FAT 表空间计算题：先算总块数，再算块号至少要多少位，最后乘每表项大小。}


\concept{FAT 表空间计算}
\begin{points}
  \item FAT 中通常“每个磁盘块/簇对应一个表项”，所以先用磁盘容量除以块大小，求出总块数。
  \item 再根据总块数判断“块号至少需要多少位”才能编号全部块。
  \item 为方便存取，实际表项长度常向上取到便于组织的字节数，如 12 位、16 位、20 位等。
  \item FAT 总空间 $=$ 表项大小 $\times$ 总块数。
\end{points}
\plain{FAT 计算题的套路很固定：先算有多少块，再算编号这些块至少要多少位，最后乘出整张 FAT 的大小。}
\exam{PPT 直接举了 `1.2MB/1KB` 和 `200MB/1KB` 这类例子，做题时不要跳过“块号要几位、表项实际取几位”这一步。}


\concept{索引分配}
\begin{points}
  \item 为每个文件建立索引块，索引块中保存该文件所有数据块号。
  \item 目录项保存索引块地址。
  \item 优点：支持直接访问，无外部碎片，文件可离散存放。
  \item 缺点：索引块本身占空间；小文件可能浪费索引块，大文件需要多级索引。
\end{points}
\plain{把文件所有数据块号集中放在索引块里，既支持离散分配，又方便随机访问。}
\exam{常考单级/多级索引最大文件长度计算：每个索引块能放多少块号，再乘数据块大小。}


\concept{多级索引与最大文件长度}
\begin{points}
  \item 若块大小为 $B$ 字节，每个块号占 $p$ 字节，则一个索引块可保存 $B/p$ 个块号。
  \item 一级索引最大可索引 $(B/p)$ 个数据块，最大文件大小为 $(B/p)\times B$。
  \item 二级索引最大可索引 $(B/p)^2$ 个数据块，最大文件大小为 $(B/p)^2\times B$。
  \item 三级索引同理为 $(B/p)^3\times B$。
\end{points}
\plain{这类题的套路固定：先算一个索引块能装多少块号，再根据几级索引逐层乘开，最后再乘数据块大小。}
\exam{不要漏掉“直接地址项”和“间接地址项”混合存在的情况，Unix/Linux 常考的是混合索引而不是纯多级索引。}


\concept{混合索引分配}
\begin{points}
  \item 混合索引把直接地址、一次间接、二次间接、三次间接结合起来，兼顾小文件效率与大文件容量。
  \item 典型 Unix inode 中常见 12 个直接地址项、1 个一次间接、1 个二次间接、1 个三次间接地址项。
  \item 小文件可直接用直接地址快速定位，大文件再逐级借助间接索引扩展容量。
  \item 计算最大文件长度时，要把各部分寻址范围分别算出后相加。
\end{points}
\plain{混合索引的设计思想是：多数文件其实不大，没必要一上来就走多级索引；但系统又必须给超大文件留扩展空间。}
\exam{考研题很爱出这种“若有若干直接地址、若干一级间接、一个二级间接”的最大文件长度计算。}

\section{空闲空间管理}
\concept{位示图/位图}
\begin{points}
  \item 用一个二进制位表示一个磁盘块是否空闲。
  \item 优点：查找连续空闲块方便，结构紧凑。
  \item 缺点：磁盘很大时位图本身也较大，需要常驻或缓存。
\end{points}
\plain{位图法的关键是“按位记录每个块是否空闲”，所以做题常要把字序号、位偏移和实际块号互相换算。}
\exam{PPT 给了大磁盘下位图本身占多大空间的估算题：先算总块数，再除以 8 得到字节数。}


\concept{空闲链表}
\begin{points}
  \item 把所有空闲块链接成链表。
  \item 优点：释放和分配单个块较简单。
  \item 缺点：查找连续空闲块效率低，链指针损坏可能影响管理。
\end{points}
\plain{链接法的主要短板不是“不能管空闲块”，而是“很难快速拿到一大片连续空闲块”。}
\exam{若题目强调连续空间分配需求，通常位图法或计数法更合适，单纯空闲链表往往不优。}


\concept{成组链接法}
\begin{points}
  \item 把空闲块分组，每组第一块记录下一组空闲块号。
  \item 兼顾链表和批量管理，是 Unix 早期常用方法。
  \item 适合大量空闲块的快速分配和回收。
\end{points}
\plain{成组链接法不是逐块找空闲块，而是把很多空闲块号成组保存，减少管理链过长的问题。}
\exam{PPT 对 Unix 成组链接法的“分配最后一个块号时先装入下一组”“表满回收时先整组写回”讲得很细，这两步是流程题重点。}


\concept{Unix 成组链接法的分配与回收}
\begin{points}
  \item 分配空闲块时，先从内存中当前专用块的空闲块号表取出下一个可用块号。
  \item 若取出的正好是当前组记录的“最后一个块号”，则要先把该块中保存的下一组空闲块号表读入内存，再把这个块分配出去。
  \item 回收空闲块时，若当前专用块中的空闲块号表未满，就直接把回收块号填入表中。
  \item 若空闲块号表已满，则先把当前整张空闲块号表写入这个新回收的块中，再把该回收块号作为新组首块号放回表中重新开始计数。
\end{points}
\plain{成组链接法流程题的关键是分清两种“临界时刻”：分配到当前组最后一个块号时要先装入下一组；回收时若表已满，要先把整组信息塞进新回收块。}
\exam{这类题通常不难算，难在顺序别写反；一旦把“先读下一组再分配”或“先写满表再重置新组”写错，后面全错。}


\concept{计数法}
\begin{points}
  \item 用“起始空闲块号 + 连续空闲块数”表示一段连续空闲区。
  \item 适合磁盘空闲块成片连续的情况。
  \item 比逐块链表更紧凑，但碎片多时记录数量增加。
\end{points}
\plain{计数法适合连续空闲区较多的磁盘，用一条记录概括一整段连续空闲块。}
\exam{判断题会考它为什么适合“频繁、连续的分配与释放”，核心原因是连续空闲段能被压缩成一条记录。}


\section{文件共享实现}
\concept{文件共享的动机}
\begin{points}
  \item 文件共享允许不同用户共同使用同一个文件。
  \item 这样做的主要动机有三类：支持用户合作、减少磁盘空间开销、减少同一数据维护多份副本带来的不一致性。
  \item 因而文件共享不仅是“方便别人读我的文件”，更是从协作效率、存储成本和一致性控制三方面做优化。
\end{points}
\plain{共享的价值不只是省空间；更关键的是多人协作时大家面对同一份数据，能少掉很多“各改各的副本然后版本打架”的问题。}
\exam{若选择题问“为什么要实现文件共享”，标准关键词就是合作、节省空间、减少不一致。}


\concept{基于索引节点的共享}
\begin{points}
  \item 多个目录项指向同一个 inode，实现硬链接共享。
  \item inode 中维护引用计数，记录有多少目录项引用它。
  \item 删除目录项只减少引用计数，不一定删除文件内容。
\end{points}
\plain{硬链接共享的是同一个 inode，所以任何一个用户修改文件内容，其他硬链接看到的都是同一份数据。}
\exam{还要会写出结论：创建硬链接时引用计数加一；只有所有硬链接都删掉后，文件才真正删除。}


\concept{索引节点中的链接计数}
\begin{points}
  \item inode 中的链接计数字段 \code{count} 表示当前有多少个目录项指向该 inode。
  \item 创建硬链接时，\code{count} 加 1；删除一个硬链接时，\code{count} 减 1。
  \item 只要 \code{count} 仍大于 0，文件数据通常不会真正删除，因为仍有其他文件名可访问它。
  \item 因而“删掉某个文件名”与“文件内容真正消失”在硬链接语义下不是同一件事。
\end{points}
\plain{硬链接题的关键是把“名字”和“文件本体”分开：目录项删掉的首先只是一个名字，只有最后一个名字也删掉时，文件才真正失去引用。}
\exam{PPT 单独把 \code{count} 拿出来讲，选择题和简答题都常考：什么时候加一、减一，以及文件真正删除的条件。}


\concept{基于符号链接的共享}
\begin{points}
  \item 符号链接文件保存目标路径名。
  \item 访问符号链接时，系统根据路径再去查找目标文件。
  \item 优点是灵活、可跨文件系统；缺点是访问多一次路径解析，目标删除后链接失效。
\end{points}
\plain{符号链接并不直接共享 inode，它只是保存“去哪里找目标文件”的路径，因此本身更像快捷方式。}
\exam{常考硬链接与软链接对引用计数的影响：硬链接会增加原文件引用计数，软链接通常不会。}


\chapter{大容量存储系统}

\section{磁盘结构与访问时间}
\concept{磁盘结构}
\begin{points}
  \item 磁盘由一个或多个盘片组成，盘片表面覆盖磁性材料。
  \item 每个盘面对应一个磁头，磁臂沿半径方向移动。
  \item 磁道是盘面上的同心圆，扇区是磁道上的固定大小区域。
  \item 柱面是所有盘面上半径相同的磁道集合。
  \item 一个物理块可用柱面号、磁头号、扇区号表示。
\end{points}
\plain{这部分先建立“盘片-盘面-磁头-磁道-扇区-柱面”的空间结构图，后面访问时间、调度算法、坏块管理都建立在这个结构上。}
\exam{概念题常考柱面、磁道、扇区的区别；计算题则常默认你已经能据此理解磁头为什么要移动。}


\concept{磁盘访问时间}
\begin{points}
  \item 寻道时间：磁臂移动到目标柱面/磁道所需时间。
  \item 旋转延迟：目标扇区旋转到磁头下方所需时间，平均为旋转一周时间的一半。
  \item 传输时间：数据从磁盘读出或写入磁盘的时间。
  \item 总访问时间约为：寻道时间 + 旋转延迟 + 传输时间。
  \item 机械磁盘优化重点通常是减少寻道时间。
\end{points}
\plain{磁盘访问时间要分解着看，机械盘里通常最贵的是寻道和旋转等待，不是数据真正传过去那一下。}
\exam{PPT 还给了完整 I/O 时间计算：定位时间 + 传输时间 + 控制器负载时间，做题别漏最后一项。}


\concept{逻辑块映射与磁盘连接方式}
\begin{points}
  \item 文件系统常把磁盘看成逻辑块的一维数组，逻辑块按一定规则顺序映射到实际扇区。
  \item 逻辑到物理的地址转换并不总是简单线性，因为可能存在坏扇区、不同磁道扇区数不完全一致等情况。
  \item 存储连接方式常见有 DAS、NAS、SAN：分别对应主机直接连接、经网络文件访问、经专用存储网络访问。
  \item 总线、主机控制器、磁盘控制器共同参与主机与磁盘的通信。
\end{points}
\plain{题目写“逻辑块号”时，不等于它天然就是磁盘上的物理位置编号；中间往往还隔着控制器和映射规则。}
\exam{PPT 单独讲了 DAS/NAS/SAN，选择题里常考三者“直接接主机 / 通过网络文件协议 / 通过专用存储网络”的差别。}


\concept{磁盘、磁带与 SSD}
\begin{points}
  \item 磁带是顺序访问设备，容量大、成本低，适合备份和归档，但随机访问极慢。
  \item 磁盘是直接访问设备，支持随机定位块，是典型通用二级存储。
  \item SSD 基于闪存，读取快、抗震、噪音低，但价格较高且擦写寿命有限。
\end{points}
\plain{这三类设备的关键区别是访问方式和适用场景，不是背一堆参数。}
\exam{常考顺序存取设备和直接存取设备代表，或比较 HDD、SSD、磁带的优缺点。}


\section{磁盘调度算法}
\concept{FCFS}
\begin{points}
  \item 按磁盘请求到达顺序服务。
  \item 优点：公平、简单。
  \item 缺点：磁头来回移动可能很大，平均寻道距离长。
\end{points}
\plain{FCFS 在磁盘里就是不做任何优化地按到达顺序走，所以磁头可能来回折返很多次。}
\exam{常考按给定请求序列计算 FCFS 的总移动磁道数。}


\concept{SSTF 最短寻道时间优先}
\begin{points}
  \item 每次选择距离当前磁头位置最近的请求。
  \item 优点：平均寻道距离通常较短。
  \item 缺点：远处请求可能长期等待，产生饥饿。
\end{points}
\plain{SSTF 每次都挑最近的请求，平均移动少，但远端请求可能一直被近处新请求压着，出现饥饿。}
\exam{常考按请求序列逐步选择“离当前磁头最近”的请求，并分析饥饿问题。}


\concept{SCAN 电梯算法}
\begin{points}
  \item 磁头像电梯一样沿一个方向移动，服务沿途请求，到达端点后反向。
  \item 优点：比 FCFS 稳定，避免频繁来回。
  \item 缺点：靠近两端的请求等待时间可能不同。
\end{points}
\plain{SCAN 的关键不是“最近优先”，而是“沿一个方向扫过去再回头”，这样能避免磁头反复折返。}
\exam{有些教材把只扫到最远请求就掉头的版本也直接叫 SCAN/C-SCAN；考试前要按老师 PPT 的定义作答。}


\concept{C-SCAN 循环扫描}
\begin{points}
  \item 磁头只沿一个方向服务，到达一端后快速返回另一端，不在返回途中服务。
  \item 提供更均匀的等待时间。
  \item 适合请求分布较均匀的系统。
\end{points}
\plain{C-SCAN 像循环单向电梯：服务时始终朝一个方向走，所以两端请求受到的对待更平均。}
\exam{算总移动距离时，别忘了它从一端“跳回”另一端那段距离通常也要计入。}


\concept{LOOK 与 C-LOOK}
\begin{points}
  \item LOOK 不一定走到磁盘端点，只走到当前方向最远请求处就反向。
  \item C-LOOK 只在有请求的范围内循环扫描，从一端最远请求跳到另一端最远请求。
  \item 它们是 SCAN/C-SCAN 的实际优化版本，减少无意义移动。
\end{points}
\plain{LOOK/C-LOOK 的改进点就是“不走冤枉路”：没有请求的尾部磁道就不必特意扫到头。}
\exam{老师 PPT 明说它们是 SCAN/C-SCAN 的改良版，若题目要求综合性能，常优先想到 LOOK 类。}


\concept{N-Step-SCAN}
\begin{points}
  \item 把磁盘请求队列按长度 $N$ 划分成若干批次。
  \item 批次之间按 FCFS 处理，每个批次内部再按 SCAN 调度。
  \item 它可以减轻普通 SCAN 的磁臂粘着问题，避免新到请求不断打断当前扫描。
  \item 当 $N$ 很大时接近普通 SCAN；当 $N=1$ 时退化为 FCFS。
\end{points}
\plain{核心思想是“先把请求分批，再对每批做电梯调度”，这样新请求不会一直把磁头吸住。}
\exam{若题目给出 $N$ 和新到请求，要先切批，再分别算每一批内部的扫描路径。}


\concept{FSCAN}
\begin{points}
  \item FSCAN 是 N-Step-SCAN 的简化版，只维护两个队列。
  \item 当前扫描只处理旧队列；扫描期间新到达的请求进入新队列。
  \item 当前队列处理完后，再把新队列作为下一轮扫描对象。
  \item 这样可避免新请求持续涌入导致当前扫描迟迟结束不了。
\end{points}
\plain{可以把 FSCAN 看成“双缓冲版 SCAN”：一边扫当前队列，一边收集下一轮请求。}
\exam{常和普通 SCAN 对比考，问它为什么更能避免磁臂粘着。}


\concept{如何选择磁盘调度算法}
\begin{points}
  \item 磁盘调度算法效果受请求数量、请求分布、文件分配方式等因素共同影响。
  \item 不存在对所有场景都绝对最优的算法，因此系统常保留可替换的调度策略。
  \item SSTF 和 LOOK 往往是较常见的默认选择。
  \item SCAN 与 C-SCAN 在高负载场景下更有利于避免饥饿、维持较稳定性能。
\end{points}
\plain{这类题不要死背“哪一个最好”，真正结论是没有通杀算法，只能按负载特征取舍。}
\exam{简答题若问“为什么没有最优磁盘调度算法”，核心要点是请求分布和系统负载是动态变化的。}


\section{磁盘存储优化分布}
\concept{记录排列会影响 I/O 总时间}
\begin{points}
  \item 在旋转型磁盘上，记录在磁道上的排列顺序会影响旋转等待时间。
  \item 如果逻辑上相邻的记录物理上也紧挨着，但处理上一条记录需要时间，则下一条记录可能已经转过磁头。
  \item 因此记录的物理分布需要结合“读出时间 + 处理时间”做错位安排。
  \item 这一优化减少的是旋转延迟，不是寻道时间。
\end{points}
\plain{上一条记录刚处理完时，下一条记录最好正好转到磁头下，而不是刚错过去。}
\exam{常考比较普通顺序分布和优化分布下的总处理时间。}


\concept{优化分布的直觉}
\begin{points}
  \item 若 CPU 处理当前记录要花若干毫秒，就应让下一条逻辑记录在若干扇区之后出现。
  \item 这样 CPU 处理结束时，下一条逻辑记录刚好到达磁头位置。
  \item 题目若给出“一圈时间、每条记录读出时间、处理时间、每道记录数”，通常按转过多少个扇区后恰好接上下一次读取来排布。
  \item SSD 不依赖机械旋转，因此不存在这种扇区优化分布问题。
\end{points}
\plain{它像排队卡节奏：不是让人全挤在窗口前，而是每次轮到时下一个人刚好走到窗口。}
\exam{常考计算题或简答：为什么优化排布能把总时间从“多等整圈”降到“几乎无额外等待”。}


\section{磁盘管理与容错}
\concept{磁盘格式化}
\begin{points}
  \item 低级格式化把磁盘划分为扇区，为每个扇区建立控制信息。
  \item 分区把磁盘划分为一个或多个逻辑区域。
  \item 高级格式化/逻辑格式化在分区上建立文件系统数据结构，如空闲空间表、根目录、inode 区等。
\end{points}
\plain{这三步是从“原始盘面”到“能存文件的文件系统分区”的构建过程：先切扇区，再切分区，再建文件系统。}
\exam{常考低级格式化、逻辑分区、高级格式化三者分别解决什么问题，别混成一步。}


\concept{扇区、块与簇}
\begin{points}
  \item 扇区 sector 是磁盘的物理存储单位。
  \item 块 block 是操作系统进行磁盘 I/O 时常用的逻辑传输单位，通常由多个扇区组成。
  \item 簇 cluster 常作为文件系统分配文件空间的单位，也可能由多个扇区或多个块组成。
  \item 做题时要先分清题目说的是物理扇区、I/O 块还是文件分配簇。
\end{points}
\plain{扇区是硬盘原生小格子，块和簇是操作系统为传输和分配方便再打的包。}
\exam{常考“磁盘 I/O 以什么为单位”“文件系统分配空间以什么为单位”这类辨析。}


\concept{原始磁盘 raw disk}
\begin{points}
  \item 原始磁盘或原始分区指不建立普通文件系统、直接由应用管理块的磁盘区域。
  \item 它减少了文件系统带来的额外管理开销，适合数据库等追求高性能的场景。
  \item 代价是通用性差，很多管理责任要交给应用自己处理。
\end{points}
\plain{raw disk 就是“别走文件系统这层，我直接控制这块盘”。}
\exam{常考 raw disk 为什么可能更快，以及它牺牲了哪些便利性。}


\concept{引导块与坏块}
\begin{points}
  \item 引导块保存启动操作系统所需的引导代码或引导信息。
  \item 坏块是无法可靠读写的磁盘块。
  \item 坏块管理可通过扇区备用、坏块链表、控制器重映射等方式处理。
\end{points}
\plain{磁盘管理不仅是读写，还包括初始化、启动、坏块处理和容错。}
\exam{PPT 还给了引导流程示例：ROM/BIOS -> MBR -> 分区引导扇区 -> 引导管理器 -> 内核，知道层次即可。}


\concept{坏块管理与重映射}
\begin{points}
  \item 发现坏块后，系统可把它记录到坏块表或坏块链表中，后续分配时直接跳过。
  \item 现代磁盘还常预留备用块或备用扇区，用于对坏块做透明重映射。
  \item 对上层来说，目标是屏蔽物理缺陷，尽量不让文件系统继续碰到坏块。
\end{points}
\plain{坏块一般不是去“修”，而是尽快标记并绕开。}
\exam{常考坏块链表、备用块、控制器重映射各自起什么作用。}


\concept{交换空间管理}
\begin{points}
  \item 交换空间用于保存被换出的进程映像或页面。
  \item 交换空间可作为专门分区，也可作为文件系统中的交换文件。
  \item 专门分区管理简单且效率高；交换文件更灵活但可能受文件系统开销影响。
\end{points}
\plain{交换空间不是普通存文件的地方，而是虚拟内存把页面或进程映像临时换出去的后备区域。}
\exam{常考交换空间的用途，以及它和普通文件存储区域的区别。}


\concept{交换文件与交换分区比较}
\begin{points}
  \item 交换文件实现简单，配置灵活，能直接建在现有文件系统里。
  \item 但访问交换文件通常还要经过文件系统的相关结构，可能带来额外开销。
  \item 专用交换分区不依赖普通文件系统，访问路径更直接，更便于专门优化。
  \item 因为交换访问很频繁，系统常更偏向专用交换分区。
\end{points}
\plain{交换区是虚拟内存的后备仓库，用得勤，所以越少绕文件系统这层越好。}
\exam{常考为什么独立交换分区通常比交换文件效率更高。}


\concept{RAID 磁盘容错}
\begin{points}
  \item RAID 通过磁盘冗余提高性能和可靠性。
  \item RAID 0 条带化提高性能但无冗余。
  \item RAID 1 镜像保存完整副本，可靠性高但空间开销大。
  \item RAID 5 使用分布式奇偶校验，兼顾容量和容错。
  \item RAID 6 可容忍两块磁盘故障，可靠性更高。
\end{points}
\plain{RAID 要分清“提速”和“容错”不是总能同时最大化；不同级别是在容量、性能、可靠性间做取舍。}
\exam{除了 0/1/5/6，PPT 还出现了 RAID 2/3/4：它们分别引入海明码校验、专用奇偶校验盘、块级奇偶校验盘，虽少用但可能出选择题。}


\concept{并行交叉存取与 RAID 2/3/4}
\begin{points}
  \item 并行交叉存取把一个盘块拆成多个子块并分散到多个磁盘上并行读写，以缩短传输时间。
  \item RAID 2 以位/字节级拆分并配合海明码校验，结构复杂，商业上较少使用。
  \item RAID 3 使用一块专用奇偶校验盘，适合大块顺序传输场景，如科学计算、图像处理。
  \item RAID 4 与 RAID 3 类似，但按块级拆分数据，仍使用专用奇偶校验盘。
\end{points}
\plain{RAID 2/3/4 的共同思路都是“条带化 + 冗余校验”，差别主要在数据拆分粒度和校验组织方式。}
\exam{若题目问“哪几级依赖专用校验盘”，要想到 RAID 3/4；若问“分布式奇偶校验”，那是 RAID 5。}


\concept{RAID 5}
\begin{points}
  \item RAID 5 采用块级条带化，并把奇偶校验信息分布到各个磁盘上，而不是集中到一块专用校验盘。
  \item 每个驱动器通常都有独立的数据通道，可独立进行读写，适合 I/O 繁忙的事务处理场景。
  \item 它兼顾了容量利用率、读写性能和容错能力，是工程上较常见的一类 RAID 方案。
  \item 与 RAID 3/4 相比，RAID 5 的关键改进是消除了专用校验盘的单点瓶颈。
\end{points}
\plain{RAID 5 的记忆重点不是“也有校验”，而是“校验分散到所有盘上”，所以既能容错，又不容易把一块校验盘压成热点。}
\exam{选择题和比较题常把 RAID 3/4/5 放一起考：3/4 有专用校验盘，5 是分布式奇偶校验。}


\concept{RAID 6 与 RAID 7}
\begin{points}
  \item RAID 6 在冗余校验方面比 RAID 5 更强，可进一步容忍更多故障，可靠性更高。
  \item 课件表述里，RAID 6 设置了一个可快速访问的异步校验盘，并具有独立数据访问通道，性能优于部分早期奇偶校验方案。
  \item RAID 7 在阵列控制层集成了更智能的实时控制软件，可独立于主机运行，尽量减少主机 CPU 负担。
  \item 记忆时抓住趋势：级别越高，通常越强调控制智能化和更强的故障容忍能力。
\end{points}
\plain{这两级不用背工程细节，考试里更重要的是知道它们是在 RAID 5 之后继续往“更强容错/更强控制能力”方向演化。}
\exam{若选择题专门问 RAID 6/7，抓“RAID 6 更强容错，RAID 7 更智能、对主机依赖更少”通常就够了。}


\concept{磁盘容错层次}
\begin{points}
  \item 低级容错主要针对磁盘表面缺陷，例如双份目录、双份文件分配表、写后读校验、热修复重定向。
  \item 中级容错主要针对驱动器或控制器故障，例如磁盘镜像和磁盘双工。
  \item 更高层容错关注系统级持续服务能力。
\end{points}
\plain{容错不只是一句 RAID，它还分“防单个坏块”“防整块盘坏”“防整套系统失效”几个层次。}
\exam{常考双份目录、写后读校验属于哪一级容错，以及镜像和双工分别冗余了什么。}

\concept{双份目录和双份文件分配表}
\begin{points}
  \item 目录和文件分配表都是文件系统运行所依赖的关键数据结构，一旦损坏会直接影响文件定位与空间管理。
  \item 为提高可靠性，可在不同磁盘或同一磁盘的不同区域保存两份：一份主目录/主文件分配表，一份备份目录/备份文件分配表。
  \item 当主副本损坏时，系统可切换到备份副本继续工作。
  \item 系统启动时通常还要检查两份数据结构的一致性。
\end{points}
\plain{这类方法保护的不是“普通文件内容”，而是文件系统自己的核心账本；账本一坏，整套文件管理都会受影响。}
\exam{简答题若问低级容错的典型措施，写出“双份目录、双份文件分配表、主备切换、一致性检查”会比只写“做备份”更稳。}

\concept{热修复重定向与写后读校验}
\begin{points}
  \item \term{写后读校验}：每次把数据块写入磁盘后，立即再读出并与原写入数据比较，确认写入结果是否正确。
  \item 若校验失败且重写后仍不一致，就可判定该盘块存在缺陷。
  \item \term{热修复重定向}：把待写数据改写到预留的热修复重定向区，并登记映射关系，供以后访问时查找。
  \item 这类方法适合缺陷数量较少时对局部坏块做补救，而不必立刻放弃整块磁盘。
\end{points}
\plain{写后读校验负责“发现写坏了没有”，热修复重定向负责“既然这块不可靠，就把数据挪到备用区并记好新地址”。}
\exam{选择题常考两者分工：前者偏检测，后者偏绕开坏块并继续提供可访问位置。}


\concept{磁盘镜像与磁盘双工}
\begin{points}
  \item \term{磁盘镜像}：两台磁盘驱动器连接到同一个磁盘控制器，主盘写入后再把同样数据写到备份盘，主盘故障时启用备份盘。
  \item 磁盘镜像主要冗余的是\key{磁盘驱动器}本身，但若共享的那一个控制器失效，两块盘仍可能一起不可用。
  \item \term{磁盘双工}：两台磁盘驱动器分别接到两个不同的磁盘控制器，同时保存完全相同的数据副本。
  \item 磁盘双工比镜像更进一步，因为它同时冗余了\key{磁盘驱动器和磁盘控制器}，可靠性更高，但成本也更高。
\end{points}
\plain{镜像和双工最稳的区分方式不是背图，而是看“两个备份盘是不是还共用同一个控制器”：共用的是镜像，不共用的是双工。}
\exam{比较题常直接问两者区别，标准答法就是“镜像冗余盘，双工冗余盘加控制器”，并顺带说明双工可靠性更高。}


\concept{第三级系统容错}
\begin{points}
  \item 第三级系统容错建立在前两级容错基础上，目标是不因单台服务器故障而中断整体服务。
  \item 典型做法是主从服务器镜像：主服务器正常对外服务，从服务器实时保存主服务器的关键内存和磁盘数据副本。
  \item 当主服务器故障时，从服务器接替成为新的服务节点；故障排除后，两台服务器再重新同步。
  \item 这一级关注点已不只是“某块盘坏了怎么办”，而是“整个系统如何持续提供服务”。
\end{points}
\plain{前两级主要还在“保住磁盘和控制器”，第三级已经上升到“即使整台主服务器挂了，业务也别停”。}
\exam{PPT 单独讲了第三级容错，简答题常考它和前两级的层次区别：前者偏系统接管，后者偏磁盘/控制器冗余。}

\chapter{设备管理}

\section{I/O 系统与设备分类}
\concept{设备管理的对象与任务}
\begin{points}
  \item 除 CPU 和内存外，大部分硬件设备称为外部设备。
  \item I/O 系统包括 I/O 设备、接口线路、控制部件、通道和管理软件。
  \item I/O 操作是主存与外围设备介质之间的信息传送。
  \item 设备管理任务：响应 I/O 请求、分配 I/O 设备、管理设备可靠性与安全性、提供友好访问接口。
\end{points}
\plain{这题通常先答“设备管理管什么”，再答“为什么需要它”：设备种类多、速度差大、还常被多个进程共享。}
\exam{常考设备管理的对象、任务和 I/O 系统包含哪些部分。}


\concept{设备管理四个功能}
\begin{points}
  \item 设备分配：根据用户 I/O 请求分配设备、控制器和通道。
  \item 设备处理：启动设备，并在 I/O 完成时处理中断。
  \item 缓冲管理：缓和 CPU 与 I/O 速度不匹配，管理缓冲区分配与释放。
  \item 设备独立性：用户编程使用逻辑设备名，不依赖具体物理设备。
\end{points}
\plain{四个功能最好能一一对上：分配管谁用设备，处理管怎么启动和完成，缓冲管速度差，独立性管逻辑名和物理名解耦。}
\exam{常考设备管理四项基本功能的对应含义。}


\concept{按使用特性分类}
\begin{points}
  \item 存储设备：保存信息，如磁盘、磁带、SSD。
  \item 传输设备：传输数据，如网卡、调制解调器。
  \item 人机交互设备：输入或输出用户信息，如键盘、显示器、打印机。
\end{points}
\plain{按使用特性分类看的是“设备拿来干什么”，不是看它能不能共享或按什么粒度传数据。}
\exam{常考存储设备、传输设备、人机交互设备的代表例子。}


\concept{按共享属性分类}
\begin{points}
  \item 独占设备：一段时间内只允许一个进程使用，如传统打印机。
  \item 共享设备：允许多个进程交替或并发使用，如磁盘。
  \item 虚拟设备：通过虚拟技术把独占设备改造成多个逻辑设备，如 SPOOLing 打印机。
\end{points}
\plain{这组分类要抓住“同一时段能否给多人用”：独占设备一次只给一个，虚拟设备则是把独占设备在逻辑上改成可共享。}
\exam{常考独占、共享、虚拟设备的例子，尤其打印机和磁盘。}


\concept{按信息交换单位分类}
\begin{points}
  \item 块设备以固定大小数据块为单位传输，如磁盘、磁带。
  \item 字符设备以字符流为单位传输，如键盘、鼠标、串口、打印机。
  \item 选择题中，磁盘通常是块设备；键盘、打印机、显示器通常是字符设备。
\end{points}
\plain{块设备和字符设备最容易从“能否随机按块定位”和“传输单位是否固定”来区分。}
\exam{常考磁盘、键盘、打印机分别属于哪类设备。}


\section{I/O 硬件}
\concept{设备控制器与寄存器}
\begin{points}
  \item 设备控制器是 CPU 与外设之间的电子部件，负责控制设备操作。
  \item 控制器通常包含数据寄存器、状态寄存器、控制寄存器。
  \item CPU 通过读写控制器寄存器与设备交互。
  \item 设备完成操作或发生错误时，通过中断通知 CPU。
\end{points}
\plain{设备控制器不是外设本体，而是 CPU 与外设之间的控制桥梁，OS 真正常接触的是这些寄存器。}
\exam{常考控制寄存器、状态寄存器、数据寄存器各自的作用。}


\concept{端口、总线与控制器接口}
\begin{points}
  \item 端口 port 是设备与计算机连接和通信的接口点。
  \item 总线 bus 是多设备共享的一组通信线路及其协议。
  \item 设备控制器内部通常既有与 CPU 的接口，也有与具体设备的接口，以及负责译码和控制的 I/O 逻辑。
\end{points}
\plain{可以把控制器理解成“翻译官 + 调度员”：一边听 CPU 的，一边管具体设备怎么动。}
\exam{常考控制器由哪几部分组成，或端口/总线/控制器分别是什么。}


\concept{I/O 地址空间}
\begin{points}
  \item 端口映射 I/O：I/O 设备有独立地址空间，CPU 使用专门 I/O 指令访问端口。
  \item 内存映射 I/O：设备寄存器映射到内存地址空间，CPU 用普通访存指令访问。
  \item 二者都需要保护，普通用户程序不能随意访问设备寄存器。
\end{points}
\plain{I/O 地址空间这题只看一点：设备寄存器到底是单独编址，还是映射进主存地址空间。}
\exam{常考端口映射 I/O 和内存映射 I/O 的区别。}


\concept{独立编址与内存映像编址比较}
\begin{points}
  \item 独立编址把 I/O 端口地址空间与内存地址空间分开，通常使用专用 I/O 指令。
  \item 内存映像编址把设备寄存器映射到主存地址空间，使用普通访存指令即可访问。
  \item 前者硬件隔离清晰；后者编程更统一，很多体系结构广泛采用。
\end{points}
\plain{一个是“设备地址单开一本账”，一个是“设备也算在内存地址里”。}
\exam{常考两种编址方式的区别，以及哪一种需要专用 IN/OUT 一类指令。}


\concept{通道}
\begin{points}
  \item 通道是一种专门负责 I/O 控制的处理部件，比普通设备控制器更强。
  \item 通道能执行通道程序，组织多个 I/O 操作，减轻 CPU 负担。
  \item 处理机和通道可以并行工作，因此通道提高了大型系统 I/O 并行能力。
\end{points}
\plain{通道比一般控制器更像一个小型 I/O 处理机，能自己按通道程序组织一串 I/O 操作。}
\exam{常考通道和 DMA 的区别，尤其“谁能执行更复杂的控制逻辑”。}


\section{I/O 控制方式}
\concept{程序直接控制方式}
\begin{points}
  \item CPU 反复查询设备状态寄存器，直到设备就绪，再进行数据传送。
  \item 优点：实现简单。
  \item 缺点：CPU 忙等，浪费处理机时间。
  \item 又称轮询方式，适合简单或低速设备。
\end{points}
\plain{程序直接控制就是 CPU 亲自盯着设备状态反复问，简单但最浪费 CPU。}
\exam{常考轮询方式为什么效率低，以及它和中断方式的根本差别。}


\concept{中断驱动方式}
\begin{points}
  \item CPU 启动 I/O 后转去执行其他任务，设备完成后发中断通知 CPU。
  \item 优点：减少 CPU 忙等，提高并发性。
  \item 缺点：每次传送仍需 CPU 参与，频繁中断会增加开销。
\end{points}
\plain{这里要抓住 I/O 主线：CPU 发出 I/O 请求后不用原地等，设备完成时再用中断把 CPU 叫回来收尾。}
\exam{常考中断处理步骤、中断与异常/系统调用区别、为什么中断能提高 CPU 与 I/O 并行度。}


\concept{DMA 方式}
\begin{points}
  \item DMA 允许设备控制器在主存和设备之间直接传送一批数据，无需 CPU 逐字节搬运。
  \item CPU 只负责初始化 DMA 控制器，设置内存地址、传输长度、方向等。
  \item 传输完成后 DMA 控制器发中断通知 CPU。
  \item 优点：大幅减少 CPU 参与，适合磁盘、网卡等高速块设备。
\end{points}
\plain{DMA 的关键不是“没有中断”，而是“搬数据这件事不用 CPU 一字节一字节亲手做”，CPU 只在首尾参与。}
\exam{常考 DMA 和中断方式的差别，尤其“中断发生在每次传送后还是一批传送结束后”。}


\concept{通道控制方式}
\begin{points}
  \item CPU 把一组 I/O 操作写成通道程序交给通道执行。
  \item 通道独立控制设备完成复杂 I/O 序列，完成后中断 CPU。
  \item 与 DMA 相比，通道更像专用 I/O 处理机，控制能力更强。
\end{points}
\plain{通道比 DMA 更进一步：DMA 主要负责一批数据搬运，通道还能按通道程序组织多步 I/O 控制逻辑。}
\exam{比较题常问 DMA 和通道谁更像专用处理机、谁能执行更复杂的 I/O 控制序列。}


\concept{轮询、中断、DMA、通道的核心区别}
\begin{points}
  \item 轮询方式中，CPU 主动反复查状态，忙等最明显。
  \item 中断方式中，每次数据传送仍通常需要 CPU 参与，但不用持续忙等。
  \item DMA 方式中，CPU 只做初始化，一批数据由 DMA 控制器直接在设备和内存间搬运。
  \item 通道方式比 DMA 更进一步，能执行通道程序，控制一组 I/O 操作序列。
\end{points}
\plain{四种方式的主线就是：CPU 参与程度越来越低，I/O 并行能力越来越强。}
\exam{常考比较题：谁在每个字节/每批数据完成后中断 CPU，谁负责真正搬数据。}


\section{I/O 软件原理}
\concept{I/O 软件的设计目标和原则}
\begin{points}
  \item I/O 软件的总体设计目标主要是\key{高效率}和\key{通用性}。
  \item 设计时要重点考虑设备无关性，即上层接口尽量不依赖具体硬件细节。
  \item 还要考虑出错处理，因为 I/O 故障既可能来自设备本身，也可能来自更高层数据结构或访问环境。
  \item 同步传输、异步传输，以及独占设备和共享设备的不同管理方式，也都属于 I/O 软件设计时必须统一处理的问题。
\end{points}
\plain{这块是整章 I/O 软件的总纲：目标不是单纯“能驱动设备”，而是既要高效，又要让不同设备在统一框架下工作。}
\exam{简答题若问 I/O 软件设计原则，标准答法通常先写“高效率、通用性”，再补“设备无关、出错处理、同步/异步、独占/共享设备管理”。}

\concept{I/O 软件层次}
\begin{points}
  \item 用户层 I/O 软件：库函数、假脱机程序等，为用户提供高级接口。
  \item 设备独立性软件：提供统一逻辑设备接口、命名、保护、缓冲、错误处理。
  \item 设备驱动程序：针对具体设备控制器发命令、处理中断。
  \item 中断处理程序：响应设备中断，完成状态检查、唤醒等待进程等。
\end{points}
\plain{I/O 软件分层的主线是“越往上越面向用户和统一接口，越往下越贴近具体硬件和寄存器”。}
\exam{常考四层软件各自负责什么，尤其设备无关软件和驱动程序的边界。}


\concept{I/O 中断处理程序}
\begin{points}
  \item I/O 中断处理程序在设备完成操作、发生错误或出现其他异步事件时被触发。
  \item 主要工作包括检查中断原因、读取状态、确认本次 I/O 是否正常结束，并唤醒等待的驱动程序或进程。
  \item 若发现异常，还要把错误信息上报给更高层软件，交由驱动或设备无关软件进一步处理。
  \item 它位于 I/O 软件的最底层之一，直接承接硬件中断信号。
\end{points}
\plain{驱动程序负责“发命令并组织 I/O”，中断处理程序负责“设备回话后先接住信号并做第一轮收尾”。}
\exam{PPT 单独列了它的功能，答层次题时不要只写“处理中断”，最好写出正常结束、异常事件、唤醒等待者这几类职责。}


\concept{设备驱动程序的职责}
\begin{points}
  \item 接收来自上层的抽象 I/O 请求。
  \item 把逻辑请求转换为设备控制器能理解的具体命令。
  \item 负责设备名到端口地址、逻辑记录到物理记录、逻辑操作到物理操作的转换。
  \item 启动设备，必要时把请求挂入等待队列，并响应控制器发来的中断。
\end{points}
\plain{驱动程序就是“上层说人话、下层说设备话”，它负责把两边翻译通。}
\exam{常考驱动程序和设备无关软件的分工，尤其是谁做地址换算、谁做通用缓冲和保护。}


\concept{设备无关软件的功能细分}
\begin{points}
  \item 设备命名：把逻辑设备名映射到实际设备和驱动程序。
  \item 设备保护：防止未授权进程直接使用设备。
  \item 提供统一逻辑块：屏蔽物理扇区大小差异。
  \item 缓冲管理、存储设备空间管理、独占设备分配与释放、通用错误处理。
\end{points}
\plain{设备无关软件做的是“所有设备都差不多要做的通用活”，不是跟某块网卡或某个磁盘死绑定。}
\exam{常考哪些功能属于设备无关软件，哪些必须由驱动程序完成。}


\concept{阻塞 I/O、非阻塞 I/O 与异步 I/O}
\begin{points}
  \item 阻塞 I/O：调用后进程挂起，直到 I/O 完成。
  \item 非阻塞 I/O：调用立即返回，可能只返回部分数据或暂时返回不能读写。
  \item 异步 I/O：I/O 在后台继续执行，完成后由 I/O 子系统再通知进程。
\end{points}
\plain{阻塞是“你等着”；非阻塞是“你先回去看看别的事”；异步是“事我替你办，办完再通知你”。}
\exam{常考非阻塞和异步的区别，这两个词不能混。}


\concept{设备独立性}
\begin{points}
  \item 设备独立性又称设备无关性，指用户程序不依赖具体物理设备。
  \item 第一层含义：独立于同类设备的具体设备号，例如打印到逻辑打印机而非某台具体打印机。
  \item 第二层含义：在一定抽象下独立于设备类型，例如统一用文件接口访问普通文件和部分设备文件。
  \item 好处：提高可移植性、灵活性和设备分配效率。
\end{points}
\plain{设备独立性的本质是“用户只说要什么设备能力，不必绑定某台具体硬件”，具体映射交给系统做。}
\exam{常考两层含义：独立于同类设备号、以及在一定抽象下独立于设备类型。}


\concept{逻辑设备表 LUT}
\begin{points}
  \item 逻辑设备表用于把逻辑设备名映射为物理设备名和驱动程序入口地址。
  \item 用户程序用逻辑设备名提出请求，系统运行时再完成映射。
  \item 这使 I/O 重定向和设备替换更容易实现。
\end{points}
\plain{LUT 就像“通讯录”：你记逻辑名字，系统去查真正该连哪台设备。}
\exam{常考设备独立性如何实现，标准回答里通常要提逻辑设备表。}


\section{缓冲技术管理}
\concept{为什么需要缓冲}
\begin{points}
  \item CPU 和 I/O 设备速度差异巨大，缓冲可缓和二者速度不匹配。
  \item 缓冲可减少中断次数和设备启动次数，提高并行程度。
  \item 缓冲还能适配不同数据粒度，例如用户按字节读写，磁盘按块传输。
\end{points}
\plain{引入缓冲不只是为了“快慢配速”，还为了把传输粒度、调用语义和并行重叠这几件事都协调起来。}
\exam{老师 PPT 把原因分成三类：速度不匹配、逻辑/物理记录大小不一致、复制语义。}


\concept{引入缓冲的另外两个原因}
\begin{points}
  \item 协调逻辑记录大小与物理记录大小不一致。
  \item 支持复制语义：系统调用返回后，即使用户缓冲区内容改变，也不应影响本次 I/O 应写出的数据。
\end{points}
\plain{缓冲不只是为了解决快慢差，还为了“先把数据稳稳拷住”，避免用户随后改内容把结果搞乱。}
\exam{常考“为什么 write 返回后还能保证写出的是调用当时的数据”，答案通常要提内核缓冲区。}


\concept{单缓冲}
\begin{points}
  \item 系统设置一个缓冲区，设备先把数据送入缓冲区，再由进程取走。
  \item 相比无缓冲，可让设备传输和进程处理部分重叠。
  \item 但一个缓冲区仍可能使设备或进程等待。
\end{points}
\plain{单缓冲下虽然能有部分重叠，但设备输入、系统传送、进程处理仍会围绕同一个缓冲区相互制约。}
\exam{常考单缓冲的时间分析，通常要判断三段时间哪些能并行、哪些必须串行。}


\concept{双缓冲}
\begin{points}
  \item 设置两个缓冲区，一个用于设备填充，另一个供进程处理。
  \item 可以更好地实现输入/输出与处理并行。
  \item 常用于数据流连续、生产和消费速度接近的场景。
\end{points}
\plain{双缓冲的核心直觉就是“一个在装，一个在用”，这样设备和进程更不容易互相等。}
\exam{常考它为什么比单缓冲并行性更强，以及大量数据处理时一块数据的时间怎么分析。}


\concept{循环缓冲}
\begin{points}
  \item 当双缓冲仍不足以平衡 I/O 与处理速度差异时，可继续增加缓冲区数量，构成首尾相连的环形缓冲队列。
  \item 循环缓冲特别适合连续数据流场景，使设备输入、系统传送和进程处理能在更多缓冲区之间流水重叠。
  \item 它更像为一类合作进程定制的连续传输通道。
\end{points}
\plain{循环缓冲的核心不是“缓冲区多”，而是把多个缓冲区按环组织起来，让数据流能一块接一块不断往前推。}
\exam{题目若问什么时候要从双缓冲进一步升级到循环缓冲，关键答“处理大量连续数据、速度差仍较大时”。}

\concept{缓冲池}
\begin{points}
  \item 缓冲池由系统统一管理一组可复用缓冲区，按需分配、回收并供多个进程共享。
  \item 与更偏专用传输通道的循环缓冲不同，缓冲池更适合多设备、多进程同时竞争缓冲资源的场景。
  \item 它既可用于输入，也可用于输出，目标是提高缓冲区整体利用率。
\end{points}
\plain{缓冲池更像操作系统统一维护的一批“公共周转箱”，谁需要谁来借，用完再还。}
\exam{选择题常比较循环缓冲和缓冲池：前者偏连续流专用，后者偏系统范围共享。}


\concept{缓冲区高速缓存}
\begin{points}
  \item 缓冲区高速缓存通常指把常用磁盘块副本保留在内存中，以减少再次访问慢速外设的次数。
  \item 它更接近 cache 的思路，主要利用局部性，而不是单纯协调传输速度。
  \item 对频繁访问的目录块、文件数据块、索引块等，命中后可直接从内存返回，提高 I/O 效率。
  \item 它与单缓冲、双缓冲、循环缓冲并不处于同一比较维度：前者偏“缓存热点块”，后者偏“协调传输过程”。
\end{points}
\plain{这部分容易和普通缓冲混淆，最稳的区分方式是看目标：缓冲区高速缓存是为了“少访盘”，普通缓冲更多是为了“顺利搬数据”。}
\exam{若题目专门问“缓冲”和“缓冲区高速缓存”关系，作答时要写出一个偏传输协调、一个偏局部性复用。}


\concept{缓冲与缓存区别}
\begin{points}
  \item 缓冲 buffer 主要解决速度不匹配和数据传输粒度不一致。
  \item 缓存 cache 主要利用局部性保存数据副本，减少慢速设备访问。
  \item 例子：磁盘块缓存是 cache；键盘输入缓冲区更偏 buffer。
\end{points}
\plain{buffer 更偏“传输时先放一放、调一调节奏”，cache 更偏“常用数据留下副本，下次少访问慢设备”。}
\exam{常考 buffer 和 cache 的目标区别，而不是只让你翻译两个英文单词。}


\concept{缓冲题的时间判断}
\begin{points}
  \item 若从设备输入一块数据到缓冲区时间为 $T$，从缓冲区送到用户区时间为 $M$，用户处理时间为 $C$。
  \item 处理大量数据时，单缓冲下每块数据的处理时间常按 $\max(T,C)+M$ 理解。
  \item 双缓冲下并行性更强，常按 $\max(T,M+C)$ 或结合教材具体流程判断。
  \item 做题时先画出“设备输入、系统传送、进程计算”三段能否重叠，不要直接背公式。
\end{points}
\plain{缓冲计算题本质是在问：这几段时间哪些能重叠，哪些必须串着来。}
\exam{常考单缓冲/双缓冲时间公式或让你自己推，先分清 T、M、C 的含义。}


\section{设备分配与 SPOOLing}
\concept{设备分配中的数据结构}
\begin{points}
  \item 设备控制表 DCT：描述一台设备的状态、类型、队列等。
  \item 控制器控制表 COCT：描述设备控制器状态。
  \item 通道控制表 CHCT：描述通道状态。
  \item 系统设备表 SDT：记录系统中所有设备及其控制表入口。
\end{points}
\plain{这些表的共同作用是把“设备、控制器、通道、等待队列”的状态逐层组织起来，便于系统按层分配资源。}
\exam{常考 DCT、COCT、CHCT、SDT 各自记录什么，以及它们之间的指针连接关系。}


\concept{设备分配过程}
\begin{points}
  \item 根据用户请求的逻辑设备名查找系统设备表。
  \item 分配设备、控制器和通道，若某一级不可用则进程等待。
  \item 分配成功后启动 I/O 操作，I/O 完成后释放相关资源并唤醒等待进程。
\end{points}
\plain{写过程题时按“查逻辑名、分配设备、再分配控制器和通道、启动 I/O、完成后释放”这个顺序展开。}
\exam{常考设备分配过程的步骤题，或让你说明某一级资源不足时进程为什么要等待。}


\concept{设备分配策略}
\begin{points}
  \item 设备分配要考虑设备使用性质、分配算法、安全性和设备独立性。
  \item 常见分配算法有先来先服务和优先级高者优先。
  \item 独占设备、共享设备、虚拟设备的分配方式并不相同。
\end{points}
\plain{设备分配不是“有空就给”，还要看设备类型、请求顺序和是否会带来死锁。}
\exam{常考设备分配时应考虑哪些因素，或区分独享分配、共享分配、虚拟分配。}

\concept{设备分配算法}
\begin{points}
  \item \term{先来先服务}：按进程提出设备请求的先后次序排队，设备优先分配给请求队列队首进程。
  \item \term{优先级高者优先}：先按进程优先级安排设备请求；若优先级相同，再按先来先服务排队。
  \item 这两类算法解决的都是“多个进程竞争同一设备时，系统先满足谁”的问题。
  \item 前者实现简单且相对公平；后者更利于保障关键任务，但可能延长低优先级进程等待时间。
\end{points}
\plain{这块不要和磁盘调度算法混淆：设备分配算法管的是“谁先拿到设备”，磁盘调度算法管的是“磁头接下来怎么走”。}
\exam{选择题常考设备分配算法有哪些，或问“同优先级请求接下来按什么规则排队”。}


\concept{独享分配、共享分配与虚拟分配}
\begin{points}
  \item \term{独享分配}：设备一旦分给某进程，就一直由该进程独占，直到完成或释放。
  \item \term{共享分配}：多个进程可共同使用同一设备，但系统必须安排合理的访问次序。
  \item \term{虚拟分配}：先给进程分配共享存储空间，再由系统在适当时机把这部分数据与真实独占设备交换。
  \item SPOOLing 就是把独占设备改造成逻辑共享设备的典型虚拟分配方法。
\end{points}
\plain{三者的根本区别在于：进程拿到的是“整台真实设备”、还是“被调度地共用设备”、还是“先拿一块共享存储区由系统代办 I/O”。}
\exam{PPT 把这三类分配方式分开讲了，常考打印机为什么更适合虚拟分配而不是直接共享分配。}


\concept{静态分配与动态分配}
\begin{points}
  \item 静态分配在作业开始前一次分配所需全部设备、控制器和通道，运行期间一直占有。
  \item 优点是不会发生设备分配死锁；缺点是利用率低。
  \item 动态分配在进程运行过程中按需分配、用完即释放。
  \item 优点是利用率高，但策略不当可能引发死锁。
\end{points}
\plain{静态分配像“开工前一次领齐工具”，动态分配像“用到什么再去借什么”。}
\exam{常考为什么静态分配安全但效率低，动态分配高效但有死锁风险。}


\concept{安全分配与不安全分配}
\begin{points}
  \item 安全分配下，进程发出 I/O 请求后进入阻塞，直到本次 I/O 完成并释放相关资源。
  \item 不安全分配下，进程发出一个 I/O 请求后还能继续执行，并继续发出后续 I/O 请求。
  \item 前者更安全但推进慢，后者并发性更强但可能造成死锁。
\end{points}
\plain{安全分配像“这件事办完再办下一件”，不安全分配像“先占着这个再去排下一个队”。}
\exam{常考设备分配中的安全性，和死锁章节形成呼应。}


\concept{单通路与多通路分配}
\begin{points}
  \item 单通路系统中，进程常直接按物理设备名申请设备，灵活性较差。
  \item 多通路系统中，用户更适合按逻辑设备名申请，系统再在多台同类设备、控制器、通道之间选择可用路径。
  \item 多通路结构提高了灵活性、可靠性和资源利用率。
\end{points}
\plain{单通路像只能指定某一台机器，多通路则像系统帮你自动找同类里空闲的那一台。}
\exam{常考为什么多通路系统更适合和设备独立性结合。}


\concept{单通路设备分配步骤}
\begin{points}
  \item 单通路分配通常假定进程按物理设备名提出请求，系统按“设备 -> 控制器 -> 通道”顺序逐级分配。
  \item 先由系统设备表 SDT 找到对应设备控制表 DCT，检查设备是否忙，必要时进入设备等待队列。
  \item 再由 DCT 找到控制器控制表 COCT，检查控制器状态；若忙则进入控制器等待队列。
  \item 最后由 COCT 找到通道控制表 CHCT，检查通道状态；若可用则完成分配并启动 I/O。
  \item 任一级资源忙碌，都可能导致进程在该级等待，而不是直接开始 I/O。
\end{points}
\plain{单通路设备分配题最稳的写法就是顺着 `SDT -> DCT -> COCT -> CHCT` 往下找，逐级判断谁忙谁闲。}
\exam{过程题常考每一级查什么状态、忙时进哪个等待队列，以及为什么必须先分到设备再分控制器和通道。}


\concept{SPOOLing 系统}
\begin{points}
  \item SPOOLing 是外部设备联机并行操作技术，常用于把独占设备改造成共享设备。
  \item 基本思想：把慢速独占设备的输入/输出先放到磁盘上的输入井/输出井中，由后台进程统一和实际设备交互。
  \item 典型做法可概括为：\key{预输入}、\key{缓输出}和\key{井管理程序}协同工作。
  \item 系统通常还设置内存中的输入缓冲区和输出缓冲区，用于暂存设备与井之间搬运的数据。
  \item 打印机例子：用户进程把打印数据写入磁盘输出井后即可继续执行，打印进程再按队列送往物理打印机。
  \item SPOOLing 组成：输入井、输出井、输入缓冲区、输出缓冲区、输入进程、输出进程、井管理程序。
  \item 直觉：用户不是直接占用打印机，而是把任务交到“打印队列”；物理打印机仍独占，但逻辑上多个用户可同时提交打印任务。
\end{points}
\plain{SPOOLing 的要害是用磁盘井加后台进程做中转，所以多个用户提交任务时，不必直接排队占物理打印机。}
\exam{常考 SPOOLing 的组成和工作流程，并说明它为什么能把独占设备改造成逻辑共享设备。}

\chapter*{总复习：高频辨析与做题模板}
\addcontentsline{toc}{chapter}{总复习：高频辨析与做题模板}

\section*{一、概念定位模板}
\addcontentsline{toc}{section}{一、概念定位模板}
\begin{points}
  \item 问 CPU 给谁用：进程调度。
  \item 问多个进程怎么不抢坏共享变量：同步互斥。
  \item 问互相等资源卡住：死锁。
  \item 问地址怎么变、页表、碎片：存储器管理。
  \item 问缺页、页面置换、抖动：虚拟存储器。
  \item 问文件名、目录、逻辑结构、访问方法：文件系统接口。
  \item 问文件块放磁盘哪里、inode、空闲块：文件系统实现。
  \item 问磁盘臂怎么移动、访问时间：大容量存储。
  \item 问 DMA、缓冲、SPOOLing、设备分配：设备管理。
\end{points}

\section*{二、地址转换题模板}
\addcontentsline{toc}{section}{二、地址转换题模板}
\begin{points}
  \item 看页大小/段大小，确定偏移位数或偏移范围。
  \item 把逻辑地址拆成页号/段号和偏移。
  \item 查页表/段表，判断是否越界或缺页。
  \item 拼物理地址：页式为物理块号 + 页内偏移；段式为段基址 + 段内偏移。
\end{points}

\section*{三、PV 题模板}
\addcontentsline{toc}{section}{三、PV 题模板}
\begin{points}
  \item 找共享资源：需要 mutex 互斥。
  \item 找先后关系：先发生的一方 V，后发生的一方 P。
  \item 找数量约束：空缓冲区用 empty，满缓冲区用 full，资源个数就是信号量初值。
  \item 写顺序：先等条件，再进互斥；先出互斥，再释放条件。
  \item 检查死锁：有没有拿着 mutex 等待别人改变条件。
\end{points}

\section*{四、调度题模板}
\addcontentsline{toc}{section}{四、调度题模板}
\begin{points}
  \item 列表：进程、到达时间、服务时间、优先级/队列。
  \item 找调度点：到达、完成、时间片到、阻塞、抢占、优先级调整。
  \item 画甘特图。
  \item 算完成时间、周转时间、等待时间、带权周转时间。
  \item 抢占式算法每个新进程到达都要判断是否抢占。
\end{points}

\section*{五、页面置换题模板}
\addcontentsline{toc}{section}{五、页面置换题模板}
\begin{points}
  \item 初始化页框为空，按引用串逐个访问。
  \item 命中：不计缺页，但更新 LRU/Clock 等状态。
  \item 缺页：若有空页框直接装入；否则按算法选牺牲页。
  \item 记录每一步页框内容和缺页次数。
  \item 判断 Belady：只针对 FIFO 这类算法可能出现，不要泛化到 LRU/OPT。
\end{points}

\section*{六、文件系统计算题模板}
\addcontentsline{toc}{section}{六、文件系统计算题模板}
\begin{points}
  \item 目录检索题：先算每块可放多少目录项，再算平均需要扫描多少块。
  \item 索引最大文件题：先算每个索引块可放多少块号，再根据索引级数求可寻址数据块数。
  \item 链式分配随机访问题：访问第 $k$ 块通常要从头顺链访问到第 $k$ 块。
  \item 连续分配题：直接定位强，但增长和外部碎片是弱点。
\end{points}

\section*{七、最后一遍背诵关键词}
\addcontentsline{toc}{section}{七、最后一遍背诵关键词}
\begin{multicols}{2}
\begin{points}
  \item OS：资源管理器、控制程序、接口。
  \item 四特征：并发、共享、虚拟、异步。
  \item 接口：命令接口、程序接口、图形接口。
  \item 进程：程序一次执行，PCB 是唯一标志。
  \item 线程：CPU 调度基本单位，共享进程资源。
  \item 调度：高级进系统，中级进内存，低级上 CPU。
  \item 临界区：空闲让进、忙则等待、有限等待、让权等待。
  \item 信号量：P 先减，V 先加。
  \item 死锁：互斥、请求保持、不可剥夺、循环等待。
  \item 分页：逻辑连续、物理离散、页表映射。
  \item 分段：按逻辑结构，便于共享保护。
  \item 虚存：局部性、请求调入、页面置换。
  \item 抖动：缺页太多，系统忙于换页。
  \item 文件：命名信息项序列，目录用于检索。
  \item inode：文件元数据和块指针。
  \item 磁盘：寻道 + 旋转延迟 + 传输。
  \item I/O：DMA、通道、缓冲、SPOOLing。
\end{points}
\end{multicols}

\end{document}
